% concatenated from icml_main.tex, icml_macro.tex, math_commands.tex, icml_main_appendix.tex
%%%%%%%% ICML 2026 EXAMPLE LATEX SUBMISSION FILE %%%%%%%%%%%%%%%%%

\documentclass{article}

% Recommended, but optional, packages for figures and better typesetting:
\usepackage{microtype}
\usepackage{graphicx}
\usepackage{subcaption}
\usepackage{subcaption}
\usepackage{orcidlink}
\usepackage{booktabs} % for professional tables
\usepackage{enumitem}
% hyperref makes hyperlinks in the resulting PDF.
% If your build breaks (sometimes temporarily if a hyperlink spans a page)
% please comment out the following usepackage line and replace
% \usepackage{icml2026} with \usepackage[nohyperref]{icml2026} above.
\usepackage{hyperref}


% Attempt to make hyperref and algorithmic work together better:
\newcommand{\theHalgorithm}{\arabic{algorithm}}

% Use the following line for the initial blind version submitted for review:
%\usepackage{icml2026}

% For preprint, use
\usepackage[preprint]{icml2026}

% If accepted, instead use the following line for the camera-ready submission:
% \usepackage[accepted]{icml2026}

\usepackage{amsmath}
\usepackage{amssymb}
\usepackage{mathtools}
\usepackage{amsthm}
\usepackage{comment}

% if you use cleveref..
\usepackage[capitalize,noabbrev]{cleveref}

%%%%%%%%%%%%%%%%%%%%%%%%%%%%%%%%
% THEOREMS
%%%%%%%%%%%%%%%%%%%%%%%%%%%%%%%%
\theoremstyle{plain}
\newtheorem{theorem}{Theorem}[section]
\newtheorem{proposition}[theorem]{Proposition}
\newtheorem{lemma}[theorem]{Lemma}
\newtheorem{corollary}[theorem]{Corollary}
\theoremstyle{definition}
\newtheorem{definition}[theorem]{Definition}
\newtheorem{assumption}[theorem]{Assumption}
\theoremstyle{remark}
\newtheorem{remark}[theorem]{Remark}

\newcommand{\red}[1]{{\color{red}#1}}
\newcommand{\orange}[1]{{\color{orange}#1}}
\newcommand{\blue}[1]{{\color{blue}#1}}
% Todonotes is useful during development; simply uncomment the next line
%    and comment out the line below the next line to turn off comments
%\usepackage[disable,textsize=tiny]{todonotes}
\usepackage[textsize=tiny]{todonotes}
\usepackage{mathrsfs}
\def\Avar{{\sf Avar}}
% \def\abs#1{\left|#1\right|}
\def\aff{{\sf aff}}
\def\argmin{{\arg\min}}
\def\argmax{{\arg\max}}
\def\avar{{\sf avar}}

\def\bA{\mathbf{A}}
\def\bB{\mathbf{B}}
\def\bC{\mathbf{C}}
\def\bD{\mathbf{D}}
% \def\bE{\mathbf{E}}
\def\bF{\mathbf{F}}
\def\bff{\mathbf{f}}
\def\bG{\mathbf{G}}
\def\bH{\mathbf{H}}
\def\bI{\mathbf{I}}
\def\bJ{\mathbf{J}}
\def\bK{\mathbf{K}}
\def\bL{\mathbf{L}}
\def\bM{\mathbf{M}}
\def\bN{\mathbf{N}}
\def\bO{\mathbf{O}}
\def\bP{\mathbf{P}}
\def\bQ{\mathbf{Q}}
\def\bR{\mathbf{R}}
\def\bS{\mathbf{S}}
\def\bT{\mathbf{T}}
\def\bU{\mathbf{U}}
\def\bV{\mathbf{V}}
\def\bW{\mathbf{W}}
\def\bX{\mathbf{X}}
\def\bY{\mathbf{Y}}
\def\bZ{\mathbf{Z}}

\def\sfA{\mathbb{A}}
\def\bbB{\mathbb{B}}
\def\bbC{\mathbb{C}}
\def\bbD{\mathbb{D}}
\def\bbE{\mathbb{E}}
\def\bbF{\mathbb{F}}
\def\bbG{\mathbb{G}}
\def\bbH{\mathbb{H}}
\def\bbI{\mathbb{I}}
\def\bbJ{\mathbb{J}}
\def\bbK{\mathbb{K}}
\def\bbL{\mathbb{L}}
\def\bbM{\mathbb{M}}
\def\bbN{\mathbb{N}}
\def\bbO{\mathbb{O}}
\def\bbP{\mathbb{P}}
\def\bbQ{\mathbb{Q}}
\def\bbR{\mathbb{R}}
\def\bbS{\mathbb{S}}
\def\bbT{\mathbb{T}}
\def\bbU{\mathbb{U}}
\def\bbV{\mathbb{V}}
\def\bbW{\mathbb{W}}
\def\bbX{\mathbb{X}}
\def\bbY{\mathbb{Y}}
\def\bbZ{\mathbb{Z}}

\def\bGamma{{\boldsymbol\Gamma}}
\def\bDelta{{\boldsymbol\Delta}}
\def\bTheta{{\boldsymbol\Theta}}
\def\bLambda{{\boldsymbol\Lambda}}
\def\bepsilon{{\boldsymbol\epsilon}}
\def\bvarepsilon{{\boldsymbol\varepsilon}}
\def\bXi{{\boldsymbol\Xi}}
\def\bPi{{\boldsymbol\Pi}}
\def\bSigma{{\boldsymbol\Sigma}}
\def\bUps{{\boldsymbol\Upsilon}}
\def\bPhi{{\boldsymbol\Phi}}
\def\bPsi{{\boldsymbol\Psi}}
\def\bOmega{{\boldsymbol\Omega}}

\def\phatdata{{\widehat{P}_{\rm data}}}
\def\pdiscdata{{\widehat{P}^{\rm disc}_{\rm data}}}


\def\bsr{{\boldsymbol{r}}}

%\def\ba{\mathbf{a}}
%\def\bb{\mathbf{b}}
% \def\bc{\mathbf{c}}
\def\bbd{\mathbf{d}}
\def\be{\mathbf{e}}
\def\bbf{\mathbf{f}}
\def\bg{\mathbf{g}}
\def\bh{\mathbf{h}}
\def\bi{\mathbf{i}}
\def\bj{\mathbf{j}}
%\def\bk{\mathbf{k}}
\def\bl{\mathbf{l}}
% \def\bm{\mathbf{m}}
\def\bn{\mathbf{n}}
\def\bo{\mathbf{o}}
\def\bp{\mathbf{p}}
\def\bq{\mathbf{q}}
% \def\br{\mathbf{r}}
% \def\bs{\mathbf{s}}
\def\bt{\mathbf{t}}
\def\bu{\mathbf{u}}
\def\bv{\mathbf{v}}
\def\bw{\mathbf{w}}
\def\bx{\mathbf{x}}
\def\by{\mathbf{y}}
\def\bz{\mathbf{z}}

\def\balpha{{\boldsymbol\alpha}}
\def\bbeta{{\boldsymbol\beta}}
\def\bgamma{{\boldsymbol\gamma}}
\def\bdelta{{\boldsymbol\delta}}
\def\beps{{\boldsymbol\epsilon}}
\def\bveps{{\boldsymbol\varepsilon}}
\def\bzeta{{\boldsymbol\zeta}}
\def\bfeta{{\boldsymbol\eta}}
\def\btheta{{\boldsymbol\theta}}
\def\bvtheta{{\boldsymbol\vartheta}}
\def\biota{{\boldsymbol\iota}}
\def\bkappa{{\boldsymbol\kappa}}
\def\blambda{{\boldsymbol\lambda}}
\def\bmu{{\boldsymbol\mu}}
\def\bnu{{\boldsymbol\nu}}
\def\bxi{{\boldsymbol\xi}}
\def\bpi{{\boldsymbol\pi}}
\def\brho{{\boldsymbol\rho}}
\def\bvrho{{\boldsymbol\varrho}}
\def\bsigma{{\boldsymbol\sigma}}
\def\btau{{\boldsymbol\tau}}
\def\bups{{\boldsymbol\upsilon}}
\def\bphi{{\boldsymbol\phi}}
\def\bvphi{{\boldsymbol\varphi}}
\def\bpsi{{\boldsymbol\psi}}
\def\bomega{{\boldsymbol\omega}}

\def\sfA{{\sf A}}
\def\sfB{{\sf B}}
\def\sfC{{\sf C}}
\def\sfD{{\sf D}}
\def\sfE{{\sf E}}
\def\sfF{{\sf G}}
\def\sfH{{\sf H}}
\def\sfI{{\sf I}}
\def\sfJ{{\sf J}}
\def\sfK{{\sf K}}
\def\sfL{{\sf L}}
\def\sfM{{\sf M}}
\def\sfN{{\sf N}}



\def\Ber{{\sf Ber}}
\def\bin{{\sf bin}}
\def\sfd{{\sf d}}
\def\sfr{{\sf r}}
\def\cA{\mathcal{A}}
\def\cB{\mathcal{B}}
\def\cC{\mathcal{C}}
\def\cD{\mathcal{D}}
\def\cE{\mathcal{E}}
\def\cF{\mathcal{F}}
\def\cG{\mathcal{G}}
\def\cH{\mathcal{H}}
\def\cI{\mathcal{I}}
\def\cJ{\mathcal{J}}
\def\cK{\mathcal{K}}
\def\cL{\mathcal{L}}
\def\cM{\mathcal{M}}
\def\cN{\mathcal{N}}
\def\cO{\mathcal{O}}
\def\cP{\mathcal{P}}
\def\cQ{\mathcal{Q}}
\def\cR{\mathcal{R}}
\def\cS{\mathcal{S}}
\def\cT{\mathcal{T}}
\def\cU{\mathcal{U}}
\def\cV{\mathcal{V}}
\def\cW{\mathcal{W}}
\def\cX{\mathcal{X}}
\def\cY{\mathcal{Y}}
\def\cZ{\mathcal{Z}}

\def\ccA{\mathscr{A}}
%\def\ccb{\mathscr{a}}
\def\ccB{\mathscr{B}}
\def\ccC{\mathscr{C}}
\def\ccD{\mathscr{D}}
\def\ccE{\mathscr{E}}
\def\ccEstar{{\mathscr{E}^{*}}}
\def\ccF{\mathscr{F}}
\def\ccG{\mathscr{G}}
\def\ccH{\mathscr{H}}
\def\ccI{\mathscr{I}}
\def\ccJ{\mathscr{J}}
\def\ccK{\mathscr{K}}
\def\ccL{\mathscr{L}}
\def\ccM{\mathscr{M}}
\def\ccN{\mathscr{N}}
\def\ccO{\mathscr{O}}
\def\ccP{\mathscr{P}}
\def\ccQ{\mathscr{Q}}
\def\ccR{\mathscr{R}}
\def\ccS{\mathscr{S}}
\def\ccT{\mathscr{T}}
\def\ccU{\mathscr{U}}
\def\ccV{\mathscr{V}}
\def\ccW{\mathscr{W}}
\def\ccX{\mathscr{X}}
\def\ccY{\mathscr{Y}}
\def\ccZ{\mathscr{Z}}

\usepackage{yfonts}
\def\gothm{\textgoth{m}}

\def\frakA{\mathfrak{A}}
\def\frakB{\mathfrak{B}}
\def\frakC{\mathfrak{C}}
\def\frakD{\mathfrak{D}}
\def\frakE{\mathfrak{E}}
\def\frakF{\mathfrak{F}}
\def\frakG{\mathfrak{G}}
\def\frakH{\mathfrak{H}}
\def\frakI{\mathfrak{I}}
\def\frakJ{\mathfrak{J}}
\def\frakK{\mathfrak{K}}
\def\frakL{\mathfrak{L}}
\def\frakM{\mathfrak{M}}
\def\frakN{\mathfrak{N}}
\def\frakO{\mathfrak{O}}
\def\frakP{\mathfrak{P}}
\def\frakQ{\mathfrak{Q}}
\def\frakR{\mathfrak{R}}
\def\frakS{\mathfrak{S}}
\def\frakT{\mathfrak{T}}
\def\frakU{\mathfrak{U}}
\def\frakV{\mathfrak{V}}
\def\frakW{\mathfrak{W}}
\def\frakX{\mathfrak{X}}
\def\frakY{\mathfrak{Y}}
\def\frakZ{\mathfrak{Z}}

 \def\xmax{x_{\sf max}}
 \def\bmax{b_{\sf max}}
 \def\pmh{\bP_{\sf MH}}
% \def\cat{{\sf cat}}
% \def\Cov{{\sf Cov}}
% \def\cov{{\sf cov}}
% \def\card{{\bf card}}
% \def\cf{{\it cf.}}
% \def\cl{{\sf cl}}
% \def\complex{{\bbC}}
% \def\cone{{\sf cone}}
% \def\conv{{\sf conv}}
% \def\deq{\overset{d}{=}}
% \def\det{{\sf det}}
% \def\diag{{\sf diag}}
% \def\diam{{\sf diam}}
% \def\dist{{\sf dist}}
% \def\dom{{\sf dom}}
% \def\doop{{\sf do}}
% \def\dto{\overset{\sf d}{\to}}
% \def\Ex{{\bbE}}
% \def\Exp{{\sf Exp}}
% %\def\exp{{\sf exp}}
% \def\eg{{\it e.g.}}
% \def\epi{{\sf epi}}
% \def\eps{{\epsilon}}
% \def\etc{{\it etc.}}
% \def\ie{{\it i.e.}}
% \def\iid{{\sf iid}}
% \def\ind{{\sf ind}}
% \def\intr{{\sf int}}
% \def\integers{{\bbZ}}

\def\halpha{\widehat{\balpha}}
\def\hbeta{\widehat{\bbeta}}
\def\hdelta{\widehat{\bdelta}}
\def\hGamma{\widehat{\bGamma}}
\def\hgamma{\widehat{\bgamma}}
\def\hmu{\widehat{\bmu}}
\def\hSigma{\widehat{\bSigma}}
\def\hsigma{\widehat{\bsigma}}
\def\hTheta{\widehat{\bTheta}}
\def\htheta{\widehat{\btheta}}
\def\hw{\widehat{\bw}}

\def\KL{{\sf KL}}
\def\Lap{{\sf Lap}}
\def\ker{{\sf ker}}
\def\lspan{{\sf span}}
\def\maximize{{\sf maximize}}
\def\med{{\sf med}}
\def\minimize{{\sf minimize}}
% \def\norm#1{\left\|#1\right\|}
\def\Norm#1{\Vert#1\Vert}
\def\innerprod#1{\left\langle#1\right\rangle}
\def\ones{\mathbf{1}}
\def\Poi{{\sf Poi}}
\def\Pr{{\bbP}}
\def\pr{{\bbP}}
\def\peq{\overset{\sf p}{=}}
\def\plim{{\sf plim}}
\def\proj{{\sf proj}}
\def\prox{{\sf prox}}
\def\pto{\overset{\sf p}{\to}}
\def\ran{{\sf ran}}
\def\rank{{\sf rank}}
\def\reals{{\bbR}}
\def\relint{{\sf relint}}
\def\sign{{\sf sign}}
\def\supp{{\sf supp}}
\def\symm{{\bbS}}

\def\st{\textrm{subject to}}
\def\d{{\rm d}}

\def\tr{{\sf tr}}
\def\tube{{\sf tube}}
\def\TV{{\sf TV}}
\def\unif{{\sf unif}}
\def\Var{{\sf Var}}
\def\var{{\sf var}}
\def\vct{{\sf vec}}
\def\veps{\varepsilon}
\def\vol{{\sf vol}}
\def\zeros{\mathbf{0}}
\def\op{{\rm op}}
\def\diam{{\sf diam}}

\def\diamT{\diam(\cT_s, \norm{\cdot}_2)}
\def\diamTsq{\diam^2(\cT_s, \norm{\cdot}_2)}

\def\diamG{\diam(\cG_s, \norm{\cdot}_{\op})}
\def\diamGsq{\diam^2(\cG_s, \norm{\cdot}_{\op})}
\def\emptyset{\varnothing}
\def\tol{\varepsilon_{n,p}}
\def\betapic{{\widehat{\beta}^{\sf{picasso}}}}
\def\betapgh{{\widehat{\beta}^{\sf{pgh}}}}
\def\ddt{{\frac{\rm d}{\rm dt}}}
\def\ddtau{{\frac{\rm d}{\rm d\tau}}}

%%%%%%%%%%%%%%%%%%%%%%%
%%%%%%% SCRAP %%%%%%%%%
%%%%%%%%%%%%%%%%%%%%%%%

\def\AIC{\textsc{aic}}
\def\Cov{\text{Cov}}
\def\Ber{{\sf {Ber}}}
\def\bss{{\sf BSS}}
\def\Gap{{\sf Gap}}
\def\Unif{{\rm Unif}}
\def\sfN{{\sf N}}
\def\sfI{{\sf I}}
\def\nspike{{\mathsf{n}_{\rm spike}}}
\def\aspike{{a_{\rm spike}}}
\def\pen{{\rm pen}}
\def\barPhi{\overline{\Phi}}
\def\col{\text{col}}
\def\dprime{{\prime\prime}}
\def\tprime{{\dprime\prime}}
% \def\kl{{\rm kl}}
\def\suff{{\rm suff}}
\def\param{{\rm par}}
\def\glm{{\rm glm}}
\def\best{{\rm best}}
\def\ctT{{\widetilde{\cT}}}
\def\ctG{{\widetilde{\cG}}}
\def\evel{\varepsilon^{\rm v}}
\def\escore{\varepsilon^{\rm s}}
\def\meanevel{\varepsilon_{\rm v}}
\def\meanescore{\varepsilon_{\rm s}}
\def\evelJone{\varepsilon^{\rm J,1}}
\def\meanevelJone{\varepsilon_{\rm J,1}}
\def\evelJtwo{\varepsilon^{\rm J,2}}
\def\meanevelJtwo{\varepsilon_{\rm J,2}}
\def\evelH{\varepsilon^{\rm H}}
\def\meanevelH{\varepsilon_{\rm H}}
\def\sqevel{{\varepsilon^2_{\rm vl}}}
\def\Dprime{{D^\prime}}
\def\tauD{{\tau^D_\eta}}
\def\tauDprime{{\tau^{D^\prime}_\eta}}
\def\gammastraight{{\gamma_{\rm st}}}
%\def\top{{ \mathrm{\scriptscriptstyle T} }}
\def\v{{\varepsilon}}
\def\PD{{\sf (PD*)}}
\def\tr{{\sf Trace}}
\def\rank{{\sf rank}}
%  \DeclareMathOperator{\E}{\mathbb{E}}
%  \DeclareMathOperator{\Surf}{\mathbb{S}}
%  %\DeclareMathOperator{\pr}{\mathbb{P}}
%  \DeclareMathOperator{\calX}{\mathcal{X}}
%  \DeclareMathOperator{\T}{\mathcal{T}}
%  \DeclareMathOperator{\calP}{\mathcal{P}}
% \def\M{{\mathcal{M}}}
% \def\R{{\mathbb{R}}}
% \DeclareMathOperator{\U}{\mathcal{U}}
%  \DeclareMathOperator{\calE}{\mathcal{E}}
% \DeclareMathOperator{\V}{\mathcal{V}}
% \DeclareMathOperator{\D}{\mathcal{D}}
% \DeclareMathOperator{\A}{\mathcal{A}}
% \DeclareMathOperator{\I}{\mathcal{I}}
% \DeclareMathOperator{\J}{\mathcal{J}}
% \DeclareMathOperator{\his}{\mathcal{H}}
%  \DeclareMathOperator{\calS}{\mathcal{S}}
\usepackage{bbm}
\def\ind{{\mathbbm{1}}}
%\DeclareMathOperator{\barPhi}{\overline{\Phi}}

%\newcommand{\ceil}[1]{\left\lceil#1\right\rceil}
%\newcommand{\floor}[1]{\left\lfloor#1\right\rfloor}
\mathchardef\mhyphen="2D
\def\dx{{\mathrm{d}x}}
\def\Ccover{C_{\rm cover}}


\newcommand{\bc}[1]{\left\{{#1}\right\}}
\newcommand{\br}[1]{\left({#1}\right)}
\newcommand{\bs}[1]{\left[{#1}\right]}
%%%%% NEW MATH DEFINITIONS %%%%%

%\usepackage{amsmath,amsfonts,bm,color}

% Mark sections of captions for referring to divisions of figures


\newcommand{\rd}{\color{red}}
\newcommand{\bk}{\color{black}}
\newcommand{\purna}[1]{\rd PS: #1\bk} 
\newcommand{\Yh}{\hat Y}
\newcommand{\figleft}{{\em (Left)}}
\newcommand{\figcenter}{{\em (Center)}}
\newcommand{\figright}{{\em (Right)}}
\newcommand{\figtop}{{\em (Top)}}
\newcommand{\figbottom}{{\em (Bottom)}}
\newcommand{\captiona}{{\em (a)}}
\newcommand{\captionb}{{\em (b)}}
\newcommand{\captionc}{{\em (c)}}
\newcommand{\captiond}{{\em (d)}}

\newcommand{\bas}[1]{\begin{align*}#1\end{align*}}
\newcommand{\ba}[1]{\begin{align}#1\end{align}}

% Highlight a newly defined term
\newcommand{\newterm}[1]{{\bf #1}}


% Figure reference, lower-case.
\def\figref#1{figure~\ref{#1}}
% Figure reference, capital. For start of sentence
\def\Figref#1{Figure~\ref{#1}}
\def\twofigref#1#2{figures \ref{#1} and \ref{#2}}
\def\quadfigref#1#2#3#4{figures \ref{#1}, \ref{#2}, \ref{#3} and \ref{#4}}
% Section reference, lower-case.
\def\secref#1{section~\ref{#1}}
% Section reference, capital.
\def\Secref#1{Section~\ref{#1}}
% Reference to two sections.
\def\twosecrefs#1#2{sections \ref{#1} and \ref{#2}}
% Reference to three sections.
\def\secrefs#1#2#3{sections \ref{#1}, \ref{#2} and \ref{#3}}
% Reference to an equation, lower-case.
%\def\eqref#1{equation~\ref{#1}}
% Reference to an equation, upper case
%\def\Eqref#1{Equation~\ref{#1}}
% A raw reference to an equation---avoid using if possible
\def\plaineqref#1{\ref{#1}}
% Reference to a chapter, lower-case.
\def\chapref#1{chapter~\ref{#1}}
% Reference to an equation, upper case.
\def\Chapref#1{Chapter~\ref{#1}}
% Reference to a range of chapters
\def\rangechapref#1#2{chapters\ref{#1}--\ref{#2}}
% Reference to an algorithm, lower-case.
\def\algref#1{algorithm~\ref{#1}}
% Reference to an algorithm, upper case.
\def\Algref#1{Algorithm~\ref{#1}}
\def\twoalgref#1#2{algorithms \ref{#1} and \ref{#2}}
\def\Twoalgref#1#2{Algorithms \ref{#1} and \ref{#2}}
% Reference to a part, lower case
\def\partref#1{part~\ref{#1}}
% Reference to a part, upper case
\def\Partref#1{Part~\ref{#1}}
\def\twopartref#1#2{parts \ref{#1} and \ref{#2}}

\def\ceil#1{\lceil #1 \rceil}
\def\floor#1{\lfloor #1 \rfloor}
\def\1{\bm{1}}
\newcommand{\train}{\mathcal{D}}
\newcommand{\valid}{\mathcal{D_{\mathrm{valid}}}}
\newcommand{\test}{\mathcal{D_{\mathrm{test}}}}

\def\eps{{\epsilon}}


% Random variables
\def\reta{{\textnormal{$\eta$}}}
\def\ra{{\textnormal{a}}}
\def\rb{{\textnormal{b}}}
\def\rc{{\textnormal{c}}}
%\def\rd{{\textnormal{d}}}
\def\re{{\textnormal{e}}}
\def\rf{{\textnormal{f}}}
\def\rg{{\textnormal{g}}}
\def\rh{{\textnormal{h}}}
\def\ri{{\textnormal{i}}}
\def\rj{{\textnormal{j}}}
\def\rk{{\textnormal{k}}}
\def\rl{{\textnormal{l}}}
% rm is already a command, just don't name any random variables m
\def\rn{{\textnormal{n}}}
\def\ro{{\textnormal{o}}}
\def\rp{{\textnormal{p}}}
\def\rq{{\textnormal{q}}}
\def\rr{{\textnormal{r}}}
\def\rs{{\textnormal{s}}}
\def\rt{{\textnormal{t}}}
\def\ru{{\textnormal{u}}}
\def\rv{{\textnormal{v}}}
\def\rw{{\textnormal{w}}}
\def\rx{{\textnormal{x}}}
\def\ry{{\textnormal{y}}}
\def\rz{{\textnormal{z}}}

% Random vectors
\def\rvepsilon{{\mathbf{\epsilon}}}
\def\rvtheta{{\mathbf{\theta}}}
\def\rva{{\mathbf{a}}}
\def\rvb{{\mathbf{b}}}
\def\rvc{{\mathbf{c}}}
\def\rvd{{\mathbf{d}}}
\def\rve{{\mathbf{e}}}
\def\rvf{{\mathbf{f}}}
\def\rvg{{\mathbf{g}}}
\def\rvh{{\mathbf{h}}}
\def\rvu{{\mathbf{i}}}
\def\rvj{{\mathbf{j}}}
\def\rvk{{\mathbf{k}}}
\def\rvl{{\mathbf{l}}}
\def\rvm{{\mathbf{m}}}
\def\rvn{{\mathbf{n}}}
\def\rvo{{\mathbf{o}}}
\def\rvp{{\mathbf{p}}}
\def\rvq{{\mathbf{q}}}
\def\rvr{{\mathbf{r}}}
\def\rvs{{\mathbf{s}}}
\def\rvt{{\mathbf{t}}}
\def\rvu{{\mathbf{u}}}
\def\rvv{{\mathbf{v}}}
\def\rvw{{\mathbf{w}}}
\def\rvx{{\mathbf{x}}}
\def\rvy{{\mathbf{y}}}
\def\rvz{{\mathbf{z}}}

% Elements of random vectors
\def\erva{{\textnormal{a}}}
\def\ervb{{\textnormal{b}}}
\def\ervc{{\textnormal{c}}}
\def\ervd{{\textnormal{d}}}
\def\erve{{\textnormal{e}}}
\def\ervf{{\textnormal{f}}}
\def\ervg{{\textnormal{g}}}
\def\ervh{{\textnormal{h}}}
\def\ervi{{\textnormal{i}}}
\def\ervj{{\textnormal{j}}}
\def\ervk{{\textnormal{k}}}
\def\ervl{{\textnormal{l}}}
\def\ervm{{\textnormal{m}}}
\def\ervn{{\textnormal{n}}}
\def\ervo{{\textnormal{o}}}
\def\ervp{{\textnormal{p}}}
\def\ervq{{\textnormal{q}}}
\def\ervr{{\textnormal{r}}}
\def\ervs{{\textnormal{s}}}
\def\ervt{{\textnormal{t}}}
\def\ervu{{\textnormal{u}}}
\def\ervv{{\textnormal{v}}}
\def\ervw{{\textnormal{w}}}
\def\ervx{{\textnormal{x}}}
\def\ervy{{\textnormal{y}}}
\def\ervz{{\textnormal{z}}}

% Random matrices
\def\rmA{{\mathbf{A}}}
\def\rmB{{\mathbf{B}}}
\def\rmC{{\mathbf{C}}}
\def\rmD{{\mathbf{D}}}
\def\rmE{{\mathbf{E}}}
\def\rmF{{\mathbf{F}}}
\def\rmG{{\mathbf{G}}}
\def\rmH{{\mathbf{H}}}
\def\rmI{{\mathbf{I}}}
\def\rmJ{{\mathbf{J}}}
\def\rmK{{\mathbf{K}}}
\def\rmL{{\mathbf{L}}}
\def\rmM{{\mathbf{M}}}
\def\rmN{{\mathbf{N}}}
\def\rmO{{\mathbf{O}}}
\def\rmP{{\mathbf{P}}}
\def\rmQ{{\mathbf{Q}}}
\def\rmR{{\mathbf{R}}}
\def\rmS{{\mathbf{S}}}
\def\rmT{{\mathbf{T}}}
\def\rmU{{\mathbf{U}}}
\def\rmV{{\mathbf{V}}}
\def\rmW{{\mathbf{W}}}
\def\rmX{{\mathbf{X}}}
\def\rmY{{\mathbf{Y}}}
\def\rmZ{{\mathbf{Z}}}

% Elements of random matrices
\def\ermA{{\textnormal{A}}}
\def\ermB{{\textnormal{B}}}
\def\ermC{{\textnormal{C}}}
\def\ermD{{\textnormal{D}}}
\def\ermE{{\textnormal{E}}}
\def\ermF{{\textnormal{F}}}
\def\ermG{{\textnormal{G}}}
\def\ermH{{\textnormal{H}}}
\def\ermI{{\textnormal{I}}}
\def\ermJ{{\textnormal{J}}}
\def\ermK{{\textnormal{K}}}
\def\ermL{{\textnormal{L}}}
\def\ermM{{\textnormal{M}}}
\def\ermN{{\textnormal{N}}}
\def\ermO{{\textnormal{O}}}
\def\ermP{{\textnormal{P}}}
\def\ermQ{{\textnormal{Q}}}
\def\ermR{{\textnormal{R}}}
\def\ermS{{\textnormal{S}}}
\def\ermT{{\textnormal{T}}}
\def\ermU{{\textnormal{U}}}
\def\ermV{{\textnormal{V}}}
\def\ermW{{\textnormal{W}}}
\def\ermX{{\textnormal{X}}}
\def\ermY{{\textnormal{Y}}}
\def\ermZ{{\textnormal{Z}}}

% Vectors
\def\vzero{{\bm{0}}}
\def\vone{{\bm{1}}}
\def\vmu{{\bm{\mu}}}
\def\vtheta{{\bm{\theta}}}
\def\va{{\bm{a}}}
\def\vb{{\bm{b}}}
\def\vc{{\bm{c}}}
\def\vd{{\bm{d}}}
\def\ve{{\bm{e}}}
\def\vf{{\bm{f}}}
\def\vg{{\bm{g}}}
\def\vh{{\bm{h}}}
\def\vi{{\bm{i}}}
\def\vj{{\bm{j}}}
\def\vk{{\bm{k}}}
\def\vl{{\bm{l}}}
\def\vm{{\bm{m}}}
\def\vn{{\bm{n}}}
\def\vo{{\bm{o}}}
\def\vp{{\bm{p}}}
\def\vq{{\bm{q}}}
\def\vr{{\bm{r}}}
\def\vs{{\bm{s}}}
\def\vt{{\bm{t}}}
\def\vu{{\bm{u}}}
\def\vv{{\bm{v}}}
\def\vw{{\bm{w}}}
\def\vx{{\bm{x}}}
\def\vy{{\bm{y}}}
\def\vz{{\bm{z}}}

% Elements of vectors
\def\evalpha{{\alpha}}
\def\evbeta{{\beta}}
\def\evepsilon{{\epsilon}}
\def\evlambda{{\lambda}}
\def\evomega{{\omega}}
\def\evmu{{\mu}}
\def\evpsi{{\psi}}
\def\evsigma{{\sigma}}
\def\evtheta{{\theta}}
\def\eva{{a}}
\def\evb{{b}}
\def\evc{{c}}
\def\evd{{d}}
\def\eve{{e}}
\def\evf{{f}}
\def\evg{{g}}
\def\evh{{h}}
\def\evi{{i}}
\def\evj{{j}}
\def\evk{{k}}
\def\evl{{l}}
\def\evm{{m}}
\def\evn{{n}}
\def\evo{{o}}
\def\evp{{p}}
\def\evq{{q}}
\def\evr{{r}}
\def\evs{{s}}
\def\evt{{t}}
\def\evu{{u}}
\def\evv{{v}}
\def\evw{{w}}
\def\evx{{x}}
\def\evy{{y}}
\def\evz{{z}}

% Matrix
\def\mA{{\bm{A}}}
\def\mB{{\bm{B}}}
\def\mC{{\bm{C}}}
\def\mD{{\bm{D}}}
\def\mE{{\bm{E}}}
\def\mF{{\bm{F}}}
\def\mG{{\bm{G}}}
\def\mH{{\bm{H}}}
\def\mI{{\bm{I}}}
\def\mJ{{\bm{J}}}
\def\mK{{\bm{K}}}
\def\mL{{\bm{L}}}
\def\mM{{\bm{M}}}
\def\mN{{\bm{N}}}
\def\mO{{\bm{O}}}
\def\mP{{\bm{P}}}
\def\mQ{{\bm{Q}}}
\def\mR{{\bm{R}}}
\def\mS{{\bm{S}}}
\def\mT{{\bm{T}}}
\def\mU{{\bm{U}}}
\def\mV{{\bm{V}}}
\def\mW{{\bm{W}}}
\def\mX{{\bm{X}}}
\def\mY{{\bm{Y}}}
\def\mZ{{\bm{Z}}}
\def\mBeta{{\bm{\beta}}}
\def\mPhi{{\bm{\Phi}}}
\def\mLambda{{\bm{\Lambda}}}
\def\mSigma{{\bm{\Sigma}}}

% Tensor
\DeclareMathAlphabet{\mathsfit}{\encodingdefault}{\sfdefault}{m}{sl}
\SetMathAlphabet{\mathsfit}{bold}{\encodingdefault}{\sfdefault}{bx}{n}
\newcommand{\tens}[1]{\bm{\mathsfit{#1}}}
\def\tA{{\tens{A}}}
\def\tB{{\tens{B}}}
\def\tC{{\tens{C}}}
\def\tD{{\tens{D}}}
\def\tE{{\tens{E}}}
\def\tF{{\tens{F}}}
\def\tG{{\tens{G}}}
\def\tH{{\tens{H}}}
\def\tI{{\tens{I}}}
\def\tJ{{\tens{J}}}
\def\tK{{\tens{K}}}
\def\tL{{\tens{L}}}
\def\tM{{\tens{M}}}
\def\tN{{\tens{N}}}
\def\tO{{\tens{O}}}
\def\tP{{\tens{P}}}
\def\tQ{{\tens{Q}}}
\def\tR{{\tens{R}}}
\def\tS{{\tens{S}}}
\def\tT{{\tens{T}}}
\def\tU{{\tens{U}}}
\def\tV{{\tens{V}}}
\def\tW{{\tens{W}}}
\def\tX{{\tens{X}}}
\def\tY{{\tens{Y}}}
\def\tZ{{\tens{Z}}}


% Graph
\def\gA{{\mathcal{A}}}
\def\gB{{\mathcal{B}}}
\def\gC{{\mathcal{C}}}
\def\gD{{\mathcal{D}}}
\def\gE{{\mathcal{E}}}
\def\gF{{\mathcal{F}}}
\def\gG{{\mathcal{G}}}
\def\gH{{\mathcal{H}}}
\def\gI{{\mathcal{I}}}
\def\gJ{{\mathcal{J}}}
\def\gK{{\mathcal{K}}}
\def\gL{{\mathcal{L}}}
\def\gM{{\mathcal{M}}}
\def\gN{{\mathcal{N}}}
\def\gO{{\mathcal{O}}}
\def\gP{{\mathcal{P}}}
\def\gQ{{\mathcal{Q}}}
\def\gR{{\mathcal{R}}}
\def\gS{{\mathcal{S}}}
\def\gT{{\mathcal{T}}}
\def\gU{{\mathcal{U}}}
\def\gV{{\mathcal{V}}}
\def\gW{{\mathcal{W}}}
\def\gX{{\mathcal{X}}}
\def\gY{{\mathcal{Y}}}
\def\gZ{{\mathcal{Z}}}

% Sets
\def\sA{{\mathbb{A}}}
\def\sB{{\mathbb{B}}}
\def\sC{{\mathbb{C}}}
\def\sD{{\mathbb{D}}}
% Don't use a set called E, because this would be the same as our symbol
% for expectation.
\def\sF{{\mathbb{F}}}
\def\sG{{\mathbb{G}}}
\def\sH{{\mathbb{H}}}
\def\sI{{\mathbb{I}}}
\def\sJ{{\mathbb{J}}}
\def\sK{{\mathbb{K}}}
\def\sL{{\mathbb{L}}}
\def\sM{{\mathbb{M}}}
\def\sN{{\mathbb{N}}}
\def\sO{{\mathbb{O}}}
\def\sP{{\mathbb{P}}}
\def\sQ{{\mathbb{Q}}}
\def\sR{{\mathbb{R}}}
\def\sS{{\mathbb{S}}}
\def\sT{{\mathbb{T}}}
\def\sU{{\mathbb{U}}}
\def\sV{{\mathbb{V}}}
\def\sW{{\mathbb{W}}}
\def\sX{{\mathbb{X}}}
\def\sY{{\mathbb{Y}}}
\def\sZ{{\mathbb{Z}}}

% Entries of a matrix
\def\emLambda{{\Lambda}}
\def\emA{{A}}
\def\emB{{B}}
\def\emC{{C}}
\def\emD{{D}}
\def\emE{{E}}
\def\emF{{F}}
\def\emG{{G}}
\def\emH{{H}}
\def\emI{{I}}
\def\emJ{{J}}
\def\emK{{K}}
\def\emL{{L}}
\def\emM{{M}}
\def\emN{{N}}
\def\emO{{O}}
\def\emP{{P}}
\def\emQ{{Q}}
\def\emR{{R}}
\def\emS{{S}}
\def\emT{{T}}
\def\emU{{U}}
\def\emV{{V}}
\def\emW{{W}}
\def\emX{{X}}
\def\emY{{Y}}
\def\emZ{{Z}}
\def\emSigma{{\Sigma}}

% entries of a tensor
% Same font as tensor, without \bm wrapper
\newcommand{\etens}[1]{\mathsfit{#1}}
\def\etLambda{{\etens{\Lambda}}}
\def\etA{{\etens{A}}}
\def\etB{{\etens{B}}}
\def\etC{{\etens{C}}}
\def\etD{{\etens{D}}}
\def\etE{{\etens{E}}}
\def\etF{{\etens{F}}}
\def\etG{{\etens{G}}}
\def\etH{{\etens{H}}}
\def\etI{{\etens{I}}}
\def\etJ{{\etens{J}}}
\def\etK{{\etens{K}}}
\def\etL{{\etens{L}}}
\def\etM{{\etens{M}}}
\def\etN{{\etens{N}}}
\def\etO{{\etens{O}}}
\def\etP{{\etens{P}}}
\def\etQ{{\etens{Q}}}
\def\etR{{\etens{R}}}
\def\etS{{\etens{S}}}
\def\etT{{\etens{T}}}
\def\etU{{\etens{U}}}
\def\etV{{\etens{V}}}
\def\etW{{\etens{W}}}
\def\etX{{\etens{X}}}
\def\etY{{\etens{Y}}}
\def\etZ{{\etens{Z}}}

% The true underlying data generating distribution
\newcommand{\pdata}{p_{\rm{data}}}
\def\phatdatadisc{{\widehat{P}_{\rm data}^{\rm disc}}}
% The empirical distribution defined by the training set
\newcommand{\rhotrain}{\hat{rho}_{\rm{data}}}
\newcommand{\Ptrain}{\hat{P}_{\rm{data}}}
% The model distribution
\newcommand{\rhomodel}{\widehat \rho_{\rm{data}}}
\newcommand{\Pmodel}{P_{\rm{model}}}
\newcommand{\ptildemodel}{\tilde{p}_{\rm{model}}}
% Stochastic autoencoder distributions
\newcommand{\pencode}{p_{\rm{encoder}}}
\newcommand{\pdecode}{p_{\rm{decoder}}}
\newcommand{\precons}{p_{\rm{reconstruct}}}

\newcommand{\laplace}{\mathrm{Laplace}} % Laplace distribution
\newcommand{\Law}{\text{Law}}
\newcommand{\Rectify}[1]{\texttt{Rectify}\br{#1}}
\newcommand{\Rectflow}[1]{\texttt{Rectflow}\br{#1}}
% \newcommand{\E}{\mathbb{E}}
\newcommand{\Ls}{\mathcal{L}}
\newcommand{\R}{\mathbb{R}}
\newcommand{\emp}{\tilde{p}}
\newcommand{\lr}{\alpha}
\newcommand{\reg}{\lambda}
\newcommand{\rect}{\mathrm{rectifier}}
\newcommand{\softmax}{\mathrm{softmax}}
\newcommand{\sigmoid}{\sigma}
\newcommand{\softplus}{\zeta}
% \newcommand{\KL}{D_{\mathrm{KL}}}
% \newcommand{\Var}{\mathrm{Var}}
% \newcommand{\standarderror}{\mathrm{SE}}
% \newcommand{\Cov}{\mathrm{Cov}}
% Wolfram Mathworld says $L^2$ is for function spaces and $\ell^2$ is for vectors
% But then they seem to use $L^2$ for  vectors throughout the site, and so does
% wikipedia.
\newcommand{\normlzero}{L^0}
\newcommand{\normlone}{L^1}
\newcommand{\normltwo}{L^2}
\newcommand{\normlp}{L^p}
\newcommand{\normmax}{L^\infty}

\newcommand{\parents}{Pa} % See usage in notation.tex. Chosen to match Daphne's book.

% \DeclareMathOperator*{\argmax}{arg\,max}
% \DeclareMathOperator*{\argmin}{arg\,min}

% \DeclareMathOperator{\sign}{sign}
\DeclareMathOperator{\Tr}{Tr}
\let\ab\allowbreak

%\usepackage{style}
% The \icmltitle you define below is probably too long as a header.
% Therefore, a short form for the running title is supplied here:
\icmltitlerunning{}

\begin{document}

\onecolumn
  \icmltitle{Low-Dimensional Adaptation of Rectified Flow: A Diffusion and Stochastic
Localization Perspective
}

  % It is OKAY to include author information, even for blind submissions: the
  % style file will automatically remove it for you unless you've provided
  % the [accepted] option to the icml2026 package.

  % List of affiliations: The first argument should be a (short) identifier you
  % will use later to specify author affiliations Academic affiliations
  % should list Department, University, City, Region, Country Industry
  % affiliations should list Company, City, Region, Country

  % You can specify symbols, otherwise they are numbered in order. Ideally, you
  % should not use this facility. Affiliations will be numbered in order of
  % appearance and this is the preferred way.
  \icmlsetsymbol{equal}{*}
   \begin{icmlauthorlist}
    \icmlauthor{Saptarshi Roy}{}
    \quad
    \icmlauthor{Alessandro Rinaldo}{}
    \quad
    \icmlauthor{Purnamrita Sarkar}{}\\
    \vspace{0.2em}
    Department of Statistics and Data Science\\
    University of Texas at Austin, TX, USA
    % \icmlauthor{Firstname4 Lastname4}{sch}
    % \icmlauthor{Firstname5 Lastname5}{yyy}
    % \icmlauthor{Firstname6 Lastname6}{sch,yyy,comp}
    % \icmlauthor{Firstname7 Lastname7}{comp}
    % %\icmlauthor{}{sch}
    % \icmlauthor{Firstname8 Lastname8}{sch}
    % \icmlauthor{Firstname8 Lastname8}{yyy,comp}
    % %\icmlauthor{}{sch}
    %\icmlauthor{}{sch}
  \end{icmlauthorlist}
  %\icmlaffiliation{}{Department of Statistics and Data Science, University of Texas at Austin, Austin, TX, USA}
  % \icmlaffiliation{comp}{Company Name, Location, Country}
  % \icmlaffiliation{sch}{School of ZZZ, Institute of WWW, Location, Country}

  \icmlcorrespondingauthor{Saptarshi Roy}{saptarshi.roy@austin.utexas.edu}
  %\icmlcorrespondingauthor{Firstname2 Lastname2}{first2.last2@www.uk}

  % You may provide any keywords that you find helpful for describing your
  % paper; these are used to populate the "keywords" metadata in the PDF but
  % will not be shown in the document
  \icmlkeywords{Machine Learning, ICML}
  \vskip 0.3in


% this must go after the closing bracket ] following \twocolumn[ ...

% This command actually creates the footnote in the first column listing the
% affiliations and the copyright notice. The command takes one argument, which
% is text to display at the start of the footnote. The \icmlEqualContribution
% command is standard text for equal contribution. Remove it (just {}) if you
% do not need this facility.

% Use ONE of the following lines. DO NOT remove the command.
% If you have no special notice, KEEP empty braces:
\printAffiliationsAndNotice{}  % no special notice (required even if empty)
% Or, if applicable, use the standard equal contribution text:
% \printAffiliationsAndNotice{\icmlEqualContribution}

\begin{abstract}
  In recent years, Rectified flow (RF) has gained considerable popularity largely due to its generation efficiency and state-of-the-art performance. In this paper, we investigate how well RF automatically adapts to the intrinsic low dimensionality of the support of the target distribution to accelerate sampling. We show that, using a carefully designed choice of the time-discretization scheme and with sufficiently accurate drift estimates, the RF sampler enjoys an iteration complexity of order $O(k/\varepsilon)$ (up to log factors), where $\varepsilon$ is the precision in total variation distance and $k$ is the intrinsic dimension of
 the target distribution. In addition, we show that the denoising diffusion probabilistic model (DDPM) procedure is equivalent to a stochastic version of RF by establishing a novel connection between these processes and stochastic localization. Building on this connection, we further design a stochastic RF sampler that also adapts to the low-dimensionality of the target distribution under mild requirements on the accuracy of the drift estimates, and also with a specific time schedule. We illustrate the efficacy of newly designed time-discretization schedules with simulations on the synthetic data and text-to-image (T2I) data experiments. 
\end{abstract}

\section{Introduction}
Recently, Rectified flow or Flow matching \cite{liu2022rectified, lipman2022flow} has gained much attention due to its state-of-the-art generative performance across various modalities including image \citep{liu2023instaflow, yan2024perflow}, audio \citep{ audio_RF_wang2024frieren, audio_RF_liu2025flashaudio}, and video \citep{video_RF_chen2025goku, video_RF_liu2025improving}. 
Its effectiveness arises from a deterministic ordinary differential equation (ODE) based sampler and the ability of reflow to learn nearly straight transport paths, resulting in faster sampling. This has naturally sparked theoretical interests \cite{bansal2025wasserstein, guan2025mirror, zhou2025an} to understand the generation quality of RF from a convergence viewpoint under a suitable metric. However, these recent results suffer from \textit{curse of dimensionality}. Therefore, they fail to explain the impressive performance of RF for \textit{high-dimensional} data, e.g., image data, even though they reside in a low-dimensional space \cite{ansuini2019intrinsic, pope2021intrinsic}. In contrast, recent works \citep{potaptchik2024linear, azangulov2024convergence,huang2024denoising,li2024adapting,liang2025low} have established low-dimensional adaptation of DDPM and denoising diffusion implicit models (DDIM)  during sampling. This explains the accelerated and high-quality generative capabilities of DDPM and DDIM in practice, even for high-dimensional data.
% However, RF still faces several challenges. 
% \begin{itemize}
%     \item The reflow process is
% computationally intensive as it requires a large number of generated pairs to model the target
% distribution.
% \item It is well known that the multiple reflow steps
% can quickly lead to performance degradation of the learned flow model \citep{model_collapse_zhu2024analyzing, reflow_bad_kim2025balanced}.
% \end{itemize}
% Therefore, \textit{understanding and mitigating the learning error in the first reflow itself is extremely important} to improve generation quality. 
% There has been prior works  that shows the low-dimensional adaptation of DDPM. Therefore, DDPM enjoys accelerated convergence to achieve improved sampling which is consistent with empirical findings even when the ambient dimension is very high. 
To the best of our knowledge such results are not currently present for RF models.
%  A recent work by \citet{lee2024improvingtrainingrectifiedflows} has empirically shown that a certain U-shaped time-discretization improves the generative quality of the first reflow model. However, proper theoretical understanding of such phenomenon remains elusive and necessitates in-depth investigation.
To this end, we make three key contributions in this paper: 
\begin{enumerate}
    \item We design a novel U-shaped time-discretization that provably allows RF to \textit{quickly achieve improved generation quality} in the \textit{first reflow} process by adapting to the intrinsic low-dimensionality of the target distribution. 
    This aligns with the recent work
    by \citet{lee2024improvingtrainingrectifiedflows}, which has empirically shown that a certain
U-shaped time-discretization improves the generative quality of the first reflow model. 
    % , thus confirming theoretically the experimental findings of \citet{lee2024improvingtrainingrectifiedflows}. 
    In particular, the RF sampler enjoys $O(k/\varepsilon)$ (upto log factors) iteration complexity, where $\varepsilon$ is the precision of the generated samples in total variation distance and $k$ is the intrinsic dimension of the target distribution. This is particularly useful because it can eliminate repeated reflow steps, which are computationally costly and often cause model collapse~\citep{model_collapse_zhu2024analyzing, reflow_bad_kim2025balanced}. 


 \item 
 Using a novel stochastic localization \cite{eldan2013thin,eldan2020taming} perspective, we demonstrate that DDPM is equivalent to the stochastic version of RF (which we refer to as STOC-RF) under a proper time change. Stochastic localization has been recently shown to be provide a unifying and effective framework for representing and analyzing diffusion-based sampling processes \citep{benton2023nearly, montanari2023sampling, el2022information}. 
% We demonstrate that DDPM is equivalent to a stochastic version of RF (which we refer to as STOC-RF) under a proper time change by establishing a novel connection between these processes and stochastic localization \cite{eldan2013thin,eldan2020taming}, which has been recently shown to be provide a unifying and effective framework for representing and analyzing diffusion-based sampling processes \citep{benton2023nearly, montanari2023sampling, el2022information}. 
 This equivalence is convenient as it allows us to view STOC-RF through the lens of well-understood DDPM processes, which are known to adapt to the low-dimensionality of the target distribution. In fact, our analysis shows that stochastic localization can be leveraged to provide a unified and convenient representation of a variety of sampling methods. %acts as a central process that bridges two different stochastic interpolants between source and target distributions. 

 \item Using the connection between DDPM and STOC-RF, we further design a stochastic RF sampler that provably also adapts to the low-dimensional structure of the target distributions under less stringent assumptions on the learned model. We illustrate the superior performance of this sampler with various simulations. This also addresses one of the themes discussed in the ICML tutorial talk by \citet{liu2025flowing}.
\end{enumerate}
%Next, we introduce some key notations used in this paper.
\textbf{Notation.} {Throughout the paper, with some abuse of notations, we will identify  a probability distribution with its Lebesgue density, which we will assume exists unless otherwise noted.} 
%For a distribution (Lebesgue density) $p$, we write $\supp(p)$ to denote the support of $p$. For a random variable $X$ we use $\Law(X)$ to denote its distribution. 
For two  distributions $P$ and $Q$ on $\mathbb{R}^d$ with Lebesgue densities $p$ and $q$ respectively, their total variation distance is denoted by $\TV(p,q):= \sup_{A} \vert{P(A) - Q(A)}\vert = 1/2 \int_{\mathbb{R}^d} |p(x) - q(x)| dx $ , where the supremum is over all Borel subsets of $\mathbb{R}^d$.

\section{Background and Preliminaries}
\label{sec: background}


\subsection{Rectified Flow}
\label{sec: Rf}
In this section, we briefly introduce the basics of Rectified flow, a generative model that transitions between two distributions $p_0$ and $p_1$ by solving ODEs. We refer the reader to \citet{Liu2024_FlowNotes,liu2022rectified} for an exhaustive treatment.  
We let $p_1 := \Law(X_1)$ be the target data distribution with its support $\cX \subseteq \bbR^d$, the data generating distribution from which we would like to draw samples. The linear-interpolation process is defined by by
\begin{equation}
\label{eq: RF linear process}
X_t = t X_1 + (1-t) X_0, \quad 0\leq t \leq 1
\end{equation}
where $p_0$ is the \textit{standard Gaussian} distribution. We will assume an independent initial coupling,  i.e., $(X_0, X_1) \sim p_0 \otimes p_1$ and let  $p_t = \Law(X_t)$, for any $t \in [0,1]$.
The RF procedure is a deterministic sampler arising from an ODE with {\it drift (or velocity) function} $v:\bbR^d \times [0, 1] \rightarrow \bbR^d$ defined as the solution to the optimization problem 
\begin{align}
    v &= \argmin_{f} \int \bbE{\Norm{X_1 - X_0 - f(X_t, t)}_2^2} dt , \label{drift_obj}
\end{align}
where the minimization is over all functions $f:\bbR^d \times [0, 1] \rightarrow \bbR^d$. 
For each $t$, The above objective in (\ref{drift_obj}) is minimized at
$
 (x,t) \mapsto   v_t(x) := v(x, t) = \bbE[X_1-X_0 \mid X_t = x]. 
$
 Under appropriate regularity conditions \citep[investigated in][]{bansal2025wasserstein, mena2025statisticalpropertiesrectifiedflow}, the ODE
\begin{align}
    dZ_t = v_t(Z_t)\ dt, \quad Z_0 \sim N(0, I_d) \label{eq: ode-true}
\end{align}
 satisfies the marginal preserving property,
 i.e., $\Law(Z_t) = \Law(X_t) = p_t$, owing to the Fokker-Planck equation
 \begin{equation}
     \label{eq: RF-FP}
     \frac{\partial p_t}{\partial t} + \nabla. (v_t p_t) = 0.
 \end{equation}
 Hence, ODE \eqref{eq: ode-true} can be used for sampling. 
 % We call $\cZ = \bc{Z_t}_{t \in [0,1]}$ the rectified flow of the coupling $(X_0, X_1)$, denoted as $\cZ = \Rectflow{(X_0, X_1)}$, and $(Z_0, Z_1)$ the rectified coupling, denoted as $(Z_0, Z_1) = \Rectify{(X_0, X_1)}$. 


%  The uniform Lipschitzness of the drift function $v_t$ for all $t \in [0, 1]$ is a sufficient condition for the rectified flow $\cZ$ to be unique \citep[Theorem 1]{murray2013existence}.  Hence, the \texttt{Rectflow} procedure rewires the trajectories of the linear interpolation process such that no two paths, corresponding to different initial conditions, intersect at the same time. After solving the ODE (\ref{eq: ode-true}), one can also apply another \texttt{Rectflow} procedure, also called Reflow or the 2-Rectified flow, to the coupling $(Z_0, Z_1)$ by learning the drift function 
% $$v^{(2)}_t(Z^{(2)}_t) = \bbE{Z_1 - Z_0 \mid tZ_1 + (1-t)Z_0 = Z^{(2)}_t}.$$ 
% This procedure can be done recursively, say $K$ times,  resulting the $K$-Rectified Flow procedure. 
% % We denote the joint distribution of $(Z^{(K)}_0, Z^{(K)}_1)$ as $\nu^{(K)}$. 
% \cite{liu2023flow} shows that $K$-Rectified Flow couplings are straight in the limit of $K \to \infty.$ We give the formal definition of straightness below.

% \begin{definition} (Straight coupling and flow)
% \label{def: straight flow}
%     A coupling $(X_0, X_1)$ is called straight or fully rectified when 
%     %its \rd linear interpolation paths do not intersect\bk, or equivalently, letting $t \sim  \text{Unif}\br{[0, 1]}$,
%     $\bbE{X_1 - X_0 \mid tX_1 + (1-t)X_0} = X_1 - X_0$ almost surely in $t \sim \Unif([0,1])$. %\red{(add definition and requirement of rectifiability)}

%     Moreover, for a straight coupling $(X_0, X_1)$, the corresponding rectified flow $\cZ = \Rectflow{(X_0, X_1)}$ has straight line trajectories, and $ (X_0, X_1) \overset{d}{=} (Z_0, Z_1) = \Rectify{(X_0, X_1)}$, i.e., they both have the same joint distribution. A flow that satisfies these properties is called a straight flow. 
% \end{definition}

% Straight flows are especially appealing because, in practice, solving the ODE \eqref{eq: ode-true} analytically is rarely feasible, necessitating the use of discretization schemes for numerical solutions. However, for straight flows, the trajectories follow straight lines, allowing for closed-form solutions without the need for iterative numerical solvers, which significantly accelerates the sampling process. 

To implement this procedure, one faces two key challenges. 
First, given a sample from the target distribution $p_1$, one has to estimate the drift function, which is naturally done by empirically minimizing the objective in (\ref{drift_obj}) over a large and expressive function class (e.g., U-Net). %and replaced in (\ref{eq: ode-true}). 
% Subsequently, the estimate $\widehat v_t \in \cF$ is used to obtain the sampling ODE
% \begin{align}
%     d \tilde Y_t = \widehat v_t(\tilde Y_t)\, dt , \quad  \tilde Y_0 \sim N(0, I_d) \label{eq: ode-approx}.
% \end{align}
% Because the solution to the ODE (\ref{eq: ode-approx}) is typically not analytically available, one must rely on discretization schemes.
Secondly, for sampling, one must rely on a time-discretization scheme of the ODE \eqref{eq: ode-true}. One such scheme is the Eulerian update rule, which, for a given discretization of the time course in $N$ intervals, with $0 = t_0 < t_1 < \dots < t_N =1$, evaluates 
\begin{align}
\label{eq: emp-ode-disc}
    Y_{t_{i+1}} =  Y_{t_{i}} + \eta_i\widehat v_{t_{i}}(Y_{t_{i}}); \; \eta_i :=  t_{i+1}-t_{i}, 0 \le i \le N-1,
\end{align}
where $Y_0 \sim N(0, I_d)$ and $\widehat v_t(\cdot)$ is the estimated velocity at time $t$. Let $\widehat p_1 := \Law(Y_1)$ denote the distribution of the final step. Other types of time discretized schemes of the ODE \eqref{eq: ode-true} can be chosen for the sampling step, and the resulting discretization errors impact the sampling fidelity differently. This has been the focus of a recent line of works \citep{benton2023nearly, chen2022sampling, li2024towards, li2024sharp, huang2024denoising, liang2025low} aimed at identifying time-discretized versions of DDPM and DDIM samplers to achieve tight convergence rates.

\paragraph{The velocity function.} 
Before proceeding, we illustrate the connection between the velocity $v_t(x)$ and the score function $s_t(x):= \nabla \log p_t(x)$. Recall that $X_0 \sim N(0, I_d)$ and independent of $X_1\sim p_1$. Therefore, Tweedie's formula \citep{efron2011tweedie, meng2021estimating} immediately yields $v_t(x) =  \frac{x}{t} + \frac{(1-t)}{t} s_t(x)$. Thus, estimating $v_t$ is essentially equivalent to estimating the score $s_t$ for $t \in (0,1)$. 
% , where $p_t$ is the density of $X_t$ w.r.t. the Lebesgue measure for $t\in [0,1)$
%  To this end, we define two key quantities:
% \[
% \mu_{1\mid t}(x) = \bbE[X_1 \mid X_t = x], 
% \; 
% \Cov_{1 \mid t}(x) = \Cov(X_1 \mid X_t = x).
% \]
% With this, Equation \eqref{eq: drift-true} can be re-written as
% \begin{equation}
% \label{eq: v_tweedie}
% v_t(x) = \bbE[X_1 - X_0 \mid X_t = x] = \frac{1}{1-t}\br{\mu_{1 \mid t}(x) - x}.
% \end{equation}

% By the first and second order Tweedie's formula \citep{efron2011tweedie, meng2021estimating}, it follows that, for any choice of $t \in (0,1]$ and $x \in \mathbb{R}^d$,
% \[
% \mu_{1 \mid t}(x) =   \frac{x}{t} + \frac{(1-t)^2}{t} s_t(x),\] 
% \[
% \Cov_{1 \mid t}(x) = \frac{(1-t)^2}{t^2}I_d + \frac{(1-t)^4}{t^2} \nabla s_t(x).
% \]
% Substituting this in \eqref{eq: v_tweedie}, we obtain that $v_t(x) = \frac{x}{t} + \frac{(1-t)}{t} s_t(x)$, and $\nabla v_t(x) = - \frac{I_d}{1-t} + \frac{t}{(1-t)^3} \Cov_{1 \mid t}(x)$. The reformulation of $v_t$ essentially shows that estimating $v_t$ is equivalent to estimating the score $s_t$ for $t \in (0,1)$. 

\subsection{Stochastic rectified flow (STOC-RF)}
Since the work of DDIM \cite{song2021denoising} and DDPM \cite{song2021scorebased}, it is well known that SDEs can be converted to ODEs to obtain deterministic samplers. Conversely, it is also possible to convert ODEs to its equivalent SDE form to obtain stochastic samplers. In fact, prior works \cite{albergo2023stochastic, xue2024unifying, Hu_2025_ICCV} have shown that stochastic samplers typically enjoy better generation quality. One reason behind the inferior performance of deterministic samplers is the possibility of error accumulation during the sampling step along the discrete trajectories stemming from inaccurate velocity estimates. In other words, if $Y_{t_i}$ differs significantly from $Z_{t_i}$ for some $t_i$, then the Euler step \eqref{eq: emp-ode-disc} accumulates error due to inaccurate evaluation of $\widehat v_{t_i}(Y_{t_i})$. Therefore, this can further reinforce the deviation in evaluating $Y_{t_{i+1}}$.
A way to mitigate the compounding error from successive Euler
steps is to introduce a stochastic Langevin correction at each step. Empirically, \citet{hu2025amo} showed that the resulting stochastic RF samplers   enjoy better text rendering quality over  deterministic samplers. We further illustrate this point with Figure \ref{fig: ODE vs SDE 2-GMM}, showing that the ODE-based RF trajectories appear to generate more outliers compared to the SDE-based STOC-RF sampler. 
\begin{figure}[h]
  \centering
  \begin{subfigure}[b]{0.45\textwidth}
    \centering
    \includegraphics[width=\textwidth]{figures/Deterministic_Sampler_ODE_sharp.pdf}
    %\caption{Caption for first subfigure}
    \label{fig:sub1}
  \end{subfigure}\hfill
  \begin{subfigure}[b]{0.45\textwidth}
    \centering
    \includegraphics[width=\textwidth]{figures/Stochastic_Sampler_SDE_sharp.pdf}
    %\caption{Caption for second subfigure}
    \label{fig:sub2}
  \end{subfigure}
  \caption{Trajectories of RF and STOC-RF samplers for a mixture of 2-Gaussian target distribution.}
  \label{fig: ODE vs SDE 2-GMM}
\end{figure}


Following e.g. \citet[Chapter 5]{Liu2024_FlowNotes}, the stochastic version of RF is specified by the SDE 
\begin{equation}
\label{eq: RF-SDE}
    d\tilde Z_t = v_t(\tilde Z_t)dt + \underbrace{\gamma_t s_t(\tilde Z_t)dt + \sqrt{2 \gamma_t} d \tilde B_t}_{\text{Langevin correction}}; \; \tilde Z_0 = Z_0,
\end{equation}
for $t \in [0,1]$. This is obtained from adding a Langevin correction to the original RF ODE \eqref{eq: ode-true}.
Above, $\gamma_t\geq  0$ is a possibly time-varying diffusion coefficient, and $\{\tilde B_t\}_{t\ge 0}$ is the standard Brownian motion in $\bbR^d$. %Classical SDE results \citep[Section 4.3]{gardiner2004handbook} 
One can check (see, e.g.,  \citet[Section 5.2]{Liu2024_FlowNotes} or \citet{song2021scorebased}) that the marginal $\tilde p_t := \Law(\tilde Z_t)$ satisfies  the continuity equation 
\begin{equation}
    \label{eq: RF-SDE-FP}
    \frac{\partial \tilde p_t}{\partial t} + \nabla. (\tilde v_t \tilde p_t) = 0,
\end{equation}
where $\tilde v_t(x) = v_t(x) + \gamma_t s_t(x) - \gamma_t \nabla \log \tilde p_t(x)$.   Since the continuity equations  \eqref{eq: RF-FP} and \eqref{eq: RF-SDE-FP} coincide, $\tilde p_t = p_t$ , i.e., $\Law(\tilde Z_t) = \Law(Z_t)$, for all $t$, demonstrating that the SDE \eqref{eq: RF-SDE} can indeed be used for sampling from $p_1$. In practice, this is accomplished via an appropriate time discretization scheme
\begin{equation}
    \label{eq: RF-SDE-disc}
    \tilde Z_{t_{i+1}} = \tilde Z_{t_i} + \tilde \eta_i \bc{v_{t_i}(\tilde Z_{t_i}) +  \gamma_{t_i} s_{t_i}(\tilde Z_{t_i})} + \sqrt{2 \tilde \eta_i \gamma_{t_i}} \xi_{t_i},
\end{equation}
where $\bc{t_i}_{i=0}^N$ is an appropriate schedule of times, the $\xi_{t_i}$'s are i.i.d. draws from $N(0, I_d)$, and the $\tilde \eta_i$'s are (typically small) step-sizes. The Langevin term essentially acts as a course-corrector to the trajectory of \eqref{eq: RF-SDE-disc} to adjust the trajectory distribution closer to $p_t$. Finally, one has freedom to choose the diffusion coefficients $\gamma_t$. In this paper we will restrict ourselves to the choice $\gamma_t = (1-t)/t$, which, as we will see below, is the \textit{appropriate noise scaling that yields equivalence between the STOC-RF SDE \eqref{eq: RF-SDE} and DDPM through the lens of Stochastic localization}. This also addresses an question raised during the ICML tutorial talk  \citet{liu2025flowing} regarding the motivation of such choice.
%As mentioned in the ICML talk \citet{liu2025flowing}, the theoretical rationale behind the choice $\gamma_t = \frac{1-t}{t}$ remained  elusive. We show that this is a natural choice that allows us to have the equivalence between STOC-RF, DDPM and Stochastic localization.



\section{Stochastic localization, RF and DDPM}
\label{sec: SL}
%\subsection{Stochastic localization: A unifying framework}
%\label{sec: SL intro}
\paragraph{Stochastic localization (SL).} 
SL  is a probabilistic technique to study the properties of high-dimensional distributions due to \citet{eldan2013thin,eldan2020taming} that provides a powerful, unifying framework for understanding, analyzing, and improving score-based generative models and sampling algorithms. See  \citet{el2022information}, \citet{montanari2023sampling} and the survey by \citet{shi2025perspectivesstochasticlocalization}. To sample from the target data distribution $p_1$, SL focuses on an appropriate measured-valued stochastic process $\{\nu_s\}_{s\ge 0}$ such that $\nu_s$ ``localizes'' towards $p_1$ as $s \to \infty$. A straightforward way to formulate this  process is by letting, for any $s \geq 0$, $\nu_s := \Law (X_1 \mid U_s)$, where $X_1 \sim p_1$ and 
\begin{equation}
\label{eq: SL process}
    U_s = s X_1 + B_s,
\end{equation}
with $\{B_s\}_{s \ge 0}$  a standard Brownian motion in $\bbR^d$ \citep{montanari2023sampling}, independent of $X_1$. Thus, $\nu_s$ is a random measure depending on $U_s$. Note that $U_s/s \to X_1$ almost surely as $s \to \infty$, i.e., $\nu_s$ does indeed localize around $p_1$. 

%\red{(SR: irrelevant paragraph)} 
While enlightening, the construction \eqref{eq: SL process} does not yield a sampler from $p_1$, as it uses $X_1 \sim p_1$ in the definition of $U_s$. Therefore, it is instead convenient to consider the alternative process $\{\tilde U_s\}_{s \geq 0}$ defined as the solution of the  SDE
\begin{equation}
    \label{eq: SL SDE}
    d\tilde U_s = a_s(\tilde U_s) + d \tilde B_s,
\end{equation}
where $\tilde B_s$ is an independent Brownian motion and, for any $s \geq 0$,  $a_s(u) = \bbE[X_1 \mid U_s = u]$. By classical SDE results -- see \citet[Theorem 8.4.3 ]{oksendal2003stochastic} or \citet[Theorem 7.12 ]{liptser2013statistics} -- $\{U_s\}_{s\ge 0}$ and $\{\tilde U_s\}_{s\ge 0}$ are equivalent in law. Therefore, \eqref{eq: SL SDE} can be used for sampling from $p_1$ as long as we have access to the regression functions $\{ a_s(\cdot), s \geq 0\}$, which can be estimated from the data. We also define the posterior covariance as $\bA_s = \Cov(X_1 \mid U_s)$. The next lemma from \citet{eldan2020taming} allows us to control the growth of the expected posterior covariance.
\begin{lemma}
    \label{lemma: eldan ddt cov lemma}
    For all $s\ge 0$, $\frac{d}{ds} \bbE\bs{\bA_s} = - \bbE\bs{\bA_s^2}$.
\end{lemma}
%This yields a very important corollary.
This is a crucial result for controlling the discretization error of ODE and SDE samplers of RF.

\paragraph{SL and DDPM.}
DDPMs are a class of diffusion-based sampling mechanism for generative AI introduced by \citet{ho2020denoising}.  DDPM first starts with a \textit{forward process} $(Y'_\tau)_{\tau\ge 0}$ starting from $Y_0' = X_1 \sim p_1$ and  evolving according to the SDE
\begin{equation}
    \label{eq: forward process}
    dY_\tau' = - \beta(\tau) Y_\tau' d\tau + \sqrt{2 \beta(\tau)} dB^\prime_\tau, \; Y_0' \sim p_1,
\end{equation}
where $\beta(\tau)$ is the diffusion coefficient and $\bc{B^\prime_\tau}_{\tau\ge 0}$ is standard Brownian motion. If $\beta(\tau) = 1$, then the SDE \eqref{eq: forward process} reduces to the classical Ornstein-Ulhenbeck (OU) process.
Typically, The forwards process is continued till a large time $\tau = N$ such that $Y_N'$ is approximately Gaussian. Let $q_\tau(y_\tau')$ be the marginals of $y_\tau'$ following \eqref{eq: forward process}. Then in order to sample from the target distribution $p_1$, a \emph{reverse process} $\{\overleftarrow Y_\tau\}_{\tau \in [0,N]} $ is formulated, starting at $\overleftarrow Y_0 = Y_N'$ and evolving according to the reverse-time SDE  
\begin{equation}
\label{eq: reverse OU process}
\begin{aligned}
  d\overleftarrow Y_{\tau} =  & \{\overleftarrow Y_\tau + 2 \nabla \log q_{N - \tau}(\overleftarrow y_\tau)\} \beta(N-\tau)d\tau\\
  & + \sqrt{2 \beta(N-\tau)} d \overleftarrow B_\tau,
  \end{aligned}
\end{equation}
where $\{\overleftarrow B_\tau\}_{\tau \ge 0}$ is a standard Brownian motion.
Classical results in SDE \cite{anderson1982reverse, haussmann1986time} show that $\overleftarrow Y_\tau \overset{d}{=} Y_{N - \tau}'$ which allows us to generate samples from $p_1$ via simulating the backward SDE \eqref{eq: reverse OU process}.

It is known that DDPMs and stochastic localization are equivalent \citep{montanari2023sampling}. 
In particular, \citet{benton2023nearly}  showed that the forward process \eqref{eq: forward process} is \textit{equivalent in law} to the SL process \eqref{eq: SL process} under an appropriate time change. In this section, we restate this result in a slightly different fashion that will allow us to easily see the connection between DDPM and RF.   
% Therefore, we can conclude that 
% \[
% \frac{ Y_\tau^\gets}{\sqrt{\bar{\alpha}_\tau}} \overset{d}{=} X_1 + B_{\frac{1 -\bar{\alpha}_\tau}{\bar \alpha_\tau}}^\prime.
% \]
To this end, we define  $\omega_{\tau} := \exp(-2 \int_0^\tau \beta(u) du)$. Then the following lemma shows that the the \textit{forward process \eqref{eq: forward process} and the SL process \eqref{eq: SL process} are equivalent:}
\begin{lemma}
\label{lemma: time change DDPM-SL}
    Let $\{\tau(s)\}_{s \ge 0}$ be sequence such that $\frac{1 -\omega_{\tau(s)}}{\omega_{\tau(s)}} = \frac{1}{s}$. Then, we have 
    %\begin{equation}
    %\label{eq: time change DDPM-SL}
    $\bc{\frac{Y'_{\tau(s)}}{\sqrt{\omega_{\tau(s)}}}}_{s\ge 0} \overset{d}{= } \bc{\frac{U_s}{s}}_{s \ge 0}.$
%\end{equation}
\end{lemma}
If $\beta(\tau) = 1$ for all $\tau$ (OU process), then the time transformation in Lemma \ref{lemma: time change DDPM-SL} simplifies to be $\tau(s)  = \frac{1}{2}\log(1 + s^{-1})$, and we get $\bc{\frac{Y'_{\tau(s)}}{e^{- \tau(s)}}}_{s\ge 0} \overset{d}{= } \bc{\frac{U_s}{s}}_{s \ge 0}$. Note, this matches the findings in \citet[Section 1.2]{benton2023nearly}. The proof is deferred to Appendix \ref{app: SL = RF = DDPM}.


\paragraph{SL, RF and STOC-RF.} Next, we will also show that the probability path of RF is also equivalent to the SL process \eqref{eq: SL process}. Recall that that RF starts by constructing the linear process $X_t = t X_1 + (1-t)X_0$, where $X_1 \sim p_1, X_0 \sim N(0, I_d)$ and $t \in [0,1]$. To this end, we consider the process 
\begin{equation}
    \label{eq: RF linear stoc process}
    \tilde X_t = t X_1 +  t W_{(1-t)^2/t^2}, \quad t \in [0,1]
\end{equation}
where $\{W_\tau\}_{\tau \ge 0}$ is a Brownian motion. Note that $\Law(\tilde X_t) = \Law (X_t)$, for all $t \in [0,1]$. In fact, as the next result show, the entire process \eqref{eq: RF linear stoc process} has the same law of the SL process \eqref{eq: SL process}, up to time-change.
\begin{lemma}
\label{lemma: SL and RF are equivalent}
    Let $\{t(s)\}_{s\ge0}$ be a sequence such that $\br{\frac{1 - t(s)}{t(s)}}^2 = \frac{1}{s}$. Then, we have 
    $
        %\label{eq: SL and RF are equivalent}
        \bc{\frac{\tilde X_{t(s)}}{t(s)}}_{s \ge 0} \overset{d}{=} \bc{\frac{U_s}{s}}_{s \ge 0}.
    $
\end{lemma}
The proof is in Appendix \ref{app: SL = RF = DDPM}. Crucially, \Cref{lemma: SL and RF are equivalent} yields the following Corollary to \Cref{lemma: eldan ddt cov lemma} which is used to control the discretization error of the RF sampler.
\begin{corollary}
    \label{cor: my ddt cov version}
    Let $\bSigma_t := \Cov(X_1\mid X_t)$ for $t \in [0,1)$. Then $\frac{d}{dt} \bbE (\bSigma_t) = - \frac{2t}{(1-t)^3} \bbE(\bSigma_t^2)$.
\end{corollary}

See \Cref{sec: RF covariance} for complete proof. Additionally, \Cref{lemma: time change DDPM-SL} and \Cref{lemma: SL and RF are equivalent} leads to the following:
\begin{equation}
    \label{eq: DDPM and RF equivalent}
    \bc{\frac{Y^{'}_{\tau}}{\sqrt{\omega_{\tau}}}}_{\tau \ge 0} \overset{d}{=} \bc{\frac{\tilde X_{t(\tau)}}{t(\tau)}}_{\tau \ge 0},
\end{equation}
where $t(\tau) = \frac{\sqrt{\omega_\tau}}{\sqrt{\omega_\tau} + \sqrt{1 - \omega_{\tau}}}$. In fact, under the same time-change the STOC-RF process \eqref{eq: RF-SDE} and the backward process \eqref{eq: reverse OU process} are equivalent.
\begin{proposition}
    \label{prop: RF-SDE and backwrd process are equivalent}
    Let $t(\tau) = \frac{\sqrt{\omega_\tau}}{\sqrt{\omega_\tau} + \sqrt{1 - \omega_{\tau}}}$, and $\{\tilde Z_t\}_{t\ge 0}$ be a solution to the SDE \eqref{eq: RF-SDE}. Then,  
    $\bc{\frac{\sqrt{\omega_{N-\tau}}\tilde Z_{t(N - \tau)}}{t(N - \tau)}}_{\tau\ge 0}$ is a solution to the SDE \eqref{eq: reverse OU process}. Conversely, if $\overleftarrow Y_\tau$ is a solution of SDE \eqref{eq: reverse OU process}, then $\bc{\frac{ \overleftarrow Y_{\tau(t)}}{\sqrt{\omega_{N-\tau(t)}} + \sqrt{1 - \omega_{N - \tau(t)}}}}_{\tau\ge 0}$ is a solution of SDE \eqref{eq: RF-SDE}, where $\tau(t)$ is the unique solution of the integral equation $\int_{0}^{N-\tau} \beta(u) du = \log\br{\sqrt{1  + \frac{(1-t)^2}{t^2}}}$.%\red{Initial condition is different. fix it.}
 \end{proposition}
The above proposition is very useful in analyzing the convergence rate of STOC-RF. Essentially, Proposition \ref{prop: RF-SDE and backwrd process are equivalent} allows uso to transit between DDPM and STOC-RF under simple time change and scaling. Therefore, convergence properties of one process can be carried over to the other one with the same transformation. This is the key insight that we will use to prove our convergence result of STOC-RF in later sections. The detailed proof is deferred to Appendix \ref{app: equivalence of backward and stoc-RF}.


\paragraph{SL and Stochastic Interpolants.}
\label{sec: RF and deterministic interpolants}
%\red{SR: the sentence doesn't look right..}
Lemma  \ref{lemma: SL and RF are equivalent} and the equivalence result in \eqref{eq: DDPM and RF equivalent} can be extended to cover general stochastic interpolation \cite{albergo2023building, albergo2023stochastic} schemes beyond RF. %\cite{albergo2023stochastic, albergo2023building}. 
In detail, instead of a linear interpolation, one can consider a general deterministic affine interpolation process $I_\theta = a_\theta X_1 + b_\theta X_0$, where $a_{(\cdot)}$ and $b_{(\cdot)}$ are smooth functions of $\theta \in [0,1]$, with $a_0 = b_1 = 0$, and $a_1 = b_0 = 1$. When $X_0 \sim N(0,I_d)$, the associated stochastic interpolant is the process $(\tilde I_\theta)_{\theta \in [0,1]}$, where $\tilde I_\theta = a_\theta X_1 + a_\theta W_{r^2_{\theta}}$, with $r_\theta = b_\theta/a_\theta$ and $\bc{W_t}_{t\ge 0}$  a Brownian motion. It is easy to see that $\tilde I_\theta \overset{d}{=} I_\theta$, for each $\theta$. Moreover, each stochastic interpolant can be rescaled and time-transformed to return the SL process, in a sense made precise by the next result. %\textcolor{violet}{ALE: state as a result?}
\begin{proposition}
\label{prop: SL and general interpolant are equivalent}
  $\bc{\frac{\tilde I_{\theta(s)}}{a_{\theta(s)}}}_{s\ge 0 } \overset{d}{=} \bc{\frac{U_s}{s}}_{s\ge 0}$, with $r^2_{\theta(s)} = \frac{1}{s}$.  
\end{proposition}


In particular, in virtue of Lemma \ref{lemma: SL and RF are equivalent}, RF is equal in law to any interpolant process $I_\theta$ under the time change $t(\theta) = \frac{a_\theta}{a_\theta + b_\theta}$. This equivalence among stochastic interpolants is not new, as it has also been shown in \citet{Liu2024_FlowNotes}  by direct arguments. Connection between diffusion and flow matching has also been studied in \citet{albergo2023stochastic, albergo2022building,gao2025diffusionmeetsflow, ma2024sit, rout2025semantic} through different perspectives.
%Our framework allows us to re-discover this by identifying the connection between $I_\theta$ and  stochastic localization. 
By leveraging the representation properties of the SL process, we establish this result in a more unified manner, which has the added benefit of directly applying to both stochastic interpolants and SDEs. The proof of \Cref{prop: SL and general interpolant are equivalent} is deferred to Appendix \ref{app: proof of SL and general interpolants equivalence}.

\section{Main Results}
\label{sec: main results}
In this section, we present our main convergence results for both deterministic and stochastic samplers. We  begin by introducing and commenting on our key assumptions. Throughout, $\bc{t_i}_{i=0}^N$ refers to the time schedule for both the RF and STOC-RF procedures, as in  \eqref{eq: emp-ode-disc} and \eqref{eq: RF-SDE-disc}, respectively, and $N$ is the number of time steps.

% \subsection{Key assumptions and quantities}
\subsection{Low-dimensionality of target distribution}
To formalize the notion of low-dimensionality of $p_1$, we turn to the geometric complexity of its support $\cX$ as measured by its metric entropy. 

\begin{assumption}[Low-dimensionality]
\label{assumption: intrinsic dimension}
    Let $\epsilon_0 = N^{-c_{\epsilon_0}}$ for some sufficiently large $c_{\epsilon_0}>0$. The $\epsilon_0$ covering number of $\cX$ with respect to the Euclidean norm, $\cN(\cX, \Norm{\cdot}_2, \epsilon_0)$, satisfies
    $
    \log \cN(\cX, \Norm{\cdot}_2, \epsilon_0) \le C_{\rm cover} k \log N,
    $
    for some constant $C_{\rm cover}>0$ and integer $k$. The quantity $k$ is referred as the intrinsic-dimension of $\cX$.
\end{assumption}
The above assumption was introduced in recent works by \citet{huang2024denoising,li2024adapting,liang2025low} on adaptation of diffusion methods to intrinsic low dimensionality of the target function. Notably, by increasing the number of steps $N$, the assumption is less and less stringer.  Informally, Assumption \ref{assumption: intrinsic dimension} requires the support of $p_1$ to be concentrated \textit{on or near} a $k$-dimensional set.
The above definition of intrinsic dimension is fairly general, as it covers a variety of important low-dimensional structures. Examples satisfying Assumption \ref{assumption: intrinsic dimension} include $k$-dimensional linear subspace in $\bbR^d$ and $k$-dimensional non-linear manifolds (provided $\cX$ is polynomially bounded as in Assumption \ref{assumption: bounded support}). Structures with doubling-dimension $k$ \citep{dasgupta2008random, kpotufe2012tree} also satisfies Assumption \ref{assumption: intrinsic dimension}. More detailed discussion on this can be found in Section 4.1.1 of \citet{huang2024denoising}. Moreover, in practice, target distribution (e.g., image data) is supported on a lower dimensional manifold and the intrinsic dimension $k \ll d$ \citep{facco2017estimating, ansuini2019intrinsic, pope2021intrinsic}.
%\textcolor{violet}{ALE: Clarify that Assumption 4.1  says that the support of  $p_1$ is concentrated on or near a $k$-dimensional set (near because $\epsilon_0$ does not tend to $0$). Later on we should also mention that RF or STOC-RF generate samples not from $p_1$ but from a distribution that has full dimensional support and is ```near" the spoort of $p_1$. This is an impprtant point poorly addressed by \cite{liang2025low}. In fact, there cannot be a vanishing TV rate if the target distribution is truly low-dimensional (because $p_0$ and $p_1$ would be singular).}
 %In the next assumption, we would like to impose boundedness of of the support $\cX$.
\begin{assumption}[Bounded support]
    \label{assumption: bounded support}
    There exists a universal constant $c_R>0$ such that 
    $
    \sup_{x \in \cX} \Norm{x}_2 \le R, \text{where $R = N^{c_R}$}.
    $
\end{assumption}
The bounded support assumption is not as stringent as it may appear at first, as the diameter of $\cX$ is allowed to scale polynomially (with arbitrarily large degree) in the
number of iterations of the sampler, which accommodates a wide range of practical applications, such as
image generation.
\subsection{Assumptions on velocity approximation error}
Our next set of assumptions are concerned with the approximation error of the estimated velocity field $\hat v_{t}$. 
% \red{
% 1. Low dimensionality of target distribution\\
% 2. Score approximation error.}
 
\begin{assumption}[Velocity approximation]
\label{assumption: velocity approximation}
    Define $(\evel_i)^2:= \bbE \Norm{v_{t_i}(X_{t_i}) - \widehat{v}_{t_i}(X_{t_i})}_2^2$. Then, 
    $\frac{1}{N} \sum_{i =1}^N (\evel_i)^2 \le \meanevel^2.$
\end{assumption}
The above assumption essentially requires a good control on the mean-squared error of the estimated velocity field. Similar assumptions have been considered in \citet{bansal2025wasserstein, guan2025mirror} to prove Wasserstein convergence of RF. In the context of score-matching algorithms, similar assumptions related to the score-approximation error have also appeared in prior works such as \citet{benton2023nearly, chen2022sampling, li2024sharp, liang2025low}. 

% Finally, to derive convergence rates of the deterministic samplers of RF that are adaptive to the intrinsic dimension of the support of the target distribution, we require the velocity estimator $\widehat v_t$ to have higher order accuracy guarantees.

\begin{assumption}[Higher-order approximations]
    \label{assumption: higher-order approximation error}
   $\;$% We assume the following:
    \begin{enumerate}[label=(\alph*)]
        \item \label{assumption: Jone} Let $(\evelJone_i)^2 := \bbE \Norm{\nabla v_{t_i}(X_{t_i}) - \nabla \widehat v_{t_i}(X_{t_i})}_F^2$. Then, $N^{-1} \sum_{i =1}^N (\evelJone_i)^2 \le \meanevelJone^2$. Also, we assume $\sup_{x \in \bbR^d} \eta_i \Norm{\nabla v_{t_i}(x) - \nabla \widehat v_{t_i}(x)}_{\op} \le 1/8$ for all $0 \le i \le N-2$.

        \item \label{assumption: Jtwo} Let $(\evelJtwo_i)^2 := \bbE [\Tr(\nabla v_{t_i}(X_{t_i}) - \nabla \widehat v_{t_i}(X_{t_i}))^2]$. Then, $N^{-1} \sum_{i =1}^N (\evelJtwo_i)^2 \le \meanevelJtwo^2$.

        \item \label{assumption: Hessian} Let $(\evelH_i)^2 := \bbE \Norm{\nabla \Tr(\nabla v_{t_i}(X_{t_i}) - \nabla \widehat v_{t_i}(X_{t_i}))}_2^2$. Then, $N^{-1} \sum_{i =1}^N (\evelJtwo_i)^2 \le \meanevelH^2$.
    \end{enumerate}
\end{assumption}
Nearly identical assumptions have been used in \citet{liang2025low} in the context of DDIM. Intuitively, an adequate control of the higher-order approximation terms mitigates the propagation of  error during sampling, as the deterministic samplers are unable to self-correct the path of the flow. 
%In contrast, Assumption \ref{assumption: higher-order approximation error} is not required to prove convergence guarantees for STOC-RF sampler as the Langevin term automatically acts as a path corrector term.
\subsection{Convergence rate of rectified flow}
% In this section, we will establish the main convergence theorem for RF sampler \eqref{eq: emp-ode-disc}.
%\subsubsection{Eulerian scheme}
% For a given time-discretization $0= t_0 < t_1<\cdots < t_N = 1$, we define $\eta_i = t_{i+1} - t_i$. 
%We recall the Euler's scheme \eqref{eq: emp-ode-disc}:
% $
% Y_{t_{i+1}} = Y_{t_i} + \eta_i \widehat v_{t_i}(Y_{t_i}), \; Y_0 \sim N(0, I_d).
% $
\paragraph{Time partition.} Our convergence guarantees for the RF sampler \eqref{eq: emp-ode-disc} depend crucially on a carefully constructed U-shaped time schedule $\{t_i\}_{i = 0}^N$, which we describe next.
For given, user-specified values $h\in (0,1)$ and $\delta \ge 0$, we consider the non-uniform time partition 
%\vspace{-1em}
% \begin{equation}
% \label{eq: gemoetric time discretization}
%     \begin{aligned}
%       & \text{
%     1. $t_0 = 0, t_1 =\delta, t_j = (1+h)t_{j-1}$ for $2\le j \le N/2$,
%     }\\ 
%     & \text{
%     2. $t_N = 0, t_{N-1} = 1- \delta, t_j = 1 - (1+h)(1-t_{j+1})$
%     }\\
%     & \quad \text{ for $N/2< j \le N-2$}.
%     \end{aligned}
% \end{equation}
\begin{equation}
\label{eq: gemoetric time discretization}
t_j =
\begin{cases}
0, & j = 0, \\[4pt]
\delta, & j = 1, \\[6pt]
(1+h)\, t_{j-1}, & 2 \le j \le \dfrac{N}{2}, \\[8pt]
1 - (1+h)\,(1 - t_{j+1}), & \dfrac{N}{2} < j \le N-2, \\[6pt]
1 - \delta, & j = N-1, \\[6pt]
1, & j = N .
\end{cases}
\end{equation}
 We set $h$ such that $t_{N/2} = \delta (1+h)^{(N-2)/2} = 1/2$. Using $N-2 > N/2$ (for $N>4$), we arrive at 
\begin{equation}
\label{eq: value of h}
\frac{h}{2} \le \log(1+h) = \frac{2\log (\frac{1}{2 \delta})}{N-2} \Rightarrow  h \le \frac{8 \log (\frac{1}{2 \delta})}{N}.
\end{equation}
 We remark that the proposed time schedule is practicable, as it does not depend on any properties of $p_1$.
Typically, $\delta$ is chosen to be small, as
this give rise to a U-shaped time discretization as shown in Figure \ref{fig: hist of geom timegrid}.  This type of choice is not unusual \cite{lee2024improvingtrainingrectifiedflows}, as it yields more accurate estimates  $\widehat v_t$ during training when $t \approx 0$ and $t \approx 1$. Prior empirical work by %\citet{lee2024improvingtrainingrectifiedflows} has shown that such U-shaped choice of partition  results in better training of the RF model. 
non-uniform time discretizations have also been considered by \citet{huang2024denoising, li2024adapting, liang2025low} in the context DDPM and DDIM to establish low-dimensional adaptation of the models. It is important to emphasize that the specific form of the time schedule \eqref{eq: gemoetric time discretization}, while conforming to best empirical practices, is theoretically grounded, and it was derived directly as a result of our technical analysis of the RF dynamics (see Step 7 of the proof of Theorem \ref{thm: RF - convergence of Eulerian scheme}, Appendix \ref{app: RF convergence}). We now present the main convergence result for the RF sampler. 
%The detailed proof of the result is deferred to Appendix \ref{app: RF convergence}.
\begin{figure}
    \centering
    \includegraphics[width=0.45\textwidth]{figures/N_200_geom_timegrid_histogram.pdf}
    \caption{Histogram and kernel density estimation (KDE) plot of time-grid \eqref{eq: gemoetric time discretization} showing U-shaped distribution.}
    \label{fig: hist of geom timegrid}
\end{figure}
\begin{theorem}
\label{thm: RF - convergence of Eulerian scheme}
   Let Assumptions \ref{assumption: intrinsic dimension}-\ref{assumption: higher-order approximation error} hold, and write $\mu_1 = \bbE_{X \sim p_1}[X]$. Let $\{Y_{t_i}\}_{i\ge0}$ denote the RF updates of \eqref{eq: emp-ode-disc} with time schedule \eqref{eq: gemoetric time discretization}. Then, for $\delta = 1/(N \vee d)$, we have
   \begin{equation}
   \label{eq: RF convergence Euler scheme}
    \begin{aligned}
    & \TV(p_{X_{t_{N-1}}}, p_{Y_{t_{N-1}}})\\
    & \lesssim   \frac{k \log^3 (\frac{1}{\delta})  }{N} +   \frac{\log^2 (\frac{1}{\delta})}{N}   \meanevelJone^2 +  \frac{\log^2 (\frac{1}{\delta})}{N} \meanevel^2 \\
    & \quad +  \log \br{1/\delta} (\meanevel + \meanevelJone + \meanevelJtwo + \meanevelH) \\
     &\quad +    \delta^2 \bbE \Norm{X_1 - \mu_1}_2^2
     +  \delta^2 (\evel_0)^2  \\
     &\quad  +  \delta^2 (\evelJone_0)^2 +  (\evel_0 + \evelJone_0 + \evelJtwo_0 + \evelH_0) +  \delta^{-2}N^{-10} .  
\end{aligned}
\end{equation}
\end{theorem}

%Some commentary on the above rate of convergence of the RF sampler is in order.

%The  above convergence rates depend on several quantities. Setting aside the higher order approximation terms for the velocity field $\widehat{v}_t$ from \Cref{assumption: higher-order approximation error}, the leading terms 



\textbf{Adaptation to low-dimensionality and accelerated convergence.}  The rate in Theorem \ref{thm: RF - convergence of Eulerian scheme} does not contain a direct dependence on the ambient dimension $d$ and instead scales linearly in the intrinsic dimension $k$ of $\cX$, indicating that the RF sampler adapts to the low-dimensional structure of the target distribution. Thus, in principle, assuming perfect knowledge of the velocity drift (i.e. $\widehat v_t = v_t$ and the higher approximation terms of \Cref{assumption: higher-order approximation error} are all zero), the RF sampler \eqref{eq: emp-ode-disc} needs $\tilde O(k/\epsilon)$ iterations to achieve the guarantee $\TV(p_{X_{t_{N-1}}}, p_{Y_{t_{N-1}}}) < \epsilon$. If the intrinsic dimension $k \ll d$, then the RF sampler automatically accelerates without any prior knowledge of the low-dimensionality of $p_1$. To the best of our knowledge, this is the first work that provides a convergence guarantee of RF that adapts to the low-dimensionality of $p_1$. 

\textbf{The dependence on $\bbE \Norm{X_1 - \mu_1}_2^2$.} %Note first leading term on the right-hand side of \eqref{eq: RF convergence Euler scheme} only scales linearly with $k$. Therefore the RF sampler \eqref{eq: emp-ode-disc} automatically adapts to the low-dimensionality of $p_1$.  
The upper bound in Theorem \ref{thm: RF - convergence of Eulerian scheme} contains a term involving $\bbE\Norm{X_1 - \mu_1}_2^2$, which even for data supported on manifolds could scale linearly in the ambient dimension $d$. However, in practice one can choose the free parameter $\delta$ to be e.g. $O(d^{-1})$, so that the term becomes negligible. With this choice of $\delta$, last term in \eqref{eq: RF convergence Euler scheme} scales as $d^2 N^{-10}$ which is also negligible in practice. 
%For example, if $d = 256 \times 256$ and $N = 80$, then $d^2 N^{-10} \approx 4\times 10^{-10}$. %It is also worth mentioning that the term can be removed if we set $Y_{t_1} \sim \Law(X_{t_1})$ 

\textbf{Guarantee on perturbed data.} The \TV~guarantee in Theorem \ref{thm: RF - convergence of Eulerian scheme} is on the penultimate update $Y_{t_{N-1}}$ of \eqref{eq: emp-ode-disc}, and not on $Y_{t_N} = Y_1$. It is worthwhile to mention that $\TV(p_{X_1}, p_{Y_{t_N}})$ might not be a useful or even meaningful quantity, as $X_1$ and $Y_1$  may have  supports of different dimensions. For example, $X_1$ could be supported on a low-dimensional space, whereas $Y_{t_N}$ would typically have full dimensional support, thereby rendering $\TV (X_1, Y_{t_N}) =1 $. However, note that for any $\delta > 0$, $X_{\delta} \overset{d}{=} (1 - \delta) X_1 + \delta X_0$ has full-dimensional  support due to the slight Gaussian perturbation. Therefore, the \TV~distance in Theorem \ref{thm: RF - convergence of Eulerian scheme} is well defined.


\textbf{Higher-order approximation error terms.} The upper bound in Theorem \ref{thm: RF - convergence of Eulerian scheme} is heavily dependent on the higher-order approximation errors of $\widehat v_t$. These terms are required to be small as deterministic samplers are unable to self-correct their trajectories. In addition, at $t = 0$, the true drift $v_0(x) = \mu_1 - x$ can be estimated by $\widehat v_0(x) = n^{-1} \sum_{i \in [n]} {X_1^{(i)}} - x$, where $\{X_1^{(i)}\}_{i\in [n]} \overset{i.i.d.}{\sim } p_1$. With this choice of $\widehat v_0$, we have $\evelJone_0 = \evelJtwo_0 = \evelH_0 = 0$, and $\evel_0 = \sqrt{\bbE \Norm{X_1 - \mu_1}_2^2/n} $.

% \paragraph{Proof sketch:}
% \red{Write proof sketch. (depends on pages left (limit - 8 pages, currently running - 9 pages ), might remove later)}

\begin{comment}
\subsubsection{Non-Eulerian scheme}
In this section, we will propose a non-Eulerian update rule that also adapts to the low-dimensionality of $\cX$. Here, we will exploit the equivalence between RF and spherical interpolants (Section \ref{sec: RF and deterministic interpolants}). We begin with the spherical interpolation process
\[
V_\tau = \sqrt{\omega_\tau} X_1  + \sqrt{1 - \omega_\tau} X_0, \quad X_0 \sim N(0, I_d),
\]
where $\omega_\tau = \prod_{j =1}^\tau \alpha_j$, and for $c_0, c_1 >0$
\begin{equation}
    \label{eq: DDPM time scheduling}
    \begin{aligned}
        %& \beta_0 : = 1 - \alpha_0 = 1,\\
        & \beta_1 := 1 - \alpha_1 = \frac{1}{N^{c_0}},\\
        & \beta_{\tau+1} := 1 - \alpha_{\tau+1}\\
        &  = \frac{c_1 \log N}{N} \min \bc{\beta_1 \br{1 + \frac{c_1 \log N}{N}}^\tau, 1}, \quad 1 \le \tau \le N.
    \end{aligned}
\end{equation}
Similar time schedule has also been used in \cite{liang2025low, li2024sharp, li2024adapting} for the convergence analysis of DDPM and DDIM. In light of Section \ref{sec: RF and deterministic interpolants}, we note that $\bc{V_\tau}_{\tau \ge 0}$ and linear process $\bc{X_t}_{t \ge 0}$ (see \eqref{eq: RF linear process}) are equivalent under the time change $t(\tau) = \frac{\sqrt{\omega_\tau}}{\sqrt{1 - \omega_\tau} + \sqrt{\omega_\tau}}$, and $\frac{V_\tau}{\sqrt{\omega_\tau}} \overset{d}{=} \frac{X_{t(\tau)}}{t(\tau)}$. Next, we begin with the DDIM sampler of $\bc{V_\tau}_{\tau\ge 0}$ considered in \cite{liang2025low}:  

\begin{equation}
    \label{eq: DDIM sampler liang2025}
    \begin{aligned}
    & \hat V_{\tau -1} = \frac{1}{\sqrt{\alpha_\tau}} \bc{ \hat V_\tau + \delta_\tau \widehat s_{V_\tau}(\hat V_\tau)}, \hat V_N \sim N(0, I_d)\\
    & \delta_\tau = \frac{1 - \alpha_\tau}{1 + \sqrt{\frac{\alpha_\tau - \omega_\tau}{1 - \omega_\tau}}}, \quad \tau   = N , N-1, \ldots, 1,
    \end{aligned}
\end{equation}
where $ \widehat s_{V_\tau}$ is an approximate score of $V_\tau$.
To make a sense of the above sampler , we recall that $1 - \alpha_t = \beta_t \approx 0$. Hence, $1/\sqrt{\alpha_\tau}  \approx 1 + \beta_\tau/2$, and $\delta_\tau/ \sqrt{\alpha_\tau} \approx \beta_\tau / 2$. Substituting this in \eqref{eq: DDIM sampler liang2025}, we get 
\[
\hat V_{\tau -1}  
 \approx \hat V_\tau + \frac{\beta_\tau}{2} \bc{\hat V_\tau  + s_\tau (\hat V_\tau)},\]
 which is essentially an Eulerian scheme with step-size $\beta_\tau/ 2$. This shows that \eqref{eq: DDIM sampler liang2025} roughly approximates a Eulerian scheme. Now, we can easily convert \eqref{eq: DDIM sampler liang2025} to a RF sampler using the equivalence between $V_\tau$ and $X_{t(\tau)}$. 
 
 To this end, we define the quantity $\sigma_t^2 = (1 -t)^2 + t^2$, and it follows that 
 \[
 \sigma_{t(\tau)}^2 = \frac{1}{(\sqrt{\omega(\tau)} + \sqrt{1 - \omega(\tau)})^2} = \frac{t^2(\tau)}{\omega_\tau}.
 \]
 Therefore, we have $V_\tau \overset{d}{ = } X_{t(\tau)}/ \sigma_{t(\tau)} $, and we have $s_{V_\tau}(y) = \sigma_{t(\tau)} s_{t(\tau)}(\sigma_{t(\tau)} y)$. We also define $t_i = t(N - i)$ for $i = 0 , 1, \ldots, N-1$, and $R_i = t_i/ \sigma_{t_i}$.
 Collecting everything above, an equivalent reformulation of \eqref{eq: DDIM sampler liang2025} for RF sampling is 
 \begin{equation}
     \label{eq: RF sampler non-Eulerian}
     \begin{aligned}
      &\frac{Y_{t_{i+1}}}{\sigma_{t_{i+1}}} = \frac{R_{i+1}}{R_i} \bc{ \frac{Y_{t_i}}{\sigma_{t_i}} + \eta_i \widehat s_{t_i}(Y_{t_i})}, \quad Y_{t_0} \sim N(0, \sigma_{t_0}^2 I_d),\\
      & i = 0, 1, \ldots, N-1,
     \end{aligned}
 \end{equation}
 where $\eta_i = \frac{1 - \frac{R_i^2}{R_{i+1}^2}}{1 + \sqrt{\frac{R_i^2}{R_{i+1}^2}.\frac{ 1 - R_{i+1}^2}{1  - R_i^2}}}$, and $\widehat s_{t_i}(x) = \frac{t_i \widehat v_{t_i}(x) - x}{1 - t_i}$.

 Finally, using the equivalence between \eqref{eq: DDIM sampler liang2025} and \eqref{eq: RF sampler non-Eulerian}, we arrive at the following theorem:
 \begin{theorem}
     \label{thm: RF sampler non-Eulerian}
     Under Assumption \ref{assumption: intrinsic dimension}-\ref{assumption: higher-order approximation error}, the RF sampler \eqref{eq: RF sampler non-Eulerian} yields, 
     \begin{align*}
     & \TV(Y_{t_{N-1}}, X_{t_{N-1}})\\
     & \lesssim \frac{k \log^3 N}{N} + \frac{\log^2 N}{N} \meanevel^2 + \frac{\log^2 N}{N } \meanevelJone^2\\
     & \quad + (\meanevel + \meanevelJone + \meanevelJtwo + \meanevelH)\log N.
     \end{align*}
 \end{theorem}
 Note that $t_{N-1} = \frac{\sqrt{\omega_1}}{\sqrt{\omega_1} + \sqrt{1 - \omega_1}} \approx 1$. Therefore, $\Law(X_{t_{N-1}}) \approx p_1$. In addition, the convergence rate is similar to the Eulerian scheme  as shown in Theorem \ref{thm: RF - convergence of Eulerian scheme}.
\end{comment}
\subsection{Convergence rate of STOC-RF}
In this section, we will leverage the equivalence between DDPM and STOC-RF to design an SDE-based sampler that also adapts to the low-dimensionality of $p_1$. We begin with the solution $\bc{Y^\prime}_{\tau\ge 0}$ to the forward process \eqref{eq: forward process}. %\textcolor{violet}{ALE: mention this is the discrete version of it}
% \[
% Y_\tau^\prime  = \sqrt{\omega_\tau} X_1 + B_{1- \omega_\tau} \overset{d}{=}\sqrt{\omega_\tau} X_1  + \sqrt{1 - \omega_\tau} \xi_\tau
% \]
% where $\bc{B_\tau}_{\tau\ge 0}$ is Brownian motion and $\xi_\tau \sim N(0, I_d)$.
We also consider the time-change map
$
\tau \in [0, \infty) \mapsto t(\tau) = \frac{\sqrt{\omega_\tau}}{\sqrt{\omega_\tau} + \sqrt{1 - \omega_\tau}}.
$
In light of \eqref{eq: DDPM and RF equivalent} and Proposition \ref{prop: RF-SDE and backwrd process are equivalent}, we note that $\bc{Y^\prime_\tau}_{\tau \ge 0}$ and $\{\tilde Z_t\}_{t \in [0,1]}$ (recall \ref{eq: RF-SDE}) are equivalent under the time change $\tau \mapsto t(\tau)$, and we have $\frac{Y^\prime_\tau}{\sqrt{\omega_\tau}} \overset{d}{=} \frac{\tilde Z_{t(\tau)}}{t(\tau)}$. %as a process (SR: not same as process. That would be same as saying forward and backward processes are same which is not true. However, Y_backward is same process as \tilde Z - Prop 3.4).
In practice, we simulate the forward process by discretizing $\tau$ to \textit{positive integers} and setting specific values for $\omega_\tau$. Here, we let $\omega_\tau = \prod_{j =1}^\tau \alpha_j$ with $\alpha_j$'s defined implicitly as 
\begin{equation}
    \label{eq: DDPM time scheduling}
    \begin{aligned}
        %& \beta_0 : = 1 - \alpha_0 = 1,\\
        & \beta_1 := 1 - \alpha_1 = 1/N^{c_0},\\
        & \beta_{\tau+1} := 1 - \alpha_{\tau+1}\\
        &  = \frac{c_1 \log N}{N} \min \bc{\beta_1 \br{1 + \frac{c_1 \log N}{N}}^\tau, 1}, \quad 1 \le \tau \le N,
    \end{aligned}
\end{equation}
for a user-specified choice of the parameters $c_0, c_1 >0$.
The time schedule above is very similar to the ones considered by \citet{potaptchik2024linear, li2024sharp, li2024adapting, huang2024denoising,liang2025low} to analyze the convergence properties of DDPM. 
% In light of Proposition \ref{prop: RF-SDE and backwrd process are equivalent}, we note that $\bc{Y^\prime_\tau}_{\tau \ge 0}$ and $\{\tilde Z_t\}_{t \ge 0}$ as defined in \eqref{eq: RF-SDE} are equivalent under the time change $t(\tau) = \frac{\sqrt{\omega_\tau}}{\sqrt{1 - \omega_\tau} + \sqrt{\omega_\tau}}$, so that $\frac{Y^\prime_\tau}{\sqrt{\omega_\tau}} \overset{d}{=} \frac{\tilde Z_{t(\tau)}}{t(\tau)}$. \textcolor{violet}{ALE: (i) ATTN $y'_{\tau}$  is a discrete process (ii) $Z_t$ has $t \in [0,1]$ and (iii) the notation $t(\cdot)$ }
Next, we recall the DDPM sampler considered in \citet{liang2025low, huang2024denoising, li2024adapting}:  
\begin{equation}
    \label{eq: DDPM sampler liang2025}
    \begin{aligned}
    & \hat Y^\prime_{\tau -1} = \frac{1}{\sqrt{\alpha_\tau}} \bc{ \hat Y^\prime_\tau + \delta_\tau \widehat s_{Y^\prime_\tau}(\hat Y^\prime_\tau) + \nu_\tau \xi_\tau}, \hat Y^\prime_N \sim N(0, I_d)\\
    & \delta_\tau = 1 - \alpha_\tau,  \nu_\tau = \sqrt{\frac{(\alpha_\tau - \omega_\tau)(1 - \alpha_\tau)}{1 - \omega_\tau}},\\
    & \tau   = N , N-1, \ldots, 1, \quad \{\xi_\tau\}_{\tau \ge 1} \overset{i.i.d.}{\sim} N(0,I_d),
    \end{aligned}
\end{equation}
where $ \widehat s_{Y^\prime_\tau}$ is an approximate score of $Y^\prime_\tau$.
To provide some intuition behind this sampler, we recall that $1 - \alpha_\tau = \beta_\tau \approx 0$. Hence, $1/\sqrt{\alpha_\tau}  \approx 1 + \beta_\tau/2, \delta_\tau/ \sqrt{\alpha_\tau} \approx \beta_\tau / 2$, and $\nu_\tau/ \sqrt{\alpha_\tau} \approx \sqrt{\beta_\tau }$. Substituting this in \eqref{eq: DDPM sampler liang2025}, we get 
\begin{equation}
\label{eq: vanilla SDE sampler}
\hat Y^\prime_{\tau -1}  
 \approx \hat Y^\prime_\tau + \frac{\beta_\tau}{2} \bc{\hat Y^\prime_\tau  + 2 \widehat s_{Y^\prime_\tau} (\hat Y^\prime_\tau)} + \sqrt{2 \br{\frac{\beta_{\tau}}{2}}} \xi_\tau,
 \end{equation}
 which is essentially the discretized version of \eqref{eq: reverse OU process} with step-size $\beta_\tau/ 2$. Therefore, the sampler \eqref{eq: DDPM sampler liang2025} roughly approximates the DDPM backward process. 
 
 Next, we convert \eqref{eq: DDPM sampler liang2025} to a STOC-RF sampler using the equivalence between $Y^\prime_\tau$ and $\tilde Z_{t(\tau)}$.  To this end, we let, for any $t \in [0,1]$, $\sigma_t^2 = (1 -t)^2 + t^2$, so that 
 \begin{equation*}
 %\label{eq: an identity}
 \sigma_{t(\tau)}^2 = \frac{1}{(\sqrt{\omega(\tau)} + \sqrt{1 - \omega(\tau)})^2} = \frac{t^2(\tau)}{\omega_\tau}.
 \end{equation*}
 Therefore, we have $Y^\prime _\tau \overset{d}{ = } \tilde Z_{t(\tau)}/ \sigma_{t(\tau)} $, and $s_{Y^\prime_\tau}(y) = \sigma_{t(\tau)} s_{t(\tau)}(\sigma_{t(\tau)} y)$. We also define $t_i := t(N - i)$ for $i = 0 , 1, \ldots, N-1$, and $R_i = t_i/ \sigma_{t_i}$.
 Collecting everything above, we arrive at an equivalent reformulation of \eqref{eq: DDPM sampler liang2025} for RF sampling of the form 
 \begin{equation}
     \label{eq: stoc-RF sampler non-Eulerian}
     \begin{aligned}
      &\frac{Y_{t_{i+1}}}{\sigma_{t_{i+1}}} = \frac{R_{i+1}}{R_i} \bc{ \frac{Y_{t_i}}{\sigma_{t_i}} + \eta_i \sigma_{t_i} \widehat s_{t_i}(Y_{t_i}) + \sqrt{\psi_i} W_{t_i}},\\
      &  \frac{Y_{t_0}}{\sigma_{t_0}} \sim N(0,  I_d), \quad \{W_{t_i}\}_{i \ge 0}\overset{i.i.d.}{\sim} N(0,I_d),\\
      & i = 0, 1, \ldots, N-1,
     \end{aligned}
 \end{equation}
 where $ \eta_i = 1- \frac{R_i^2}{R_{i+1}^2},   \psi_i = \frac{R_i^2}{R_{i+1}^2} \frac{ 1 - R_{i+1}^2}{1  - R_i^2} \br{1 - \frac{R_i^2}{R_{i+1}^2}}$, and $\widehat s_{t_i}(x) = \frac{t_i \widehat v_{t_i}(x) - x}{1 - t_i}$. See Appendix \ref{app: Stoc-RF convergence} for details. Lastly, we impose an assumption on the approximation error of $\widehat s_{t_i}$.
 \begin{assumption}
     \label{assumption: score approximation error stoc-RF}
     Define $(\escore_i)^2:= \bbE \Norm{s_{t_i}(X_{t_i}) - \widehat{s}_{t_i}(X_{t_i})}_2^2$. Then, 
    $\frac{1}{N} \sum_{i =1}^N (\escore_i)^2 \le \meanescore^2.$
 \end{assumption}

 Leveraging the equivalence between \eqref{eq: DDPM sampler liang2025} and \eqref{eq: stoc-RF sampler non-Eulerian}, we obtain convergence guarantee for STOC-RF, which the second main result of this paper.
 \begin{theorem}
     \label{thm: stoc-RF sampler non-Eulerian}
     Under Assumption \ref{assumption: intrinsic dimension}, \ref{assumption: bounded support} and \ref{assumption: score approximation error stoc-RF}, the stoc-RF sampler \eqref{eq: stoc-RF sampler non-Eulerian} satisfies 
     \begin{align*}
     & \TV(Y_{t_{N-1}}, X_{t_{N-1}}) \lesssim \frac{k \log^3 N}{N} +  \meanescore \sqrt{\log N}.
     \end{align*}
 \end{theorem}
 Note that $t_{N-1} = \frac{\sqrt{\omega_1}}{\sqrt{\omega_1} + \sqrt{1 - \omega_1}} \approx 1$. Therefore, $\Law(X_{t_{N-1}}) \approx p_1$. In comparison with Theorem \ref{thm: RF - convergence of Eulerian scheme}, STOC-RF sampler does not require any conditions on higher-order derivatives (\Cref{assumption: higher-order approximation error}) of the score or velocity field. That is, STOC-RF sampler can achieve similar performance under less stringent requirements. This is primarily due to the Gaussian perturbation term in \eqref{eq: stoc-RF sampler non-Eulerian} which acts as a corrector step. The proof is in Appendix \ref{app: Stoc-RF convergence}.
%\vspace{0.5em}

\begin{remark}
\label{remark: new RF sampler}
Using the equivalence relation \eqref{eq: DDPM and RF equivalent}, it is also possible to directly relate the deterministic RF sampler to the DDIM sampler studied by \citet[Theorem 1]{liang2025low} (Appendix \ref{app: new RF sampler}) via the same time change transformation considered in this section. The resulting time schedule also follows a U-shaped distribution and the sampler will also adapts to the low-dimensionality of $p_1$.
%See Appendix \ref{app: new RF sampler} for details. 
%more unwieldly  -- sOne can show that the new sampler also adapts to the low-dimensionality of $p_1$. However, we omit this discussion for brevity as the RF sampler in Theorem \ref{thm: RF - convergence of Eulerian scheme} already adapts to low-dimension with simpler time-steps.    
\end{remark}






\section{Experiments}
We perform extensive simulations to corroborate our theoretical findings for both synthetic data and T2I experiments. Codes are available in \href{https://github.com/roysaptaumich/Low-dim-RF}{https://github.com/roysaptaumich/Low-dim-RF}. 
%More T2I experimental results can be found in \Cref{app: T2I exp} and \Cref{app: RF and Stoc-RF T2I}.
\begin{figure*}[t!]
  \centering
  \begin{subfigure}[b]{0.24\textwidth}
    \centering
    \includegraphics[width=\textwidth]{figures/N_100_low_rank_TV_vs_dimension.pdf}
    \caption{}
    \label{fig:sub1}
  \end{subfigure}\hfill
  \begin{subfigure}[b]{0.24\textwidth}
    \centering
    \includegraphics[width=\textwidth]{figures/N_200_low_rank_TV_vs_dimension.pdf}
    \caption{}
    \label{fig:sub2}
  \end{subfigure}\hfill
  \begin{subfigure}[b]{0.24\textwidth}
    \centering
    \includegraphics[width=\textwidth]{figures/N_100_low_rank_TV_geom_vs_sde_vs_dimension.pdf}
    \caption{}
    \label{fig:sub3}
    \end{subfigure}\hfill
  \begin{subfigure}[b]{0.24\textwidth}
    \centering
    \includegraphics[width=\textwidth]{figures/N_200_low_rank_TV_geom_vs_sde_vs_dimension.pdf}
    \caption{}
    \label{fig:sub4}
    \end{subfigure}
  \caption{(a)-(b) \textit{Estimated} \TV~ distance between \textit{blurred} low-rank Gaussian distribution and the generated samples by the RF sampler \eqref{eq: emp-ode-disc} under uniform and non-uniform time-grid under varying $d$. (c)-(d) \textit{Estimated} \TV~ distance between \textit{blurred} low-rank Gaussian distribution and the generated samples by the RF and STOC-RF samplers under uniform and non-uniform time-grid under varying $d$.}
  \label{fig: TV distance comparison low-rank Gaussian}
\end{figure*}


\begin{figure*}[h]
    \centering
    
    % --- ROW 1: Two Subfigures Separated by \hfill ---
    \begin{subfigure}[t]{0.3\textwidth}
        \centering
        \includegraphics[width=\textwidth]{figures/flux_figures/grape_vine_N_50.pdf}
        \caption{\scriptsize\sf \textbf{Prompt:} Grape vines in the shape of text \orange{`open your mind’} sprouting out of a head with flowers and butterflies. DSLR photo. (\underline{Better text rendering}, $N = 50$)}
        %\label{fig:sub1}
    \end{subfigure}\hfill
    % Note: No \hfill after the second subfigure in the row
    \begin{subfigure}[t]{0.3\textwidth}
        \centering
        % Assuming the fourth figure is unique, even if the filename is duplicated
        \includegraphics[width=\textwidth]{figures/flux_figures/chimp_arguing_N_50.pdf}
        \caption{\scriptsize \sf \textbf{Prompt:} A sign that says \orange{"Please refrain from arguing with the
chimpanzees"}. (\underline{Minimal mistakes}, $N = 50$)}
        %\label{fig:sub4}
    \end{subfigure}\hfill
    \begin{subfigure}[t]{0.3\textwidth}
        \centering
        % Assuming the fourth figure is unique, even if the filename is duplicated
        \includegraphics[width=\textwidth]{figures/flux_figures/robo_teacher_N_50.pdf}
        \caption{\scriptsize \sf \textbf{Prompt:} Photo of a robot lecturer writing the words \orange{`Representation Learning'} in cursive on a blackboard, with math formulas and diagrams. (\underline{Correct spelling}, $N = 50$)}
        %\label{fig:sub6}
    \end{subfigure}
    
    
    \caption{Generated images using Flux for specific prompts via Euler's scheme \eqref{eq: emp-ode-disc} under \red{non-uniform} time-grid \eqref{eq: gemoetric time discretization} (\red{Red frame}) and \blue{uniform} time-grid (\blue{Blue frame}). Generation quality of RF is better under \red{non-uniform} time-grid \eqref{eq: gemoetric time discretization}: less hallucinations, enjoys better text rendering and achieves better semantic completeness.}
    \label{fig: Flux experiments mini}
\end{figure*}

\textbf{Synthetic data.}
In this section we demonstrate the usefulness of non-uniform time-grid for generation task from low-dimensional distribution.
We set the intrinsic dimension $k =8$, and set our target distribution to be $p_1 = N(\mu, \Sigma)$, where $\mu = (8, 8, \ldots, 8)^\top \in \bbR^d$  and $\Sigma$ is a block-diagonal matrix with blocks $I_{k}$ and $ \boldsymbol{0}_{(d-k) \times (d-k)}$. 
Tweedie's formula (Section \ref{sec: Rf}) yields
\begin{equation}
\label{eq: low-dim velocity field}
\begin{aligned}
 v_{t_i}(x) 
 = \begin{bmatrix}
    \frac{2t_i-1}{\sigma_{t_i}^2} I_k & \boldsymbol{0}_{k \times (d-k)}\\
    \boldsymbol{0}_{(d-k) \times k} & - \frac{1}{1-t_i}I_{d-k}
\end{bmatrix} x
+ \begin{bmatrix}
    \frac{1-t_i}{\sigma_{t_i}^2} I_k  & \boldsymbol{0}_{k \times (d-k)}\\
    \boldsymbol{0}_{(d-k) \times k} & \frac{1}{1 - t_i} I_{d-k}
\end{bmatrix} \mu.
\end{aligned}
\end{equation}

% \rd Then, using Tweedie's formula (see Section~\ref{sec: Rf}) the drift $v_{t_i}(x)$ is given by $Bx$ (see Appendix Section~\ref{}). The $d\times d$ matrix $B$ is block-diagonal with blocksizes $k$ and $d-k$. Each block is an appropriate scaling of the identity matrix. 
% \bk

We use \eqref{eq: low-dim velocity field} for generating new samples from $p_1$. We vary $d$ between 10 and 800, and set $N \in \bc{100, 200}$. We use both a non-uniform time-grid \eqref{eq: gemoetric time discretization} (with $\delta = \min\{1/N, 1/d\}$) and a uniform time-grid for the generation task. We generate 2000 samples using \eqref{eq: gemoetric time discretization} and measure the \TV~distance with respect to another batch of 2000 samples independently generated from the slightly blurred target distribution $\Law((1 - \delta) X_1 + \delta Z)$, where $Z \sim N(0, I_d)$. We use the discriminative approach proposed in \citet{tao2024discriminative} to estimate the \TV~distance between the distributions via 10 independent Monte Carlo rounds. Figure \ref{fig: TV distance comparison low-rank Gaussian}(a)-(b) show that nonuniform discretization has better adaptation to the low-dimensional structure of $p_1$ for large $d$.
% \begin{figure}[h]
%   \centering
%   \begin{subfigure}[b]{0.24\textwidth}
%     \centering
%     \includegraphics[width=\textwidth]{ICML2026/figures/N_100_low_rank_TV_vs_dimension.pdf}
%     \caption{}
%     \label{fig:sub1}
%   \end{subfigure}\hfill
%   \begin{subfigure}[b]{0.24\textwidth}
%     \centering
%     \includegraphics[width=\textwidth]{ICML2026/figures/N_200_low_rank_TV_vs_dimension.pdf}
%     %\caption{Caption for second subfigure}
%     \label{fig:sub2}
%   \end{subfigure}
%   \caption{\TV~ distance between \textit{blurred} low-rank Gaussian distribution and the generated samples by the RF sampler under uniform and non-uniform time-grid with varying $d$.}
%   \label{fig: TV distance comparison low-rank Gaussian}
% \end{figure}
For the SDE sampler \eqref{eq: stoc-RF sampler non-Eulerian}, we also compare its \TV~convergence with that of the deterministic sampler \eqref{eq: emp-ode-disc} with the non-uniform time grid \eqref{eq: gemoetric time discretization}. Figure \ref{fig: TV distance comparison low-rank Gaussian}(c)-(d) show that both samplers have similar performance across varying dimensions.

\textbf{T2I experiments.}~We perform experiments using Flux \citep{flux2024, labs2025flux1kontextflowmatching} to demonstrate the usefulness of the U-shaped time schedule \eqref{eq: gemoetric time discretization}. 
% For our experiments, we use the prompts designed by
% \citet{liu2023character, hu2025amo} to test and assess the generation quality. 
Figure \ref{fig: Flux experiments mini} shows that generation quality of RF is superior under a non-uniform time-grid \eqref{eq: gemoetric time discretization} compared to uniform time-grid. Figure \ref{fig: Flux SDE/RF experiments mini} compares the generation quality of STOC-RF sampler \eqref{eq: stoc-RF sampler non-Eulerian} and the vanilla SDE sampler \eqref{eq: vanilla SDE sampler} and ODE samplers. We see that the generation quality is generally better for STOC-RF. Details are deferred to \Cref{sec:appendix:exp}.

% \begin{figure}[h]
%   \centering
%   \begin{subfigure}[b]{0.24\textwidth}
%     \centering
%     \includegraphics[width=\textwidth]{ICML2026/figures/N_100_low_rank_TV(geom vs sde)_vs_dimension.pdf}
%     %\caption{Caption for first subfigure}
%     \label{fig:sub1}
%   \end{subfigure}\hfill
%   \begin{subfigure}[b]{0.24\textwidth}
%     \centering
%     \includegraphics[width=\textwidth]{ICML2026/figures/N_200_low_rank_TV(geom vs sde)_vs_dimension.pdf}
%     %\caption{Caption for second subfigure}
%     \label{fig:sub2}
%   \end{subfigure}
%   \caption{\TV~ distance between \textit{blurred} low-rank Gaussian distribution and the generated samples by the RF sampler and the stoc-RF sampler with varying $d$.}
%   \label{fig: TV distance comparison low-rank Gaussian ode vs sde}
% \end{figure}

\begin{figure}[t!]
    \centering

    % --- ROW 1: Two Subfigures Separated by \hfill ---
    \begin{subfigure}[t]{0.7\textwidth}
        \centering
        \includegraphics[width=\textwidth]{figures/flux_figures/cat_banner_N_50_all_samplers-flat.pdf}
        \caption{\sf \textbf{Prompt:} A photo of a cat holding a sign \orange{Rectified flow}. ($N = 50$)}
        %\label{fig:sub1}
    \end{subfigure}\hfill
    % Note: No \hfill after the second subfigure in the row
    %\vspace{1em}
    \begin{subfigure}[t]{0.7\textwidth}
        \centering
        \includegraphics[width=\textwidth]{figures/flux_figures/grape_vine_N_50_all_samplers-flat.pdf}
        \caption{\sf \textbf{Prompt:} Grape vines in the shape of text \orange{`open your mind’} sprouting out of a head with flowers and butterflies. DSLR photo. ($N = 50$)}
       %\label{fig:sub2}
    \end{subfigure}
    \vspace{1em}
    \begin{subfigure}[t]{0.7\textwidth}
        \centering
        % Assuming the fourth figure is unique, even if the filename is duplicated
        \includegraphics[width=\textwidth]{figures/flux_figures/robo_teacher_N_50_all_samplers-flat.pdf}
        \caption{\sf \textbf{Prompt:} Photo of a robot lecturer writing the words \orange{`Representation Learning'} in cursive on a blackboard, with math formulas and diagrams. ($N = 50$)}
        %\label{fig:sub5}
    \end{subfigure}
    \vspace{1em}
    \begin{subfigure}[t]{0.7\textwidth}
        \centering
        % Assuming the fourth figure is unique, even if the filename is duplicated
        \includegraphics[width=\textwidth]{figures/flux_figures/cheese_car_N_80_all_samplers-flat.pdf}
        \caption{\sf \textbf{Prompt:} A \orange{car made of cheese} on a wooden table. ($N = 80$)}
        %\label{fig:sub6}
    \end{subfigure}
   
    
    
    \caption{Generated images using Flux for specific prompts via Euler's scheme \eqref{eq: emp-ode-disc} under non-uniform time-grid \eqref{eq: gemoetric time discretization} (\red{Columns 1}),  uniform time-grid (Column 2), STOC-RF sampler (\red{Column 3}), and vanilla SDE sampler (Column 4). 
    % \red{Red} color indicates the proposed methods in this paper. The generation quality of Stoc-RF is generally better out of the four samplers. More figures can be found in \Cref{app: RF and Stoc-RF T2I}.
    }
    \label{fig: Flux SDE/RF experiments mini}
\end{figure}

\textbf{Caution:}~The \TV~convergence rates in both \Cref{thm: RF - convergence of Eulerian scheme} and \Cref{thm: stoc-RF sampler non-Eulerian} depend highly on the approximation errors. 
% In our experiments, we find that achieving good approximation error in high-dimensional regime is difficult even for the aforementioned simple low-dimensional Gaussian distribution case. We demonstrate this in a low-dimensional Gaussian example in Figure \ref{fig: Generation quality with varying d} where we set $k = 2$. The figure shows the generated data along the 3rd dimension for both the RF and STOC-RF samplers. We can see that the growing dimension has an adverse effect on the generation quality. In fact, STOC-RF appears to be more brittle, presumably due to its inherent stochasticity. 
Difficulty may arise in estimating $v_{t}(\cdot)$ along the direction of low-dimensionality when $t \approx 1$ due to non-smoothness (see \eqref{eq: low-dim velocity field}). Such non-smoothness is usually hard to capture via standard neural networks (e.g., MLP, U-Net) that are typically smooth, and necessitate the adoption of more expressive architecture. We defer this to future research.



\section{Conclusions}
In this paper, we design a new time-discretization scheme that allows any off-the-shelf pretrained RF sampler to achieve accelerated convergence by adapting to the low-dimensional structure of the target distribution. This allows the RF model to achieve better sampling quality in the first reflow step itself and also corroborates the empirical findings in \citet{lee2024improvingtrainingrectifiedflows}. Next, we show that the DDPM is equivalent to a stochastic version of RF by establishing a novel connection between these processes and the stochastic localization. Building on this connection, we further design a stochastic RF sampler that also adapts to
the low-dimensionality of the target distribution under less stringent assumptions on the learned RF model. 





\section*{Acknowledgements}
We thank Prof. Eric Vanden-Eijnden\textsuperscript{\orcidlink{0000-0001-7982-5854}} at NYU, Prof. Sanjay Shakkottai\textsuperscript{\orcidlink{0000-0002-4325-9050}} at UT Austin, and Dr. Christian Rau\textsuperscript{\orcidlink{0000-0003-0852-5015}} for helpful discussions and pointing out important references. AR, PS and SR gratefully acknowledge the NSF grants CCF-2019844 and CCF-2505865. PS and SR gratefully acknowledges NSF grant 2217069.
% PS also acknowledges
%DMS 2109155.

%\nocite{langley00}

\bibliography{icml}
\bibliographystyle{icml2026}

%%%%%%%%%%%%%%%%%%%%%%%%%%%%%%%%%%%%%%%%%%%%%%%%%%%%%%%%%%%%%%%%%%%%%%%%%%%%%%%
%%%%%%%%%%%%%%%%%%%%%%%%%%%%%%%%%%%%%%%%%%%%%%%%%%%%%%%%%%%%%%%%%%%%%%%%%%%%%%%
% APPENDIX
%%%%%%%%%%%%%%%%%%%%%%%%%%%%%%%%%%%%%%%%%%%%%%%%%%%%%%%%%%%%%%%%%%%%%%%%%%%%%%%
%%%%%%%%%%%%%%%%%%%%%%%%%%%%%%%%%%%%%%%%%%%%%%%%%%%%%%%%%%%%%%%%%%%%%%%%%%%%%%%
\newpage
\appendix
\onecolumn
The appendix is organized as follows.
\begin{enumerate}
    \item Section~\ref{sec:appendix:exp} contains additional experimental results.
    \item Section~\ref{app: proof of main results} contains detailed proofs from Sections~\ref{sec: main results} and the description of the RF sampler.
    \item Section~\ref{app: SL lemmas} contains detailed proofs of results about Section~\ref{sec: SL}.
    \item Finally, Section~\ref{app: auxiliary results} contains all other helpful lemmas and their proofs.
\end{enumerate}

\section{Additional Experiments}\label{sec:appendix:exp}
% \subsection{Comparison of RF sampler under uniform and non-uniform grid in T2I setting}
% \label{app: T2I exp}
We have performed various experiments to demonstrate the usefulness of the U-shaped time schedule \eqref{eq: gemoetric time discretization} in the context of T2I. We use the Flux model \citep{flux2024, labs2025flux1kontextflowmatching} which is a pretrained RF model. For our experiments, we use the prompts designed by
\citet{liu2023character, hu2025amo} to test and assess the generation quality. In these examples, we always set $\delta = 1/N$  for a specified choices of $N$. In Figure \ref{fig: Flux experiments} we see that generation quality of RF is superior under a non-uniform time-grid \eqref{eq: gemoetric time discretization} compared to uniform time-grid. In particular, the experiments indicate that a non-uniform grid helps to reduce hallucinations, enjoys better text rendering and achieves better semantic completeness.

We also compare the generation quality of STOC-RF sampler \eqref{eq: stoc-RF sampler non-Eulerian} and the vanilla SDE sampler \eqref{eq: vanilla SDE sampler} in Figure \ref{fig: Flux SDE/RF experiments}. The generated plots show that the quality is generally better for STOC-RF sampler \eqref{eq: stoc-RF sampler non-Eulerian} compared to the RF sampler and the vanilla SDE sampler \eqref{eq: vanilla SDE sampler}. 



\begin{figure}[h!]
    \centering
    
    % --- ROW 1: Two Subfigures Separated by \hfill ---
    \begin{subfigure}[t]{0.3\textwidth}
        \centering
        \includegraphics[width=\textwidth]{figures/flux_figures/cat_banner_N_50.pdf}
        \caption{\scriptsize\sf \textbf{Prompt:} A photo of a cat holding a sign \orange{Rectified flow}. (\underline{Less hallucination}, $N = 50$)}
        \label{fig:sub1}
    \end{subfigure}\hfill
    % Note: No \hfill after the second subfigure in the row
    \begin{subfigure}[t]{0.3\textwidth}
        \centering
        \includegraphics[width=\textwidth]{figures/flux_figures/grape_vine_N_50.pdf}
        \caption{\scriptsize\sf \textbf{Prompt:} Grape vines in the shape of text \orange{`open your mind’} sprouting out of a head with flowers and butterflies. DSLR photo. (\underline{Better text rendering}, $N = 50$)}
        \label{fig:sub2}
    \end{subfigure}\hfill
    \begin{subfigure}[t]{0.3\textwidth}
        \centering
        % Assuming the third figure is unique, even if the filename is duplicated
        \includegraphics[width=\textwidth]{figures/flux_figures/hubble_space_N_50.pdf}
        \caption{\scriptsize \sf \textbf{Prompt:} the \orange{hubble telescope} and the milky way, with the text \orange{``the universe is a mystery, but we are here to solve it"} (\underline{Better fidelity/ prompt adherence}, $N = 50$)}
        \label{fig:sub3}
    \end{subfigure}
    % �� KEY FIX: Force a new line with a mandatory break (\\) 
    % and add vertical space (\vspace) for separation.
    \vspace{1em} 
    % --- ROW 2: Two Subfigures Separated by \hfill ---
    \begin{subfigure}[t]{0.3\textwidth}
        \centering
        % Assuming the fourth figure is unique, even if the filename is duplicated
        \includegraphics[width=\textwidth]{figures/flux_figures/chimp_arguing_N_50.pdf}
        \caption{\scriptsize \sf \textbf{Prompt:} A sign that says \orange{"Please refrain from arguing with the
chimpanzees"}. (\underline{Minimal mistakes}, $N = 50$)}
        \label{fig:sub4}
    \end{subfigure}\hfill
    \begin{subfigure}[t]{0.3\textwidth}
        \centering
        % Assuming the fourth figure is unique, even if the filename is duplicated
        \includegraphics[width=\textwidth]{figures/flux_figures/robo_teacher_N_50.pdf}
        \caption{\scriptsize \sf \textbf{Prompt:} Photo of a robot lecturer writing the words \orange{`Representation Learning'} in cursive on a blackboard, with math formulas and diagrams. (\underline{Correct spelling}, $N = 50$)}
        \label{fig:sub5}
    \end{subfigure}\hfill 
    \begin{subfigure}[t]{0.3\textwidth}
        \centering
        % Assuming the fourth figure is unique, even if the filename is duplicated
        \includegraphics[width=\textwidth]{figures/flux_figures/cheese_car_N_80.pdf}
        \caption{\scriptsize \sf \textbf{Prompt:} A \orange{car made of cheese} on a wooden table. (\underline{Semantic completeness}, $N = 80$)}
        \label{fig:sub6}
    \end{subfigure}
    
    
    \caption{Generated images using Flux for specific prompts via Euler's scheme \eqref{eq: emp-ode-disc} under non-uniform time-grid \eqref{eq: gemoetric time discretization} (\red{Red frame}) and uniform time-grid (\blue{Blue frame}).}
    \label{fig: Flux experiments}
\end{figure}

\newpage
% \subsection{Comparison of generation quality of samplers with growing dimension}

% \begin{figure}[h!]
%   \centering
%   \begin{subfigure}[b]{0.3\textwidth}
%     \centering
%     \includegraphics[width=\textwidth]{ICML2026/figures/RF_sampler_+_Geometric_grid_(d_=_3).pdf}
%     %\caption{Caption for first subfigure}
%     \label{fig:sub1}
%   \end{subfigure}\hfill
%   \begin{subfigure}[b]{0.3\textwidth}
%     \centering
%     \includegraphics[width=\textwidth]{ICML2026/figures/RF_sampler_+_Geometric_grid_(d_=_10).pdf}
%     %\caption{Caption for second subfigure}
%     \label{fig:sub2}
%   \end{subfigure}\hfill
%   \begin{subfigure}[b]{0.3\textwidth}
%     \centering
%     \includegraphics[width=\textwidth]{ICML2026/figures/RF_sampler_+_Geometric_grid_(d_=_100).pdf}
%     %\caption{Caption for second subfigure}
%     \label{fig:sub3}
%   \end{subfigure}
%   \vspace{1em}
%   \centering
%   \begin{subfigure}[b]{0.3\textwidth}
%     \centering
%     \includegraphics[width=\textwidth]{ICML2026/figures/Stoc-RF_sampler_(d_=_3).pdf}
%     %\caption{Caption for first subfigure}
%     \label{fig:sub1}
%   \end{subfigure}\hfill
%   \begin{subfigure}[b]{0.3\textwidth}
%     \centering
%     \includegraphics[width=\textwidth]{ICML2026/figures/Stoc-RF_sampler_(d_=_10).pdf}
%     %\caption{Caption for second subfigure}
%     \label{fig:sub2}
%   \end{subfigure}\hfill
%   \begin{subfigure}[b]{0.3\textwidth}
%     \centering
%     \includegraphics[width=\textwidth]{ICML2026/figures/Stoc-RF_sampler_(d_=_100).pdf}
%     %\caption{Caption for second subfigure}
%     \label{fig:sub3}
%   \end{subfigure}
%   \caption{ The generation quality of RF sampler and Stoc-RF sampler along the low-dimensional direction gradually degrades with growing dimension}
%   \label{fig: Generation quality with varying d}
% \end{figure}





% \newpage

% \subsection{Comparison of RF sampler and STOC-RF sampler in T2I setting}
% \label{app: RF and Stoc-RF T2I}
% \begin{figure}[h!]
%     \centering
    
%     % --- ROW 1: Two Subfigures Separated by \hfill ---
%     \begin{subfigure}[t]{0.7\textwidth}
%         \centering
%         \includegraphics[width=\textwidth]{figures/flux_figures/cat_banner_N_50_all_samplers-flat.pdf}
%         \caption{\sf \textbf{Prompt:} A photo of a cat holding a sign \orange{Rectified flow}. ($N = 50$)}
%         %\label{fig:sub1}
%     \end{subfigure}\hfill
%     % Note: No \hfill after the second subfigure in the row
%     %\vspace{1em}
%     \begin{subfigure}[t]{0.7\textwidth}
%         \centering
%         \includegraphics[width=\textwidth]{figures/flux_figures/grape_vine_N_50_all_samplers-flat.pdf}
%         \caption{\sf \textbf{Prompt:} Grape vines in the shape of text \orange{`open your mind’} sprouting out of a head with flowers and butterflies. DSLR photo. ($N = 50$)}
%        %\label{fig:sub2}
%     \end{subfigure}
%     \vspace{1em}
%     \begin{subfigure}[t]{0.7\textwidth}
%         \centering
%         % Assuming the fourth figure is unique, even if the filename is duplicated
%         \includegraphics[width=\textwidth]{figures/flux_figures/robo_teacher_N_50_all_samplers-flat.pdf}
%         \caption{\sf \textbf{Prompt:} Photo of a robot lecturer writing the words \orange{`Representation Learning'} in cursive on a blackboard, with math formulas and diagrams. ($N = 50$)}
%         %\label{fig:sub5}
%     \end{subfigure}
%     \vspace{1em}
%     \begin{subfigure}[t]{0.7\textwidth}
%         \centering
%         % Assuming the fourth figure is unique, even if the filename is duplicated
%         \includegraphics[width=\textwidth]{figures/flux_figures/cheese_car_N_80_all_samplers-flat.pdf}
%         \caption{\sf \textbf{Prompt:} A \orange{car made of cheese} on a wooden table. ($N = 80$)}
%         %\label{fig:sub6}
%     \end{subfigure}
    
    
%     \caption{Generated images using Flux for specific prompts via Euler's scheme \eqref{eq: emp-ode-disc} under non-uniform time-grid \eqref{eq: gemoetric time discretization} (\red{Columns 1}),  uniform time-grid (Column 2), STOC-RF sampler (\red{Column 3}), and vanilla SDE sampler (Column 4). \blue{Red} color indicates the proposed methods in this paper. The generation quality of Stoc-Rf is generally better out of the four samplers.}
%     \label{fig: Flux SDE/RF experiments}
% \end{figure}

\newpage
\section{Proof of main results}
\label{app: proof of main results}



%%%%%%%%%%%%%%%%%%%%%%%%%%
%%% Proof of RF (ODE) %%%%
%%%%%%%%%%%%%%%%%%%%%%%%%%
\subsection{Velocity function and its first order identities}
\label{app: v first order identities}
Before proceeding, we illustrate the connection between the velocity $v_t(x)$ and the score function $s_t(x):= \nabla \log p_t(x)$. 
% , where $p_t$ is the density of $X_t$ w.r.t. the Lebesgue measure for $t\in [0,1)$
 To this end, we define two key quantities:
\[
\mu_{1\mid t}(x) = \bbE[X_1 \mid X_t = x], 
\; 
\Cov_{1 \mid t}(x) = \Cov(X_1 \mid X_t = x).
\]
With this, true drift be re-written as
\begin{equation}
\label{eq: v_tweedie}
v_t(x) = \bbE[X_1 - X_0 \mid X_t = x] = \frac{1}{1-t}\br{\mu_{1 \mid t}(x) - x}.
\end{equation}

By the first and second order Tweedie's formula \citep{efron2011tweedie, meng2021estimating}, it follows that, for any choice of $t \in (0,1]$ and $x \in \mathbb{R}^d$,
\[
\mu_{1 \mid t}(x) =   \frac{x}{t} + \frac{(1-t)^2}{t} s_t(x),\] 
\[
\Cov_{1 \mid t}(x) = \frac{(1-t)^2}{t^2}I_d + \frac{(1-t)^4}{t^2} \nabla s_t(x).
\]
Substituting this in \eqref{eq: v_tweedie}, we obtain that 
$$
v_t(x) = \frac{x}{t} + \frac{(1-t)}{t} s_t(x), \quad  \text{and} \quad \nabla v_t(x) = - \frac{I_d}{1-t} + \frac{t}{(1-t)^3} \Cov_{1 \mid t}(x).
$$

\subsection{Proof of convergence of RF sampler (Theorem \ref{thm: RF - convergence of Eulerian scheme})}
\label{app: RF convergence}
\paragraph{Notational conventions.}
We begin with some notational conventions for the ease of presentation. 

\begin{enumerate}
    \item We define transition map for the Euler's scheme as 
\[
\Phi_{t_i}(x):= x + \eta_{i} \widehat{v}_{t_i}(x) \quad \text{with $\eta_{i} = t_{i+1} - t_i$},
\]
where $\widehat v_{t}(x)$ is the velocity estimate. 

\item For notational brevity, we will use index $i$ to denote the time $t_i$. In this index notation we have $X_N \equiv X_1 \sim p_1$ 

\item With the above notations in mind, we have 
\begin{equation}
\label{eq: interp in modified index}
X_i = (1-t_i)X_0 + t_iX_N, \quad \text{and}\quad Y_{i+1}=\Phi_{i}\left(Y_{i}\right).
\end{equation}
\end{enumerate}


\paragraph{Step 1: TV recursion.}
We start by obtaining a recursion relation between the \TV distances between the updates $Y_i$ and the ground truth $X_i$. For any $i \le N-2$, we have
\begin{equation}
\label{eq: TV recursion}
\begin{aligned}
\TV\left(p_{X_{i+1}}, p_{Y_{i+1}}\right) & =\sup _{\mathcal{A} \subseteq \bbR^d}\left\{\mathbb{P}_{X_{i+1}}(\mathcal{A})-\mathbb{P}_{Y_{i+1}}(\mathcal{A})\right\}=\sup _{\mathcal{A} \subseteq \bbR^d}\left\{\mathbb{P}_{X_{i+1}}(\mathcal{A})-\mathbb{P}_{Y_{i}}\left(\Phi_{i}^{-1}(\mathcal{A})\right)\right\} \\
& \leq \sup _{\mathcal{A} \subseteq \bbR^d}\left\{\mathbb{P}_{X_{i+1}}(\mathcal{A})-\mathbb{P}_{X_{i}}\left(\Phi_{i}^{-1}(\mathcal{A})\right)\right\}+\sup _{\mathcal{A} \subseteq \bbR^d}\left\{\mathbb{P}_{X_{i}}\left(\Phi_{i}^{-1}(\mathcal{A})\right)-\mathbb{P}_{Y_{i}}\left(\Phi_{i}^{-1}(\mathcal{A})\right)\right\}  \\
& \leq \sup _{\mathcal{A} \subseteq \bbR^d}\left\{\mathbb{P}_{X_{i+1}}(\mathcal{A})-\mathbb{P}_{\Phi_{i}\left(X_{i}\right)}(\mathcal{A})\right\}+\TV\left(p_{X_{i}}, p_{Y_{i}}\right) \\
& =\TV\left(p_{X_{i+1}}, p_{\Phi_{i}\left(X_{i}\right)}\right)+\TV\left(p_{X_{i}}, p_{Y_{i}}\right)
\end{aligned}
\end{equation}

where the first identity arises from the basic property of the TV distance.

\paragraph{Step 2: Simplifying $\TV(p_{X_{i+1}}, p_{\Phi_i(X_i)})$.}
% Next, we define a set $\mathcal{E}_{t}$ as

% \begin{equation}
% \label{eq: def of E_i}
% \mathcal{E}_{i}:=\left\{x \in \mathcal{X}_{\text {data }}: \eta_{i}\left\|\varepsilon_{i}^{J}(x)\right\| \leq \frac{1}{8}\right\}
% \end{equation}


% Based on the set $\cE_i$, we define an auxiliary map $\widetilde{\Phi}_{i}: \mathcal{X}_{\text {data }} \rightarrow \mathcal{R}^{d} \cup\{\infty\}$ as
% $$
% \widetilde{\Phi}_{i}(x)= \begin{cases}\Phi_{i}(x), & \text { if } x \in \mathcal{E}_{i}  \\ \infty, & \text { otherwise }\end{cases}
% $$
Next, we recall that $v_i(x) = \frac{x}{t_i} + \frac{1-t_i}{t_i} s_i(x)$. Therefore, by \citet[Lemma 13]{Liu2024_FlowNotes} we have 
\begin{equation}
\label{eq: partial v_x}
\frac{\partial v_i(x)}{\partial x} = -  \frac{1}{1-t_i} I + \frac{ t_i}{(1-t_i)^3} \Cov_{1 \mid t_i}(x_i)
\end{equation}
Further, for any vector $u \in \mathbb{R}^{d}$ with $\|u\|_{2}=1$ 
%and any $x \in \mathcal{E}_{i}$,
Assumption \ref{assumption: higher-order approximation error}\ref{assumption: Jone} tells us that for $i \le N-2$
$$
\begin{aligned}
% v^{\top} \frac{\partial \widetilde{\Phi}_{i}(x)}{\partial x} v 
& u^{\top} \frac{\partial \Phi_{i}(x)}{\partial x} u=u^{\top}\left(I+\eta_{i} \frac{\partial \widehat v_{i}(x)}{\partial x}\right) u=u^{\top}\left(I+\eta_{i} \frac{\partial v_{i}(x)}{\partial x}\right) u+\eta_{i} u^{\top} \varepsilon_{i}^{J}(x) u \\
& \geq u^{\top}\left(I+\eta_{i} \frac{\partial v_{i}(x)}{\partial x}\right) u-\frac{1}{8}\|u\|_{2}^{2}=u^{\top}\left\{\left(1-\frac{\eta_{i} }{1-t_i}\right) I+\frac{ \eta_{i} t_i}{(1-t_i)^3} \operatorname{Cov}_{1 \mid t_i}(x)\right\} u-\frac{1}{8}\|u\|_{2}^{2} \\
& \geq\left(1-h-\frac{1}{8}\right)\|u\|_{2}^{2}>\frac{3}{8} , \quad \text{if $h <  1/2$.}
\end{aligned}
$$

Here, the penultimate relation follows from the positive semi-definiteness of the covariance matrix $\operatorname{Cov}_{1 \mid t}(x)$ and the fact that $\frac{\eta_i}{1 - t_i} = \frac{h t_i}{1 - t_i} \ind\bc{t_i \le 1/2} + \frac{h(1-t_{i+1})}{1 - t_i} \ind\bc{t_i>1/2} \le h$ (see \Cref{lemma: properties of time-steps}). Here $\ind \bc{\cdot}$ denotes the indicator function.

%a consequence of property (46). 
\textbf{Invertibility of $\Phi_i(\cdot)$:} The above inequality implies that $\frac{\partial {\Phi}_{i}(x)}{\partial x}$ is positive definite uniformly over all $x$. This entails that $\Phi_i(\cdot)$ is \textit{proper}, i.e, $\lim_{\Norm{x}_2 \to \infty} \Norm{\Phi(x)}_2 = \infty$. To see this, we define $\phi(t):=  u^\top \Phi_i(t u)$ for a fixed \textit{unit vector} $u$.Then, we have $\phi^\prime(t) = u^\top \nabla \Phi_i(t u) u > 3/8$. This shows that 
\[
\phi(t) - \phi(0) = \int_{0}^t \phi^\prime(u) \; du  \ge 3t/8.
\]
An application Cauchy-Schwarz inequality yields
\[
\Norm{\Phi_i(tu)}_2 \ge u^\top \Phi_i(t u) \ge 3t/8 + u^\top \Phi_i(0) \ge 3t/8 - \Norm{\Phi_i(0)}_2.
\]
Setting $t = \Norm{x}_2$ and $u = x/\Norm{x}_2$, we have $\Norm{\Phi_i(x)}_2 \ge 3\Norm{x}_2/8 - \Norm{\Phi_i(0)}_2$. This implies that $\Phi_i(\cdot)$ is proper. Then, by Theorem 2.2 of \citet{ruzhansky2015global} shows that $\Phi_i(\cdot)$ is a diffeomorphism.
%According to the inverse mapping theorem, we know that ${\Phi}_{i}$ is a bijection on $\mathcal{E}_{i}$. 
Therefore, we can find a unique $x_{i} $ such that ${\Phi}_{i}\left(x_{i}\right)=x_{i+1}$, which in turn allows us to derive
$$
p_{{\Phi}_{i}\left(X_{i}\right)}\left(x_{i+1}\right)=p_{X_{i}}\left({\Phi}_{i}^{-1}\left(x_{i}\right)\right) \cdot \operatorname{det}\left(\frac{\partial {\Phi}_{i}^{-1}\left(x_{i+1}\right)}{\partial x_{i+1}}\right)=p_{X_{i}}\left(x_{i}\right) \cdot \operatorname{det}\left(\frac{\partial x_{i}}{\partial x_{i+1}}\right)
$$

Consequently, we see that: for any $i \geq 1$ %and any $x_{i+1} \in \widetilde{\Phi}_{i}\left(\mathcal{E}_{i}\right)$
,
\begin{equation}
\label{eq: TV distance p_phi_i and p_i+1}
\begin{aligned}
p_{{\Phi}_{i}\left(X_{i}\right)}\left(x_{i+1}\right) & -p_{X_{i+1}}\left(x_{i+1}\right)=p_{X_{i}}\left(x_{i}\right) \operatorname{det}\left(\frac{\partial x_{i}}{\partial x_{i+1}}\right)-p_{X_{i+1}}\left(x_{i+1}\right) \\
& =\int\left\{p_{X_{t} \mid X_N}\left(x_{i} \mid x_N\right) \operatorname{det}\left(\frac{\partial x_{i}}{\partial x_{i+1}}\right)-p_{X_{i+1} \mid X_N}\left(x_{i+1} \mid x_N\right)\right\} p_{X_N}\left(x_N\right) \mathrm{d} x_N \\
& =\int \Big\{1-\underbrace{\frac{p_{X_{i+1} \mid X_N}\left(x_{i+1} \mid x_N\right)}{p_{X_{i} \mid X_N}\left(x_{i} \mid x_N\right)} \operatorname{det}\left(\frac{\partial x_{i+1}}{\partial x_{i}}\right)}_{=: \mathcal{T}\left(x_{i}, x_N\right)}\Big\} p_{X_{i}, X_N}\left(x_{i}, x_N\right) \operatorname{det}\left(\frac{\partial x_{i}}{\partial x_{i+1}}\right) \mathrm{d} x_N  \\
& =\int\left\{1-\mathcal{T}\left(x_{i}, x_N\right)\right\} p_{X_{i}, X_N}\left(x_{i}, x_N\right) \operatorname{det}\left(\frac{\partial x_{i}}{\partial x_{i+1}}\right) \mathrm{d} x_N
\end{aligned}
\end{equation}

\paragraph{Step 3: Analysis of $\cT(x_i, x_N)$.}
The next step is then to analyze $\mathcal{T}\left(x_{i}, x_N\right)$.
Controlling $\mathcal{T}\left(x_{i}, x_N\right)$ : Given how $X_{i+1}$ and $X_{i}$ are generated, we observe that

\begin{align*}
\mathcal{T}\left(x_{i+1}, x_N\right) & =\frac{p_{X_{i+1} \mid X_N}\left(x_{i+1} \mid x_N\right)}{p_{X_{i} \mid X_{N}}\left(x_{i} \mid x_N\right)} \operatorname{det}\left(\frac{\partial x_{i+1}}{\partial x_{i}}\right)=\frac{\left(\frac{1}{1-t_{i+1}}\right)^{d} \exp \left\{\frac{-\left\|x_{i+1}-t_{i+1}x_N\right\|_{2}^{2}}{2\left(1-t_{i+1}\right)^2}\right\}}{\left(\frac{1}{1-t_i}\right)^{d} \exp \left\{\frac{-\left\|x_{i}-t_i x_N\right\|_{2}^{2}}{2\left(1-t_i\right)^2}\right\}} \operatorname{det}\left(\frac{\partial x_{i+1}}{\partial x_{i}}\right) \\
& =\underbrace{\left(\frac{1-t_i}{1-t_{i+1}}\right)^{d} \operatorname{det}\left(\frac{\partial x_{i+1}}{\partial x_{i}}\right)}_{=: \mathcal{T}_{1}\left(x_{i}, x_N\right)} \underbrace{\exp \left\{\frac{\left\|x_{i}-t_i x_N\right\|_{2}^{2}}{2\left(1-t_i\right)^2}-\frac{\left\|x_{i+1}-t_{i+1}x_N\right\|_{2}^{2}}{2\left(1-t_{i+1}\right)^2}\right\}}_{=: \mathcal{T}_{2}\left(x_{t}, x_{0}\right)}
\end{align*}

leaving us with two terms to control.
Let us first study the term $\mathcal{T}_{1}\left(x_{t}, x_{0}\right)$, towards which we see that

\begin{align*}
\mathcal{T}_{1}\left(x_{t}, x_{0}\right) & =\left(\frac{1-t_i}{1-t_{i+1}}\right)^{d} \operatorname{det}\left(\frac{\partial (x_i + \eta_i \widehat v_i(x_i))}{\partial x_{i}}\right) \\
& =\left(\frac{1-t_i}{1-t_{i+1}}\right)^{d} \operatorname{det}\left(I+\eta_{t} \frac{\partial}{\partial x_{t}} v_i\left(x_{i}\right)+\eta_{i}\left(\frac{\partial}{\partial x_{i}} \widehat v_{i}\left(x_{i}\right)-\frac{\partial}{\partial x_{i}} v_{i}\left(x_{i}\right)\right)\right) \\
& \stackrel{(a)}{=}\left(\frac{1-t_i}{1-t_{i+1}}\right)^{d} \operatorname{det}\left(\left(1-\frac{\eta_{i}}{1-t_i}\right) I+\frac{ \eta_{i} t_i}{(1-t_i)^3} \operatorname{Cov}_{1 \mid t_i}(x)+\eta_{t} \varepsilon_i^{J}\left(x_{i}\right)\right)  \\
& =\operatorname{det}\left(\frac{1 - t_i - \eta_i}{1 - t_{i+1}} I+ \frac{\eta_i t_i}{(1-t_{i+1}) (1-t_i)^2} \operatorname{Cov}_{1 \mid t_i}\left(x_{i}\right)+ \frac{\eta_i(1-t_i)}{(1-t_{i+1})} \varepsilon_i^{J}\left(x_{t}\right)\right)
\\
& \stackrel{(b)}{=} \operatorname{det}\left( I+ \frac{\eta_i t_i}{(1-t_{i+1}) (1-t_i)^2} \operatorname{Cov}_{1 \mid t_i}\left(x_{i}\right)+ \frac{\eta_i(1-t_i)}{(1-t_{i+1})} \varepsilon_i^{J}\left(x_{i}\right)\right).
\end{align*}


 Here, (a) arises from Tweedie's formula, whereas (b) follows since

\begin{equation*}
\frac{1 - t_i - \eta_i}{1 - t_{i+1}} = \frac{1 -t_{i+1}}{1 -t_{i+1}} = 1.
\end{equation*}


Next, we turn attention to the term $\mathcal{T}_{2}\left(x_{t}, x_{0}\right)$, which satisfies

\begin{equation}
\label{eq: logT2}
\begin{aligned}
\log \mathcal{T}_{2}\left(x_{i}, x_{N}\right) & =\frac{\left\|x_{i}-t_i x_N\right\|_{2}^{2}}{2\left(1-t_i\right)^2}-\frac{\left\|x_{i+1}-t_{i+1}x_N\right\|_{2}^{2}}{2\left(1-t_{i+1}\right)^2}\\
& =\frac{\left\|x_{i}-t_i x_N\right\|_{2}^{2}}{2\left(1-t_i\right)^2}-\frac{\left\|x_{i} + \eta_i \widehat v_i(x_i)-t_{i+1}x_N\right\|_{2}^{2}}{2\left(1-t_{i+1}\right)^2}
\end{aligned}
\end{equation}
We focus on the term $x_{i} + \eta_i \widehat v_i(x_i)-t_{i+1}x_N$:
$$
\begin{aligned}
x_{i} + \eta_i \widehat v_i(x_i)-t_{i+1}x_N & =\left(1-\frac{\eta_{i}}{1-t_i}\right) x_{i}+\frac{ \eta_{i}}{1-t_i} \mu_{1 \mid t_i}\left(x_{i}\right) - (t_{i} + \eta_i) x_N + \eta_{i}\left(\widehat v_{i}\left(x_{i}\right)-v_{i}\left(x_{i}\right)\right)
\\
& =\left(1-\frac{\eta_{i}}{1-t_i}\right)\left(x_{i}- t_{i} x_N\right)+\frac{ \eta_{i}}{1-t_i} \mu_{1 \mid t_i}\left(x_{i}\right)-\frac{\eta_{i}}{1-t_i} t_{i} x_N - \eta_i x_N + \eta_{t}\left(\widehat v_{i}\left(x_{i}\right)-v_{i}\left(x_{i}\right)\right) 
\\
& =\left(\frac{1 - t_{i+1}}{1-t_i}\right)\left(x_{i}-t_i x_N\right)+\frac{\eta_i}{1-t_i}\left(\mu_{1 \mid t_i}\left(x_{i}\right)-x_{N}\right)
+
\eta_{i} \varepsilon_{i}^{\rm v}\left(x_{i}\right) 
\end{aligned}
$$

Substitution into \eqref{eq: logT2} yields

\begin{equation}
\label{eq: logT2_expanded}
\begin{aligned}
& \log \mathcal{T}_{2}\left(x_{i}, x_{N}\right)\\
& = -\frac{\eta_i}{(1-t_i)^2 (1-t_{i+1})}(x_i - t_i x_N)^\top(\mu_{1\mid t_i}(x_i) - x_N) - \frac{\eta_i}{(1-t_i)(1-t_{i+1})}(x_i - t_i \mu_{1 \mid t_i}(x_i))^\top \varepsilon_i^{\rm v}(x_i)
\\
& - \frac{\eta_i t_{i+1}}{(1-t_{i+1})^2} \left(\mu_{1\mid t_i}(x_i) - x_N\right) \varepsilon_i^{\rm v}(x_i)
 - \frac{\eta_i^2 \Norm{\mu_{1\mid t_i}(x_i) - x_N}_2^2 }{2(1-t_i)^2 (1-t_{i+1})^2} - \frac{\eta_i^2 \Norm{\varepsilon_i^{\rm v}(x_i)}_2^2}{2(1-t_{i+1})^2}.
\end{aligned}
\end{equation}
Now Define:
\begin{align*}
\xi(x_i, x_N)& := -\frac{\eta_i(x_i - t_i x_N)^\top(\mu_{1\mid t_i}(x_i) - x_N)}{(1-t_i)^2 (1-t_{i+1})} - \frac{\eta_i t_{i+1}\left(\mu_{1\mid t_i}(x_i) - x_N\right) \varepsilon_i^{\rm v}(x_i)}{ (1-t_{i+1})^2} \\
 & - \frac{\eta_i^2 \Norm{\mu_{1\mid t_i}(x_i) - x_N}_2^2 }{2(1-t_i)^2 (1-t_{i+1})^2} + \frac{\eta_i}{(1-t_i)^2 (1-t_{i+1})} \left(t_i + \frac{\eta_i}{2(1-t_{i+1})}\right) \Tr(\Cov_{1\mid t_i})
\end{align*}
It is easy to see that 
\[
\int \xi(x_i, x_N) p_{X_N\mid X_i = x_i}(x_N) \; dx_N = 0.
\]


Thus, based on \eqref{eq: logT2_expanded}, we can further simplify $\log \mathcal{T}_{2}\left(x_{i}, x_{N}\right)$ as follows:

\begin{align*}
\log \mathcal{T}_{2}\left(x_{i}, x_{N}\right)= & 
\xi(x_i, x_N)-\frac{\eta_i (x_i - t_i \mu_{1\mid t_i}(x_i))^\top \varepsilon_i^{\rm v}(x_i)}{(1-t_i)(1-t_{i+1})} - \frac{\eta_i^2 \Norm{\varepsilon_i^{\rm v}(x_i)}_2^2}{2(1-t_{i+1})^2}\\
& - \frac{\eta_i}{(1-t_i)^2 (1-t_{i+1})} \left(t_i + \frac{\eta_i}{2(1-t_{i+1})}\right) \Tr(\Cov_{1\mid t_i})
\end{align*}

Now, for ease of presentation, define
$$
\begin{aligned}
W\left(x_{i}\right) & :=\log \operatorname{det}\left( I+ \frac{\eta_i t_i}{(1-t_{i+1}) (1-t_i)^2} \operatorname{Cov}_{1 \mid t_i}\left(x_{i}\right)+ \frac{\eta_i(1-t_i)}{(1-t_{i+1})} \varepsilon_i^{J}\left(x_{i}\right)\right)
-  \frac{\eta_i^2 \Norm{\varepsilon_i^{\rm v}(x_i)}_2^2}{2(1-t_{i+1})^2} \\
& - \frac{\eta_i}{(1-t_i)^2 (1-t_{i+1})} \left(t_i + \frac{\eta_i}{2(1-t_{i+1})}\right) \Tr(\Cov_{1\mid t_i})
-
\frac{\eta_i (x_i - t_i \mu_{1\mid t_i}(x_i))^\top \varepsilon_i^{\rm v}(x_i)}{(1-t_i)(1-t_{i+1})}
\end{aligned}
$$

Then for any $x_{i} \in \mathcal{X}_{\text {data }}$, it holds that

\begin{equation}
\label{eq: integral inequality of W}
\int_{x_{N}}\left(1-e^{\xi\left(x_{t}, x_{N}\right)}\right) e^{W\left(x_{i}\right)} p_{X_{N} \mid X_{i}}\left(x_{N} \mid x_{i}\right) \mathrm{d} x_{N} \leq-e^{W\left(x_{i}\right)} \int_{x_{N}} \xi\left(x_{i}, x_{N}\right) p_{X_{N} \mid X_{i}}\left(x_{N} \mid x_{i}\right) \mathrm{d} x_{N}=0
\end{equation}

where the inequality results from the elementary inequality $1-e^{x} \leq-x$ for all $x \in \mathbb{R}$. Using the above quantities and \eqref{eq: TV distance p_phi_i and p_i+1} we get: for any $\mathcal{A} \subseteq \mathcal{X}_{\text {data }} \cap \mathcal{E}_{i}$,

\begin{align*}
&\mathbb{P}_{\Phi_{i}\left(X_{i}\right)}(\mathcal{A})-\mathbb{P}_{X_{i+1}}(\mathcal{A}) \\
& =\int_{\mathcal{A}}\left\{1-\mathcal{T}\left(x_{i}, x_N\right)\right\} p_{X_{i}, X_N}\left(x_{i}, x_N\right) \operatorname{det}\left(\frac{\partial x_{i}}{\partial x_{i+1}}\right) \mathrm{d} x_N \; \dx_{i+1}\\
& =\int_{\mathcal{A} \times \mathcal{X}_{\text {data }}}\left\{1-\mathcal{T}\left(x_{i}, x_{N}\right)\right\} p_{ X_{N} \mid X_i = x_i}\left( x_{N}\right) p_{X_i}(x_i)\operatorname{det}\left(\frac{\partial x_{i}}{\partial x_{i+1}}\right) \mathrm{d} x_{N} \;  \dx_{i+1} \\
& =\int_{\Phi_{i}^{-1}(\mathcal{A}) \times \mathcal{X}_{\text {data }}}\left\{1-\mathcal{T}\left(x_i, x_N\right)\right\} p_{ X_{N} \mid X_i = x_i}\left( x_{N}\right) p_{X_i}(x_i)\mathrm{d} x_N \mathrm{~d} x_{i} \\
& =\int_{\Phi_{i}^{-1}(\mathcal{A}) \times \mathcal{X}_{\text {data }}}\left\{1-e^{\xi\left(x_i, x_N\right)} \cdot e^{W\left(x_i\right)}\right\} p_{ X_{N} \mid X_i = x_i}\left( x_{N}\right) p_{X_i}(x_i) \mathrm{d} x_N \mathrm{~d} x_{i}  \\
& = \int_{\Phi_{i}^{-1}(\mathcal{A}) \times \mathcal{X}_{\text {data }}}\left\{\left(1-e^{\xi\left(x_i, x_N\right)}\right) e^{W\left(x_i\right)}+\left(1-e^{W\left(x_i\right)}\right)\right\} p_{ X_{N} \mid X_i = x_i}\left( x_{N}\right) p_{X_i}(x_i)\mathrm{d} x_N \mathrm{~d} x_i \\
& \stackrel{(a)}{\leq} \int_{x_i, x_N \in \Phi_{i}^{-1}(\mathcal{A}) \times \cX_{\rm data}}\left\{1-e^{W\left(x_i\right)}\right\} p_{X_i, X_N}\left(x_i, x_N\right) \mathrm{d} x_i\; \dx_N \\
& \leq \int_{x_i\in \Phi_{i}^{-1}(\mathcal{A}) }-W\left(x_i\right) p_{X_i}\left(x_i       \right) \mathrm{d} x_i, 
\end{align*}
where, ($a$) invokes \eqref{eq: integral inequality of W}, and the last inequality follows from the elementary inequality $1-e^{x} \leq-x$.


\paragraph{Step 4: Analyzing $W(x_i)$.}
First, we introduce the following lemma from \cite{liang2025low}.

\begin{lemma}[Lemma 5, \citet{liang2025low}]
\label{lemma: lower bound on log det}
    Let $A\in \bbR^{d \times d}$ be a positive-definite matrix, and $\Delta \in \bbR^{d \times d}$ be any square matrix. Suppose $\eta \Norm{\Delta}_2 \le 1/4$, where $ \eta \in (0,1)$. Then it holds that 
    \[
    \log \det (I + \eta A + \eta \Delta ) \ge \eta (\Tr(A) + \Tr(\Delta)) - 4 \eta^2 (\Norm{A}_F^2 + \Norm{\Delta}_F^2).
    \]
\end{lemma}
Using the above lemma we have 
\begin{align*}
& - \log \operatorname{det}\left( I+ \frac{\eta_i t_i}{(1-t_{i+1}) (1-t_i)^2} \operatorname{Cov}_{1 \mid t_i}\left(x_{i}\right)+ \frac{\eta_i(1-t_i)}{(1-t_{i+1})} \varepsilon_i^{J}\left(x_{i}\right)\right) \\
& \le  \frac{4 \eta_i^2 t_i^2}{(1-t_{i+1})^2 (1-t_i)^4} \Norm{\Cov_{1 \mid t_i}}_F^2 + \frac{4 \eta_i^2 (1-t_i)^2}{(1-t_{i+1})^2} \Norm{\varepsilon_i^J(x_i)}_F^2\\
& \quad - \frac{\eta_i t_i}{(1-t_{i+1}) (1-t_i)^2} \Tr(\Cov_{1 \mid t_i}) - \frac{\eta_i (1-t_i)}{(1-t_{i+1})} \Tr(\varepsilon_i^J(x_i)),
\end{align*}
provided that the inequality

$$\frac{\eta_i (1-t_i)}{(1-t_{i+1})} \Norm{\varepsilon_i^J(x_i)}_2\le \frac{1}{4}
$$
holds for all $i$. This is indeed the case in our setting, due to Assumption \ref{assumption: higher-order approximation error}\ref{assumption: Jone},  and 
$$\frac{1-t_i}{1 - t_{i+1}} = \frac{1}{1 - \frac{\eta_i}{1 - t_i}} \le 2$$
due to the inequality $\frac{\eta_i}{1 - t_i} \le 2h \le 1/2$, where $h$ is defined in \eqref{eq: gemoetric time discretization}.
This yields 
\begin{equation}
\label{eq: bound on W}
\begin{aligned}
- W\left(x_{i}\right) & = - \log \operatorname{det}\left( I+ \frac{\eta_i t_i}{(1-t_{i+1}) (1-t_i)^2} \operatorname{Cov}_{1 \mid t_i}\left(x_{i}\right)
+ \frac{\eta_i(1-t_i)}{(1-t_{i+1})} \varepsilon_i^{J}\left(x_{i}\right)\right)
+  \frac{\eta_i^2 \Norm{\varepsilon_i^{\rm v}(x_i)}_2^2}{2(1-t_{i+1})^2} \\
& + \frac{\eta_i}{(1-t_i)^2 (1-t_{i+1})} \left(t_i + \frac{\eta_i}{2(1-t_{i+1})}\right) \Tr(\Cov_{1\mid t_i})
+
\frac{\eta_i (x_i - t_i \mu_{1\mid t_i}(x_i))^\top \varepsilon_i^{\rm v}(x_i)}{(1-t_i)(1-t_{i+1})}\\
& \le \frac{4 \eta_i^2 t_i^2}{(1-t_{i+1})^2 (1-t_i)^4} \Norm{\Cov_{1 \mid t_i}(x_i)}_F^2 + \frac{4 \eta_i^2 (1-t_i)^2}{(1-t_{i+1})^2} \Norm{\varepsilon_i^J(x_i)}_F^2 + \frac{\eta_i^2 \Norm{\varepsilon_i^{\rm v}(x_i)}_2^2}{2(1-t_{i+1})^2} \\
& \quad + \frac{\eta_i (x_i - t_i \mu_{1\mid t_i}(x_i))^\top \varepsilon_i^{\rm v}(x_i)}{(1-t_i)(1-t_{i+1})} +  \frac{\eta_i}{(1-t_i)^2 (1-t_{i+1})} \left(t_i + \frac{\eta_i}{2(1-t_{i+1})} -  t_i\right) \Tr(\Cov_{1\mid t_i}(x_i) )\\
& \quad - \frac{\eta_i^2 (1-t_i)}{(1-t_{i+1})} \Tr(\varepsilon_i^J(x_i))\\
& \le \frac{4 \eta_i^2 t_i^2}{(1-t_{i+1})^2 (1-t_i)^4} \Norm{\Cov_{1 \mid t_i}(x_i)}_F^2 + \frac{4 \eta_i^2 (1-t_i)^2}{(1-t_{i+1})^2} \Norm{\varepsilon_i^J(x_i)}_F^2 + \frac{\eta_i^2 \Norm{\varepsilon_i^{\rm v}(x_i)}_2^2}{2(1-t_{i+1})^2} \\
& \quad  +  \frac{\eta_i^2}{2(1-t_i)^2 (1-t_{i+1})^2} \Tr(\Cov_{1\mid t_i}(x_i) ) + \underbrace{\frac{\eta_i (x_i - t_i \mu_{1\mid t_i}(x_i))^\top \varepsilon_i^{\rm v}(x_i)}{(1-t_i)(1-t_{i+1})}}_{\Delta(\varepsilon_i^{\rm v}(x_i))}\\
& \quad - \frac{\eta_i^2 (1-t_i)}{(1-t_{i+1})} \Tr(\varepsilon_i^J(x_i))\\
& 
\le \frac{4 \eta_i^2 t_i^2}{(1-t_{i+1})^2 (1-t_i)^4} \Norm{\Cov_{1 \mid t_i}(x_i)}_F^2 + \frac{4 \eta_i^2 (1-t_i)^2}{(1-t_{i+1})^2} \Norm{\varepsilon_i^J(x_i)}_F^2 + \frac{\eta_i^2 \Norm{\varepsilon_i^{\rm v}(x_i)}_2^2}{2(1-t_{i+1})^2} \\
& \quad  +  \frac{\eta_i^2}{2(1-t_i)^2 (1-t_{i+1})^2} \Tr(\Cov_{1\mid t_i}(x_i) ) + \underbrace{\frac{\eta_i (1-t_i)}{(1-t_{i+1})}\frac{(x_i - t_i \mu_{1\mid t_i}(x_i))^\top \varepsilon_i^{\rm v}(x_i)}{(1-t_i)^2}}_{\Delta(\varepsilon_i^{\rm v}(x_i))}\\
& \quad - \frac{\eta_i^2 (1-t_i)}{(1-t_{i+1})} \Tr(\varepsilon_i^J(x_i))
\end{aligned}
\end{equation}

\paragraph{Step 5: Controlling the error arising from $\Delta(\varepsilon_i^{\rm v}(x_i))$.}
We start with the following integral quantity for any set $\cA$:
\begin{equation}
\label{eq: Delta(e_i) simplified}
\begin{aligned}
    & \int_{x_i \in \cA} \frac{1}{(1-t_i)^2}(x_i - t_i \mu_{1 \mid t_i}(x_i))^\top \varepsilon^{\rm v}_i(x_i)p_{X_i}(x_i) \; \dx_i\\
    & = \int_{x_i, x_N \in \cA \times \cX_{\rm data}} \frac{1}{(1-t_i)^2}(x_i - t_i \mu_{1 \mid t_i}(x_i))^\top \varepsilon^{\rm v}_i(x_i)p_{X_i, X_N}(x_i, x_N) \; \dx_i\; \dx_N\\
    %&\rd  = \int_{x_i, x_N \in \cA \times \cX_{\rm data}} \frac{1}{(1-t_i)^2}(x_i - t_i \mu_{1 \mid t_i}(x_i))^\top \varepsilon^{\rm v}_i(x_i)
    % p_{X_N \mid X_i = x_i}(x_N) p_{X_i}(x_i) \; \dx_N\; \dx_i\\
    &  = \int_{x_i, x_N \in \cA \times \cX_{\rm data}} \frac{1}{(1-t_i)^2}(x_i - t_i \mu_{1 \mid t_i}(x_i))^\top \varepsilon^{\rm v}_i(x_i)
    p_{X_i \mid X_N = x_N}(x_i) p_{X_N}(x_N) \; \dx_i\; \dx_N
    \\
    & = \int_{x_i, x_N \in \bbR^d \times \cX_{\rm data}} \frac{1}{(1-t_i)^2}\left\{(x_i - t_i x_N)^\top \varepsilon^{\rm v}_i(x_i) p_{X_i \mid X_N = x_N}(x_i)\; \dx_i \right\} p_{X_N}(x_N) \; \dx_N\\
    & \le \int_{x_i, x_N \in \bbR^d \times \cX_{\rm data}} \frac{1}{(1-t_i)^2}\left\{\left\vert(x_i - t_i x_N)^\top \varepsilon^{\rm v}_i(x_i)\right\vert p_{X_i \mid X_N = x_N}(x_i)\; \dx_i  \right\} p_{X_N}(x_N) \; \dx_N \\
    & \le \left(\int_{x_i, x_N \in \bbR^d \times \cX_{\rm data}} \left\{\frac{1}{(1-t_i)^4}\left\vert(x_i - t_i x_N)^\top \varepsilon^{\rm v}_i(x_i)\right\vert ^2 p_{X_i \mid X_N = x_N}(x_i)\; \dx_i  \right\} p_{X_N}(x_N) \; \dx_N\right)^{1/2}\\
    &  = \left(\int_{x_i, x_N \in \bbR^d \times \cX_{\rm data}} \left\{\frac{1}{(1-t_i)^4}\left \langle (x_i - t_i x_N) (x_i - t_i x_N)^\top , \varepsilon^{\rm v}_i(x_i) \varepsilon^{\rm v}_i(x_i)^\top \right\rangle  p_{X_i \mid X_N = x_N}(x_i)\; \dx_i  \right\} p_{X_N}(x_N) \; \dx_N\right)^{1/2} 
\end{aligned}
\end{equation}
Now, recall that $X_i \mid X_N = x_N \sim N(t_i x_N , (1-t_i)^2 I)$. Using standard techniques from \citet[Section B.5]{liang2025low}, we can obtain bound on \eqref{eq: Delta(e_i) simplified}:
%{\color{red} (SR: Check this) } 
\[
\int_{x_i \in \cA} \frac{1}{(1-t_i)^2}(x_i - t_i \mu_{1 \mid t_i}(x_i))^\top \varepsilon^{\rm v}_i(x_i)p_{X_i}(x_i) \; \dx_i \le \frac{2 \evel_i}{(1-t_i)} + 
\evelJone_i + \evelJtwo_i + \evelH_i.
\]
Substituting this in \eqref{eq: bound on W}, we have 
\begin{align*}
& \int_{x_i \in \cA}
\left\{\Delta(\evel_i(x_i))- \frac{\eta_i^2 (1-t_i)}{(1-t_{i+1})} \Tr(\varepsilon_i^J(x_i))\right\} p_{X_i}(x_i) \\
& \le  \frac{2 \evel_i}{(1-t_i)} + 
\evelJone_i + \evelJtwo_i + \evelH_i +  \frac{\eta_i^2 (1-t_i)}{(1-t_{i+1})} \br{\int \Tr(\varepsilon_i^J(x_i))^2 p_{X_i}(x_i)\dx_i}^{1/2}\\
& \le \frac{2 \eta_i }{(1-t_{i+1})}( \evel_i + \evelJone_i + \evelJtwo_i + \evelH_i) .
\end{align*}
Here, we used the fact that $\frac{\eta_i^2 (1-t_i)}{(1 - t_{i+1})} \le \frac{\eta_i}{(1-t_{i+1})}$. 


\paragraph{Step 6: Bounding the terms involving $\Cov_{1 \mid t_i}(x_i)$.} 
Here we always assume that $k\ge \log d$.\footnote{If $k <\log d$, then we can redefine $k$ as $\log d$. This does not change the results in Theorem \ref{thm: RF - convergence of Eulerian scheme} and Theorem \ref{thm: stoc-RF sampler non-Eulerian}.
}
We begin defining some quantities first. 
\begin{itemize}
    \item Let $\{x_i^\ast\}_{1 \le i \le N_{\epsilon_0}}$ be an $\epsilon_0$-net of $\mathcal{X}_{\mathrm{data}}$, with $N_{\epsilon_0}$ denoting its cardinality. Let $\{B_i\}_{1 \le i \le N_{\epsilon_0}}$ be a disjoint $\epsilon_0$-cover for $\mathcal{X}_{\mathrm{data}}$ such that $x_i^\ast \in B_i$ for each $i$. 
    
    \item Define the following two sets:
    \begin{equation}
        \mathcal{I} := \left\{ 1 \le i \le N_{\epsilon_0} : \mathbb{P}(X_0 \in B_i) \ge \exp(-C_1 k \log T) \right\}
        \label{eq:I}
    \end{equation}
    and
    \begin{equation}
        \mathcal{G} := \left\{ w \in \mathbb{R}^d : \|w\|_2 \le 2\sqrt{d} + \sqrt{C_1 k \log T}, \;
        \left| (x_i^\ast - x_j^\ast)^\top w \right|
        \le \sqrt{C_1 k \log T}\, \|x_i^\ast - x_j^\ast\|_2,
        \; \forall\, 1 \le i,j \le N_{\varepsilon_0} \right\}.
        \label{eq:G}
    \end{equation}
    for some sufficiently large universal constant $C_1 > 0$. As it turns out, $\bigcup_{i \in \mathcal{I}} B_i$ and $\mathcal{G}$ form certain high-probability sets related to the random vector $X_0 \sim p_{\mathrm{data}}$ and the standard Gaussian random vector in $\mathbb{R}^d$, respectively.
    
    \item Recall \eqref{eq: interp in modified index}, we can express
    \begin{equation}
        X_{t_i} \equiv X_{i} = t _i  X_N + (1 - t_i) X_0,
        \label{eq:xt}
    \end{equation}
    for some random vector $X_0 \sim \mathcal{N}(0, I_d)$. 
    % By introducing
    % \begin{equation}
    %     V_\alpha := \sqrt{\alpha}\, V_1 + \sqrt{1 - \alpha}\, Z
    %     \qquad \text{with } V_1 := X_0,
    %     \label{eq:valpha}
    % \end{equation}
\end{itemize}
Next, we define the set  $\cT_t := \left\{t v_1 + (1-t) \omega: v_1 \in \cup_{i \in \cI}\cB_i, \omega \in \cG \right\}$. %\red{(SR: need to define things here.)}

Borrowing standard techniques and results from the proofs of \citet[Lemma 1]{liang2025low} and \citet[Corollary 1]{liang2025low}, we have the following results

\begin{equation}
    \label{eq: prob of exclusion}
    \pr\br{X_{i} \notin \cT_{t_i} } \le \exp\br{- C_2 k \log N}, \quad \text{for $C_2 \gg C_{\rm cover}$, and $0 \le t_i<1$.}
\end{equation}
\begin{equation}
    \label{eq: trace_Cov bound}
    \Tr(\Cov_{1 \mid t_i}(x_i) )\le \frac{C_3 (1 - t_i)^2}{t_i^2} k \log N, \quad \text{for all $t_i>0$ and $x_i \in \cT_{t_i}$.}
\end{equation}


\paragraph{Step 7: Bound on the term involving squared Frobenius norm using linearization:}
One can use \eqref{eq:  trace_Cov bound} directly to bound the term $\bbE_{X_i}\left(\Norm{\Cov_{1\mid t_i}(X_i)}_F^2\right)$ in \eqref{eq: bound on W}. However, this will lead to a quadratic dependence on $k$. Therefore, we present the following lemma that will yield a linear dependence on $k$. The proof is deferred to Appendix \ref{app: proof of Cov Frobenius bound}.
% One can use \eqref{eq:  trace_Cov bound} directly to bound the term in 
% We present the following bound on $\bbE_{X_i}\left(\Norm{\Cov_{1\mid t_i}(X_i)}_F^2\right)$.
\begin{proposition}
\label{prop: Main cov bound via trace}
    For all $i\ge 1$, the following holds:
    \begin{equation*}
    %\label{eq: Main cov bound via trace}
    \frac{ \eta_{i}t_i}{(1-t_{i+1})^2 (1-t_i)^2}  \bbE(\Norm{\Cov_{1\mid t_i}(X_i)}_F^2)  \le \frac{3}{4} \left(\bbE [\Tr(\Cov_{1\mid t_i}(X_i))] - \bbE[\Tr(\Cov_{1\mid t_{i+1}}(X_{i+1}))]\right) + \frac{6 h \eta_{i} }{(1-t_{i+1})^2 N^{10}}.
\end{equation*}
    
\end{proposition}
The 

%Also, note that $\Norm{\Cov_{1\mid t_i}(x_i)}_F^2 \le \left(\Tr(\Cov_{1 \mid t_i}(x_i) )\right)^2$. Putting all of the above fact together, we have 
Using Proposition \ref{prop: Main cov bound via trace} and \eqref{eq: trace_Cov bound}, we have
\begin{align*}
& \mathbb{P}_{\Phi_{i}\left(X_{i}\right)}(\mathcal{A})-\mathbb{P}_{X_{i+1}}(\mathcal{A}) \\
% &= \pr(\Phi_i(X_i) \in \cA, X_i \in \cE_i) - \pr(X_{i+1} \in \cA) + \pr (X_i \notin \cE_i)\\
% & \le \pr(\Phi_i(X_i) \in \cA, \Phi_i(X_i) \in \Phi_i(\cE_i) ) - \pr(X_{i+1} \in \cA) + \pr(X_i \notin \cE_i)\\
% & \le \pr (\Phi_i(X_i) \in \cA \cap \Phi_i(\cE_i)) - \pr(X_{i+1} \in \cA \cap \Phi_i(\cE_i)) + \pr(X_i \notin \cE_i)\\
 & \le \frac{3 \eta_i t_i }{ (1-t_i)^2} \left(\frac{(1-t_i)^2}{t_i^2} - \frac{(1-t_{i+1})^2}{t_{i+1}^2}\right) k \log N + \frac{\eta_i^2}{2(1-t_{i+1})^2 t_i^2} k \log N + \frac{4 \eta_i^2 (1-t_i)^2}{(1-t_{i+1})^2} (\evelJone_i)^{2} + \frac{\eta_i^2}{2(1-t_i)^2}(\evel_i)^2\\
  &\quad  + \frac{2 \eta_i }{(1-t_{i+1})}( \evel_i + \evelJone_i + \evelJtwo_i + \evelH_i) + \pr(X_i \notin \cE_i) + \frac{24 h \eta_i^2 t_i}{(1-t_{i+1})^2(1-t_i)^2 N^{10}}.
\end{align*}

First, Lemma \ref{lemma: simple identity 1} yields that 
\begin{equation}
\label{eq: simple identity 1}
\frac{(1-t_i)^2}{t_i^2} - \frac{(1-t_{i+1})^2}{t_{i+1}^2} \le \frac{2 \eta_i}{t_i^3}.
\end{equation}
% We also use the Markov's inequality to control the following:
% \[
% \pr(X_i \notin \cE_i) = \pr\left\{\frac{\eta_i (1-t_i)}{(1-t_{i+1}) } \Norm{\varepsilon_i^{\rm J}(x_i)} > \frac{1}{8}\right\} \le \frac{64 \eta_i^2 (1-t_i)^2}{(1-t_{i+1})^2}\bbE\left\{ \Norm{\varepsilon_i^{\rm J}(x_i)}^2\right\} \le \frac{64 \eta_i^2 (1-t_i)^2}{(1-t_{i+1})^2} (\evelJone_i)^2. 
% \]
Using the above inequality along with \eqref{eq: simple identity 1} yields 
\begin{align*}
  & \mathbb{P}_{\Phi_{i}\left(X_{i}\right)}(\mathcal{A})-\mathbb{P}_{X_{i+1}}(\mathcal{A}) \\
&  \le \frac{6 C_3 \eta_i^2}{(1-t_{i+1})^2 t_i^2} k \log N + \frac{C_3 \eta_i^2}{2(1-t_{i+1})^2 t_i^2} k \log N + \frac{C \eta_i^2 (1-t_i)^2}{(1-t_{i+1})^2} (\evelJone_i)^{2} + \frac{\eta_i^2}{2(1-t_i)^2}(\evel_i)^2\\
  &\quad  + \frac{2 \eta_i }{(1-t_{i+1})}( \evel_i + \evelJone_i + \evelJtwo_i + \evelH_i) + \frac{24 h \eta_i^2 t_i}{(1-t_{i+1})^2(1-t_i)^2 N^{10}} .
\end{align*}
\paragraph{Recalling properties of  time scheduling \eqref{eq: gemoetric time discretization} (see Lemma \ref{lemma: properties of time-steps}):}
\begin{enumerate}  
    \item For all $i>1$, we have $\frac{\eta_i}{1-t_i} \le h$ and $\frac{\eta_i}{1-t_{i+1}} \le h$.
    
     \item For all $i>1$, we have $\frac{ \eta_i^2 (1-t_i)^2 }{(1-t_{i+1})^2} = \frac{(1-t_i)^2}{(\frac{1-t_i}{\eta_i} - 1)^2} \le  h^2 $, as $\frac{\eta_i}{(1-t_i)} \le h$ and $(1-t_i) \le 1$.
     
     \item $\sum_{i : t_i >0}\frac{ \eta_i^2 }{(1-t_{i+1})^2 t_i^2} = \sum_{i: 0 < t_{i} < 1/2 }\frac{ \eta_i^2 }{(1-t_{i+1})^2 t_i^2} + \sum_{i: t_{i}\ge 1/2 }\frac{ \eta_i^2 }{(1-t_{i+1})^2 t_i^2} \le 4 h^2 N$.
\end{enumerate}


Putting all of these together and solving \eqref{eq: TV recursion}, we get
\begin{equation*}
    \begin{aligned}
        & \TV(p_{X_{N-1}}, p_{Y_{N-1}})\\
        & \le C_1^\prime h^2 N k \log N + + C_2^\prime h^2  N \meanevelJone^2 + C_3^\prime h^2 N\meanevel^2 \\
        & \quad + C_4^\prime h \times \sum_{i = 1}^{N-1} (\evel_i + \evelJone_i + \evelJtwo_i + \evelH_i) + C_5^\prime \frac{h^2}{\delta^2 N^{10}} + \TV(p_{X_1}, p_{Y_1})\\
        & \le C_1^\prime h^2 N k \log N  + C_2^\prime h^2  N \meanevelJone^2 + C_3^\prime h^2 N\meanevel^2 \\
        & \quad + C_4^\prime h N (\meanevel + \meanevelJone + \meanevelJtwo + \meanevelH) + C_5^\prime \frac{h^2}{\delta^2 N^{10}}+ \TV(p_{X_1}, p_{Y_1}),
    \end{aligned}
\end{equation*}
for sufficiently large $N$.
\paragraph{Handling $\TV(p_{X_1}, p_{Y_1})$ separately:}
Recall that \eqref{eq: bound on W} is also valid for $t_{i+1} = \delta, t_i = 0$. Therefore, we have 
\begin{align*}
& - W(x_0) \\
& \le \frac{4 \delta^2 }{(1-\delta)^2} \Norm{\varepsilon_0^J(x_0)}_F^2 + \frac{\delta^2 \Norm{\varepsilon_0^{\rm v}(x_0)}_2^2}{2(1-\delta)^2}  +  \frac{\delta^2}{2(1-\delta)^2} \tr(\Cov_{1} ) + \underbrace{\frac{\delta}{1 - \delta} x_0 ^\top \varepsilon_0^{\rm v}(x_0)}_{\Delta(\varepsilon_0^{\rm v}(x_0))}\\
& \quad - \frac{\delta^2 }{(1-\delta)} \tr(\varepsilon_0^J(x_0)). 
\end{align*}
Therefore, applying similar argument, we get 

\begin{align*}
    \TV(p_{X_1}, p_{Y_1}) \le A_1 \delta^2 \tr(\Cov_1) + A_2 \delta^2 (\evel_0)^2 + A_3 \delta^2 (\evelJone_0)^2 + A_4 (\evel_0 + \evelJone_0 + \evelJtwo_0 + \evelH_0) + \underbrace{\TV(p_{X_0},p_{ Y_0})}_{=0}.
\end{align*}

% We recall that we need $\delta \alpha^{N/2} = 1/2$. Now, we set $\alpha = 1 + h$ and $\delta = 1/N$, which entails that
% \begin{equation}
% \label{eq: value of N}
% N = \frac{2\log (N/2)}{\log (1 + h)} \le \frac{4 \log(N/2)}{h}. \quad \text{(as $\log(1 + h)\ge h/2$ for $h \in (0, 1]$)}
% \end{equation}
Now we recall that $\delta < 1/N$.
Therefore, combining all of the above inequalities and \eqref{eq: value of h} gives 
\begin{align*}
    & \TV(p_{X_{N-1}}, p_{Y_{N-1}})\\
    & \le C_1''  \frac{k \log^3 (1/\delta)  }{N} + C_2''  \frac{\log^2 (1/\delta)}{N}   \meanevelJone^2 + C_3'' \frac{\log^2 (1/\delta)}{N} \meanevel^2 \\
    & \quad + C_4'' \log (1/\delta) (\meanevel + \meanevelJone + \meanevelJtwo + \meanevelH) + C_5'' \delta^{-2} N^{-10}\\
     &\quad +   A_1 \delta^2 {\Tr(\Cov_1)} + A_2 \delta^2 (\evel_0)^2 + A_3 \delta^2 (\evelJone_0)^2 + A_4 (\evel_0 + \evelJone_0 + \evelJtwo_0 + \evelH_0) , 
\end{align*}
where $\Cov_1:= \Cov_{X_1 \sim p_1}(X_1)$.
Now, note that one can estimate $v_0(x) = \mu_{p_{X_1}} - x$ by $\widehat v_0(x):= \widehat \mu_1 - x$, where $\widehat \mu_1 = \sum_{i=1}^n x_{i,1}/n$ and $n$ is the (training) sample size. In this case, we have 
\begin{itemize}
    \item $\evelJone_0 = \evelJtwo_0 = \evelH = 0$.
    \item $(\evel_0)^2 = \bbE \Norm{\widehat \mu_1 - \mu_{p_{X_1}}}_2^2 = O\left( \frac{\tr(\Cov_1)}{n }\right)$.
\end{itemize}
This finishes the proof of Theorem \ref{thm: RF - convergence of Eulerian scheme}.



%%%%%%%%%%%%%%%%%%%%%%%%
%%% Proof of Stoc-RF %%%
%%%%%%%%%%%%%%%%%%%%%%%%

\subsection{Proof of convergence of STOC-RF (Theorem \ref{thm: stoc-RF sampler non-Eulerian})}
\label{app: Stoc-RF convergence}
The proof of this result essentially follows by reduction and leverages key results of \cite{liang2025low}. We start by recalling some key facts:
 $t(\tau) = \frac{\sqrt{\omega_\tau}}{\sqrt{1 - \omega_\tau} + \sqrt{\omega_\tau}}$, and $\frac{Y^\prime_\tau}{\sqrt{\omega_\tau}} \overset{d}{=} \frac{\tilde Z_{t(\tau)}}{t(\tau)}$. With this time transformation, we have 
 \begin{equation}
 \label{eq: alpha_t expression stoc-RF}
 \alpha_\tau = \frac{\omega_\tau}{\omega_{\tau -1}} = \frac{(t_i/ \sigma_{t_i})^2}{(t_{i+1}/\sigma_{t_{i+1}})^2} = \frac{R_i^2}{R_{i+1}^2}, \quad \text{ where $t_i = t(\tau)$}.
 \end{equation}
 We also have $s_{Y^\prime_\tau }(y) = \sigma_t s_{\tilde Z_t}(\tilde z)$, where $\tilde z = \sigma_t y$.
 Next, we recall the DDIM sampler considered in \citet{liang2025low}:  

\begin{equation}
    \label{eq: DDPM sampler liang2025 2}
    \begin{aligned}
    & \hat Y^\prime_{\tau -1} = \frac{1}{\sqrt{\alpha_\tau}} \bc{ \hat Y^\prime_\tau + \delta_\tau \widehat s_{Y^\prime_\tau}(\hat Y^\prime_\tau) + \nu_\tau \xi_\tau}, \hat Y^\prime_N \sim N(0, I_d)\\
    & \delta_\tau = 1 - \alpha_\tau,  \nu_\tau = \sqrt{\frac{(\alpha_\tau - \omega_\tau)(1 - \alpha_t)}{1 - \omega_\tau}},\\
    & \tau   = N , N-1, \ldots, 1, \quad \{\xi_\tau\}_{\tau \ge 1} \overset{i.i.d.}{\sim} N(0,I_d),
    \end{aligned}
\end{equation}

Now, if we substitute all the quantities of the above sampler by their corresponding RF versions, i.e.,  we get 
$\hat Y_\tau \rightarrow  Y_{t}/\sigma_t, \widehat s_{Y^\prime_\tau }(y) \rightarrow \sigma_t \widehat s_{t_i}(Y_{t_i})$, and $\alpha_\tau \rightarrow \frac{R_i^2}{R_{i+1}^2} $, then it leads to the following sampler:
\begin{equation}
     %\label{eq: stoc-RF sampler non-Eulerian}
     \begin{aligned}
      &\frac{Y_{t_{i+1}}}{\sigma_{t_{i+1}}} = \frac{R_{i+1}}{R_i} \bc{ \frac{Y_{t_i}}{\sigma_{t_i}} + \eta_i \sigma_{t_i} \widehat s_{t_i}(Y_{t_i}) + \sqrt{\psi_i} W_{t_i}},\\
      &  \frac{Y_{t_0}}{\sigma_{t_0}} \sim N(0,  I_d), \quad \{W_{t_i}\}_{i \ge 0}\overset{i.i.d.}{\sim} N(0,I_d),\\
      & i = 0, 1, \ldots, N-1,
     \end{aligned}
 \end{equation}
 where $ \eta_i = 1- \frac{R_i^2}{R_{i+1}^2},   \psi_i = \frac{R_i^2}{R_{i+1}^2}.\frac{ 1 - R_{i+1}^2}{1  - R_i^2}.\br{1 - \frac{R_i^2}{R_{i+1}^2}}$, and $\widehat s_{t_i}(x) = \frac{t_i \widehat v_{t_i}(x) - x}{1 - t_i}$. The above is exactly the sampler described in \eqref{eq: stoc-RF sampler non-Eulerian}. In other words, we have equivalence between two above samplers. Therefore, we can use the convergence result of \eqref{eq: DDPM sampler liang2025 2}. 

 \paragraph{Defining few auxiliary processes:}
 We define the auxiliary processes in exact way defined in \citet[Section C.1]{liang2025low}.
 \begin{itemize}
    \item First, we define $\{Y^\star_\tau\}_{\tau = N}^1$ by 
    \[
    Y^\star_N = Y^\prime_N \sim N(0, I_d), \quad Y^\star_{\tau - 1} = \frac{1}{\sqrt{\alpha_\tau}} \br{Y^\star_\tau + \delta_\tau s_{Y^\prime_\tau}(Y^\star_\tau) + \nu_\tau \xi_\tau}.
    \]
Note that, the above update rule is just the analog of $\{\hat Y_\tau\}_\tau$ process with the true score.
 
  \item We construct an auxiliary process $\bar{Y}_{\tau}$ that follows the same transition dynamics as $Y_{\tau}^{\star}$:
\end{itemize}

\begin{equation}
\label{eq: aux process bar_Y}
\bar{Y}_{\tau - 1}^{-}\left|\bar{Y}_{\tau} \sim Y_{\tau - 1}^{\star}\right| Y_{\tau}^{\star}, \quad \bar{Y}_{\tau - 1} \left\lvert\,\left\{\bar{Y}_{\tau - 1}^{-}=y_{\tau - 1}^{-}\right\}= \begin{cases}y_{\tau - 1}^{-}, & \text {with prob. } \frac{p_{Y^\prime_{\tau - 1}}\left(y_{\tau - 1}^{-}\right)}{p_{\bar{Y}_{\tau - 1}^{-}}\left(y_{\tau - 1}^{-}\right)} \wedge 1  \\ \infty, & \text { otherwise }\end{cases}\right.
\end{equation}

for any $y_{\tau - 1}^{-} \neq \infty$, where we recall that $a \wedge b:=\min \{a, b\}$. It is straightforward to show that


\begin{equation}
\label{eq: pdf of aux process bar_Y}
p_{\bar{Y}_{\tau}}\left(y_{\tau}\right)=\int_{\mathbb{R}^{d}}\left(p_{Y^\prime_{\tau}}\left(y_{\tau}^{-}\right) \wedge p_{\bar{Y}_{\tau}^{-}}\left(y_{\tau}^{-}\right)\right) \delta\left(y_{\tau}-y_{\tau}^{-}\right) \mathrm{d} y_{\tau}^{-}=p_{Y^\prime_{\tau}}\left(y_{\tau}\right) \wedge p_{\bar{Y}_{\tau}^{-}}\left(y_{\tau}\right) 
\end{equation}


for any $y_{\tau} \neq \infty$, where $\delta(\cdot)$ denotes the Dirac measure.

\begin{itemize}
  \item To account for the score estimation error, we introduce another auxiliary process $\widetilde{Y}_{\tau}$ based on the dynamics of $Y_{\tau}$ :
\end{itemize}

\begin{equation}
\label{eq: aux process hat_Y}
\widetilde{Y}_{\tau - 1}^{-}\left|\widetilde{Y}_{\tau} \sim \hat Y^\prime_{\tau - 1}\right| \hat Y^\prime_{\tau}, \quad \widetilde{Y}_{\tau - 1} \left\lvert\,\left\{\widetilde{Y}_{\tau - 1}^{-}=y_{\tau - 1}^{-}\right\}= \begin{cases}y_{\tau - 1}^{-}, & \text {with prob. } \frac{p_{Y^\prime_{\tau - 1}}\left(y_{\tau - 1}^{-}\right)}{p_{\bar{Y}_{\tau - 1}^{-}}\left(y_{\tau - 1}^{-}\right)} \wedge 1 \\ \infty, & \text { otherwise. }\end{cases}\right.
\end{equation}

% It is seen that the probability densities of $Y_{\tau}$ and $\widehat{Y}_{i}$ satisfy the properties stated in the following lemma, whose proof can be found in \red{Appendix C.3.}


 Next, we observe that for $t_i = t(\tau)$ 
 % \begin{equation}
 %     \label{eq: score error equivalence}
 %      (\varepsilon^s_\tau)^2 := \bbE \Norm{\widehat s_{Y^\prime_{\tau}}(Y^\prime_{\tau}) - s_{Y^\prime_{\tau}}(Y^\prime_{\tau})}_2^2  = \frac{\sigma_{t_i}^2 t_i^2}{(1-t_i)^2}\bbE \Norm{ \widehat v_{t_i}(X_{t_i}) - v_{t_i}(X_{t_i})}_2^2  \le \frac{t_i^2}{(1-t_i)^2}(\evel_i)^2 \le \frac{\omega_{\tau_i}}{1 - \omega_{\tau_i}} (\evel_i)^2.
 % \end{equation}
 \begin{equation}
     \label{eq: score error equivalence}
      (\varepsilon^s_\tau)^2 := \bbE \Norm{\widehat s_{Y^\prime_{\tau}}(Y^\prime_{\tau}) - s_{Y^\prime_{\tau}}(Y^\prime_{\tau})}_2^2  = \sigma_{t_i}^2 \bbE \Norm{ \widehat s_{t_i}(X_{t_i}) - s_{t_i}(X_{t_i})}_2^2  \le (\escore_i)^2.
 \end{equation}
 Now, we can exactly replicate the analysis in \citet[Section C.2]{liang2025low} till Equation (107) in Step 6, which yields the following for a universal constant $L>0$:

 \[
 \KL\br{p_{\bar Y_1} \| p_{\widetilde{Y}_1} } \le \sum_{\tau = 2}^N L(1 - \omega_\tau) (\varepsilon_\tau^s)^2 \lesssim \frac{\log N}{N}\sum_{t_i = t(2)}^{t(N)} (\escore_i)^2 \lesssim  \meanescore^2 \log N.
 \]
 After this, Step 7 of the proof follows verbatim to yield the final bound in Theorem \ref{thm: stoc-RF sampler non-Eulerian}.


\subsection{Description of the new RF sampler in Remark \ref{remark: new RF sampler}}
\label{app: new RF sampler}
The DDIM sampler considered in \citet{liang2025low} takes the form

\begin{equation}
    \label{eq: DDIM sampler liang2025 2}
    \begin{aligned}
    & \hat Y^\prime_{\tau -1} = \frac{1}{\sqrt{\alpha_\tau}} \bc{ \hat Y^\prime_\tau + \delta_\tau \widehat s_{Y^\prime_\tau}(\hat Y^\prime_\tau)}, \; \hat Y^\prime_N \sim N(0, I_d)\\
    & \delta_\tau = \frac{1 - \alpha_\tau}{1 + \sqrt{\frac{\alpha_\tau - \omega_\tau}{1 - \omega_\tau}}},\\
    & \tau   = N , N-1, \ldots, 1.
    \end{aligned}
\end{equation}
Using the same transformations detailed in Section \ref{app: Stoc-RF convergence} and plugging them in \eqref{eq: DDIM sampler liang2025 2}, we arrive at the RF sampler

\begin{equation}
     \label{eq: DDIM sampler non-Eulerian}
     \begin{aligned}
      &\frac{Y_{t_{i+1}}}{\sigma_{t_{i+1}}} = \frac{R_{i+1}}{R_i} \bc{ \frac{Y_{t_i}}{\sigma_{t_i}} + \eta_i \sigma_{t_i} \widehat s_{t_i}(Y_{t_i}) },\\
      &  \frac{Y_{t_0}}{\sigma_{t_0}} \sim N(0,  I_d),\\
      & i = 0, 1, \ldots, N-1,
     \end{aligned}
 \end{equation}
where $\eta_i = \frac{1 - \frac{R_i^2}{R_{i+1}^2}}{1 + \sqrt{\frac{R_i^2}{R_{i+1}^2}. \frac{1 - R_{i+1}^2}{1 -  R_{i}^2}}}, R_i = t_i/\sigma_{t_i}$. 


\begin{figure}[h!]
    \centering
    \begin{subfigure}{0.45\textwidth}
        \includegraphics[width=\textwidth]{figures/N_200_geom_timegrid_histogram.pdf}
        \caption{Histogram and kernel density estimation (KDE) plot of time-grid \eqref{eq: gemoetric time discretization} showing U-shaped distribution.}
    \end{subfigure}\hfill
    \begin{subfigure}{0.45\textwidth}
       \includegraphics[width=\textwidth]{figures/N_200_sde_timegrid_histogram.pdf} 
       \caption{Histogram and KDE plot of the transformed time-grid \eqref{eq: DDPM time scheduling} showing U-shaped distribution.}
    \end{subfigure}
    \caption{ Histogram of non-uniform time grids.}
    
    \label{fig: hist of sde timegrid}
\end{figure}

\paragraph{Algebraic simplification.}
Using $R_i=t_i/\sigma_{t_i}$, we have
\[
\frac{R_{i+1}}{R_i}
=
\frac{t_{i+1}}{t_i}
\frac{\sigma_{t_i}}{\sigma_{t_{i+1}}},
\qquad
\frac{R_i^2}{R_{i+1}^2}.\frac{1 - R_{i+1}^2}{1 - R_i^2} = \frac{t_i^2 (1-t_{i+1})^2}{t_{i+1}^2(1-t_i)^2},
\qquad
\eta_i
=
\frac{1-\frac{t_i^2\sigma_{t_{i+1}}^2}{t_{i+1}^2\sigma_{t_i}^2}}{1 + \frac{t_i (1-t_{i+1})}{t_{i+1}(1-t_i)}}.
\]
Multiplying \eqref{eq: DDIM sampler non-Eulerian} by $\sigma_{t_{i+1}}$ yields
\[
Y_{t_{i+1}}
=
\frac{t_{i+1}}{t_i}
\Bigl(
Y_{t_i}
+
\eta_i\,
\sigma_{t_i}^2\,
\widehat s_{t_i}(Y_{t_i})
\Bigr).
\]
A direct calculation using $\sigma_t^2=t^2+(1-t)^2$ shows that
\[
\eta_i\sigma_{t_i}^2
=
\frac{
(t_{i+1}-t_i)
\bigl(1-t_i\bigr)
}{
t_{i+1}
}.
\]
Therefore, the deterministic sampler can be written equivalently as
\begin{equation}
\label{eq:ode-sampler-simplified}
Y_{t_{i+1}}
=
\frac{t_{i+1}}{t_i}\,Y_{t_i}
+
\frac{(t_{i+1}-t_i)(1-t_i)}{t_i}
\widehat s_{t_i}(Y_{t_i}) = Y_{t_i} + \Delta_i \underbrace{\bc{\frac{Y_{t_i}}{t_i} + \frac{1 - t_i}{t_i} \widehat s_{t_i}(Y_{t_i})}}_{ = \widehat v_{t_i}(Y_{t_i})},
\end{equation}
where $ \Delta_i = t_{i+1} - t_i$. This expression matches exactly the form of the RF sampler in \eqref{eq: emp-ode-disc}. However, the time steps $
\Delta_i$ are different from those we propose in \eqref{eq: gemoetric time discretization}, which we used to obtain the RF convergence guarantees given in \Cref{thm: RF - convergence of Eulerian scheme}. Although the time-schedule is still U-shaped as shown in \Cref{fig: hist of sde timegrid} (b), however, it is not symmetric and, as a result, might take more discretization steps at either the beginning or the end of the sampling stage. Whether a symmetric or asymmetric schedule is better remains an open problem. Furthermore, we observe that empirically it appears to be relatively sensitive to the choice of the parameters $c_0$ and $c_1$ in \eqref{eq: DDPM time scheduling}.


\section{Proof of Lemmas related to Stochastic Localization (Section \ref{sec: SL})}
\label{app: SL lemmas}


\subsection{Proof of Lemma \ref{lemma: time change DDPM-SL} and Lemma \ref{lemma: SL and RF are equivalent}}
\label{app: SL = RF = DDPM}
We first show that the solution to the SDE \eqref{eq: forward process} takes the  form
\begin{equation}
\label{eq: forwards process DDS solution}
 \frac{Y'_{\tau}}{\sqrt{\omega_{\tau}}} =  X_1 + \widetilde B_{\frac{1 -\omega_{\tau}}{\omega_{\tau}}}^\prime,  
\end{equation}
where $\omega_{\tau} = \exp(-2 \int_0^\tau \beta(u) du)$ and $\bc{\widetilde B^\prime_\tau}_{\tau\ge 0}$ is a standard Brownian motion. First, it is easy to see that the solution of SDE \eqref{eq: forward process} is 
\begin{equation}
\label{eq: forward process solution}
\frac{Y'_{\tau}}{\sqrt{\omega_{\tau}}} = X_1 + \underbrace{\int_0^\tau \sqrt{2 \beta(t)} \exp \br{\int_0^t \beta(u) du} dB^\prime_t}_{:= M_\tau}.
\end{equation}
We note that $M_\tau$ is continuous local Martingale with $M_0 = 0$. Also, the quadratic variation of $M_\tau$ is 
$$
\innerprod{M}_\tau = \int_{0}^\tau 2 \beta(t) \exp \br{2 \int_{0}^t \beta(u) du} dt = \exp \br{2 \int_{0}^\tau \beta(u) du} -1 = \frac{1 - \omega_\tau}{\omega_\tau}.
$$
Then, by Dambis–Dubins–Schwarz theorem \citep[Theorem 1.6]{revuz2013continuous} there exists a Brownian motion $\bc{\widetilde B_u^\prime}_{u \ge 0}$ such that $\bc{M_\tau}_{\tau \ge 0} = \bc{\widetilde B_{\frac{1 -\omega_{\tau}}{\omega_{\tau}}}^\prime}_{\tau \ge 0}$. Therefore, we have \eqref{eq: forwards process DDS solution} by substituting $M_\tau$ in \eqref{eq: forward process solution}.


Now, we define two processes : $\hat U_s:= s Y^\prime_{\tau(s)}/\sqrt{\omega_{\tau(s)}} = s X_1 + s \widetilde B^\prime_{\frac{1 - \omega_{\tau(s)}}{\omega_{\tau(s)}}}$, and $\tilde U_s := s \tilde X_{t(s)}/t(s) = s X_1 + s \tilde B_{\frac{(1-t(s))^2}{t^2(s)}}$. 
Now, under the prescribed time change in Lemma \ref{lemma: time change DDPM-SL} and Lemma \ref{lemma: SL and RF are equivalent}, we have 
\[
s \widetilde B^\prime_{\frac{1 - \omega_{\tau(s)}}{\omega_{\tau(s)}}} = s \widetilde B^\prime_{1/s}, \quad \text{and} \quad  s \tilde B_{\frac{(1-t(s))^2}{t^2(s)}} = s \tilde B_{1/s}.
\]
Note, that both $\bc{s \widetilde B^\prime_{1/s}}_{s \ge0}$ and $\bc{s \tilde B_{1/s}}_{s \ge 0}$ are Brownian motions. Hence, $\bc{\hat U_s}_{s\ge0}$ and $\bc{\tilde U_s}_{s \ge 0 }$ are equivalent to SL process.


\subsection{Proof of Lemma \ref{prop: RF-SDE and backwrd process are equivalent}: Equivalence of SDE \eqref{eq: RF-SDE} and and SDE \eqref{eq: reverse OU process}}
\label{app: equivalence of backward and stoc-RF}

\subsection*{SDE \eqref{eq: RF-SDE} to SDE \eqref{eq: reverse OU process}:}
We begin with the solution $\tilde Z_t$ of STOC-RF SDE~\eqref{eq: RF-SDE}:
\begin{equation}
\label{eq: d tilde Z_t}
d \tilde Z_t = \bc{\frac{\tilde Z_t}{t} + 2 \br{\frac{1-t}{t}} s_t(\tilde Z_t)}dt + \sqrt{2 \br{\frac{1-t}{t}}} d B_t,
\end{equation}
where $s_t(\cdot)$ is the score of $\tilde Z_t$, and $\bc{B_t}_{t \ge 0}$ is Brownian motion.

Recall the transformation $t(\tau) = \frac{\sqrt{\omega_{\tau}}}{\sqrt{\omega_\tau} + \sqrt{1 - \omega_\tau}}$, where $\omega(\tau)  = \exp \br{- 2 \int_{0}^\tau \beta(u) du}$. 

Now, we define $\overleftarrow Y_\tau:= \frac{\sqrt{\omega_{N-\tau}} \tilde Z_{t(N-\tau)}}{t(N - \tau)}$. We first begin by noting that 
\begin{align}\label{eq:anidentity}
\frac{\sqrt{\omega_{N-\tau}}}{t(N- \tau)} = \frac{1}{ \sigma_{t(N - \tau)}}
\end{align}
We also have $\log \omega_{N- \tau} = - 2 \int_{0}^{N - \tau} \beta(u) du$. Therefore, we have 
\begin{equation}
\label{eq: dtau/dt}
    \begin{aligned}
        &\quad  \frac{1}{\omega_{N - \tau}}. \frac{d \omega_{N-\tau}}{dt} = 2 \beta(N- \tau) \frac{d \tau}{dt}\\
        & \Rightarrow \frac{d \tau}{dt} \stackrel{(i)}{=} \frac{1}{2 \beta(N - \tau) \omega(N - \tau)} . \frac{d\omega_{N - \tau}}{dt} = \frac{1}{2 \beta(N - \tau) \omega(N - \tau)} . \frac{d(t^2/\sigma_t^2)}{dt}\\
        & \Rightarrow \frac{d \tau}{d t } = \frac{1}{\beta(N - \tau)}. \frac{1 - t}{t \sigma_t^2}
    \end{aligned}
\end{equation}
Step (i) follows from~\eqref{eq:anidentity}.
Next, we define the backward time $t_\tau:= t(N- \tau)$, and then we can write $\tilde Z_{t_\tau} = \sigma_{t_\tau} \overleftarrow Y_\tau$. This also yields that the score of $\overleftarrow Y_\tau$ satisfies $s^{\leftarrow}_\tau(\overleftarrow Y_\tau) = \sigma_{t_\tau} s_t(\tilde Z_{t_\tau})$. Moreover, $\eqref{eq: DDPM and RF equivalent}$ yields that $\sigma_{t_\tau} s_t(\tilde Z_{t_\tau}) = \nabla \log q_{N - \tau}(\overleftarrow Y_\tau)$.
Next, with \eqref{eq: dtau/dt} and \eqref{eq: d tilde Z_t} in mind, we can write 

\begin{equation*}
\label{eq: d Y_tau part 1}
\begin{aligned}
d \overleftarrow Y_\tau & = \frac{1}{\sigma_{t_\tau}} d \tilde Z_{t_\tau} - \frac{2t-1}{\sigma_{t_\tau}^3} \tilde Z_{t_\tau} dt_\tau\\
& = \bc{\frac{\tilde Z_{t_\tau}}{t_\tau \sigma_{t_\tau}} + 2 \br{\frac{1 - t_{\tau}}{t_{\tau }\sigma_{t_\tau}} } s_{t_\tau}(\tilde Z_{t_\tau}) - \frac{2 t_{\tau} - 1}{\sigma_{t_\tau}^3} \tilde Z_{t_\tau}} dt_\tau + \sqrt{2 \br{\frac{1 - t_{\tau}}{\sigma_{t_\tau}^2 t_\tau}}} d B_{t_\tau}\\
& = \bc{\frac{1 - t_{\tau}}{t_\tau \sigma^3_{t_\tau}} \tilde Z_{t_\tau} + 2 \br{\frac{1 - t_{\tau}}{t_{\tau }\sigma_{t_\tau}} } s_{t_\tau}(\tilde Z_{t_\tau}) } dt_\tau + \sqrt{2 \br{\frac{1 - t_{\tau}}{\sigma_{t_\tau}^2 t_\tau}}} d B_{t_\tau}\\
& = \bc{\frac{1 - t_{\tau}}{t_\tau \sigma^2_{t_\tau}} \overleftarrow Y_\tau + 2 \br{\frac{1 - t_{\tau}}{t_{\tau }\sigma^2_{t_\tau}} } s^\leftarrow_{t_\tau}(\overleftarrow Y_\tau) } \times \beta(N - \tau) \times \frac{t_{\tau} \sigma_{t_\tau}^2}{1 - t_\tau} d\tau + \sqrt{2 \br{\frac{1 - t_{\tau}}{\sigma_{t_\tau}^2 t_\tau}}} \times \sqrt{\beta(N - \tau) \times \frac{t_{\tau} \sigma_{t_\tau}^2}{1 - t_\tau}}  d \overleftarrow B_{\tau}\\
& = \beta(N - \tau) \bc{\overleftarrow Y_\tau + 2 s_\tau^\leftarrow (\overleftarrow Y_\tau)}d\tau + \sqrt{2 \beta(N  - \tau)}d \overleftarrow B_\tau\\
& = \beta(N - \tau) \bc{\overleftarrow Y_\tau + 2 \nabla \log q_{N - \tau}(\overleftarrow Y_\tau)}d\tau + \sqrt{2 \beta(N  - \tau)}d \overleftarrow B_\tau,
\end{aligned}
\end{equation*}
where $\bc{\overleftarrow B_\tau}_{\tau\ge 0}$ is another Brownian motion.
\paragraph{Initial condition check:} Recall that $\overleftarrow Y_0 = \frac{\sqrt{\omega_{N}} \tilde Z_{t(N)}}{t(N)} = (\sqrt{\omega_N} + \sqrt{1 - \omega_N})\tilde Z_{t(N)}$. Therefore, we have 
$$\Law (\overleftarrow Y_0) = \Law\bs{(\sqrt{\omega_N} + \sqrt{1 - \omega_N}) . \bc{t(N) X_1 + (1-t(N)) X_0}} = \Law(\sqrt{\omega_N} X_1 + \sqrt{1 - \omega_N} X_0) = \Law(Y^\prime _N).$$

This shows that $\overleftarrow Y_\tau$ satisfies the backward DDPM SDE \eqref{eq: reverse OU process}.  

\subsection*{SDE \eqref{eq: reverse OU process} to SDE \eqref{eq: RF-SDE}:}
For the other direction, we let $\bc{\overleftarrow Y_\tau}_{\tau\ge 0}$ be a solution of SDE \eqref{eq: reverse OU process}. Now, to generate a solution of SDE \eqref{eq: RF-SDE}, we need to define 
\[
\tilde Z_t = \frac{t(N - \tau(t)) \overleftarrow Y_{\tau(t)}}{\sqrt{\omega_{N - \tau(t)}}} = \frac{\overleftarrow Y_{\tau(t)}}{\sqrt{\omega_{N - \tau(t)}} + \sqrt{1 - \omega_{N - \tau(t)}}},
\]
where $\tau(t)$ should satisfy the following: 
\begin{equation}
\label{eq: dtau/dt}
\begin{aligned}
  t(N - \tau(t)) =  & \quad \frac{\sqrt{\omega_{N - \tau(t)}}}{\sqrt{\omega_{N - \tau(t)}} + \sqrt{1  - \omega_{N - \tau(t)}}} = t\\
    & \Rightarrow \omega_{N - \tau(t)} = \frac{t^2}{t^2 + (1-t)^2}\\
    & \Rightarrow \exp \br{- 2 \int_0^{N - \tau(t)} \beta(u) du} = \frac{t^2}{t^2 + (1-t)^2}\\
    & \Rightarrow \int_{0}^{N - \tau(t)} \beta(u) du = \log \br{\sqrt{1 + \frac{(1-t)^2}{t^2}}}.
\end{aligned}
\end{equation}
This shows that $\tilde Z_t = \sigma_{t} \overleftarrow Y_{\tau(t)}$. Additionally, \eqref{eq: DDPM and RF equivalent} yields that the score of $\tilde Z_t$ is $s_t(\cdot)$.
Now, for convenience, we now write $\tau$ to denote $\tau(t)$.
Taking derivative w.r.t. $t$ on both sides of \eqref{eq: dtau/dt} yields
\[
\frac{d \tau}{dt } =  \frac{1}{\beta(N - \tau)}. \frac{(1-t)}{t \sigma_t^2}.
\]

Using these facts, we finally get 
\begin{equation*}
    \begin{aligned}
        d \tilde Z_t & = \frac{(2t - 1)}{\sigma_t} \overleftarrow Y_\tau + \sigma_t d\overleftarrow Y_{\tau} \\
        & = \frac{(2t -1)}{\sigma_t^2} \tilde Z_t + \sigma_t \bc{\overleftarrow Y_\tau + 2 s_\tau^\leftarrow (\overleftarrow Y_\tau)} \beta(N - \tau) d\tau + \sigma_t \sqrt{2 \beta(N - \tau)} d \overleftarrow B_\tau\\
        & = \frac{(2t -1)}{\sigma^2_t} \tilde Z_t + \sigma_t \bc{\frac{\tilde Z_t}{\sigma_t} + 2 \sigma_t s_t(\tilde Z_t)} \frac{(1-t)}{t \sigma_t^2} dt + \sqrt{2 \br{\frac{1-t}{t}}} d \tilde B_t\\
        & = \br{ \frac{2t - 1}{\sigma_t^2} + \frac{1-t}{t \sigma_t^2}} \tilde Z_t + 2 \br{\frac{1-t}{t}} s_t(\tilde Z_t)+ \sqrt{2 \br{\frac{1-t}{t}}} d \tilde B_t\\
        & = \frac{\tilde Z_t}{t}+  2 \br{\frac{1-t}{t}} s_t(\tilde Z_t)+ \sqrt{2 \br{\frac{1-t}{t}}} d \tilde B_t.
    \end{aligned}
\end{equation*}
%The initial conditions can also be checked as before. 
Therefore, we $\tilde Z_t$ satisfies SDE \eqref{eq: RF-SDE}.

\subsection{Proof of Proposition \ref{prop: SL and general interpolant are equivalent}}
\label{app: proof of SL and general interpolants equivalence}
Recall that $\frac{\tilde I_{\theta(s)}}{a_{\theta(s)}} = X_1 + W_{r^2_{\theta(s)}}$, where $r^2_{\theta(s)} = 1/s$. Now, we define $\hat U_s:= \frac{\tilde I_{\theta(s)}}{a_{\theta(s)}}$. Then, we have $s \hat U_s = s X_1 + s W_{1/s}$. Similar to the proofs of Lemma \ref{lemma: time change DDPM-SL} and Lemma \ref{lemma: SL and RF are equivalent}, the process $\bc{\hat W_s}_{s \ge 0}:= \bc{s W_{1/s}}_{s \ge 0}$ is also a Brownian motion. This finishes the proof by the description of the SL process in \eqref{eq: SL process}. 

\section{Auxiliary results}
\label{app: auxiliary results}

\subsection{Controlling posterior covariance of RF}
\label{sec: RF covariance}

% \textcolor{violet}{ALE: I wonder if we really need this in the main text. This is an important but technical result but I think it belong in the appendix}
We recall the SL process \eqref{eq: SL process} to study the dynamics of the conditional covariance matrix $ \bSigma_t := \Cov_{1 \mid t}(X_t)$ where $X_t = t X_1 + (1-t)X_0$ and $X_0 \sim N(0, I_d)$. Note that $\bSigma_t$ is a random matrix. Recall the defintion of $\tilde X_t$ in \eqref{eq: RF linear stoc process}. It is evident that $\tilde \bSigma_t:= \Cov_{1 \mid t}(\tilde X_t) \overset{d}{=}\bSigma_t  $. Due to Lemma  \ref{lemma: SL and RF are equivalent}, the SL process \eqref{eq: SL process} is equivalent to the process \eqref{eq: RF linear stoc process} under time change $t(s):= \frac{\sqrt{s}}{1 + \sqrt{s}}$. Then, one can substitute this in Lemma \ref{lemma: eldan ddt cov lemma} to get the following result:
\begin{lemma}[\Cref{cor: my ddt cov version}]
    \label{lemma: RF ddt cov lemma}
    for $t \in [0,1)$, we have $\frac{d}{dt} \bbE (\bSigma_t) = - \frac{2t}{(1-t)^3} \bbE(\bSigma_t^2)$.
\end{lemma}
\begin{proof}
    Now, we note that $\bA_s$ and $\tilde \bSigma_{t(s)}$ (so is $\bSigma_{t(s)}$) are marginally equal in law. Therefore, by chain rule we have the following:

\begin{align*}
    \frac{d}{dt} \bbE(\bSigma_t) & = \frac{d}{dt} \bbE(\tilde \bSigma_t)\\
    &= \frac{d}{dt} \bbE (\bA_s)\\
    & = \frac{2t}{(1-t)^3} \frac{d}{ds} \bbE (\bA_s)\\
    & = - \frac{2t}{(1-t)^3} \bbE(\bA_s^2)\\
    & = - \frac{2t}{(1-t)^3} \bbE(\bSigma_t^2).
\end{align*}
\end{proof}
As mentioned in previous section, the above result allows to control the discretization error of the ODE and SDE sampler of RF. More concretely, Lemma \ref{lemma: RF ddt cov lemma} allows us to bound the spectra of $\bSigma_t^2$ in terms of the spectra of $\bSigma_t$, thereby producing tighter bounds for convergence rate.



\subsection{Proof of Proposition \ref{prop: Main cov bound via trace}}
\label{app: proof of Cov Frobenius bound}
% We recall Lemma \ref{lemma: RF ddt cov lemma}. i.e, 
% \[
% \frac{d}{dt} \bbE (\bSigma_t) = - \frac{2t}{(1-t)^3} \bbE(\bSigma_t^2).
% \]

% The above display leads us to the following lemma:
% \begin{lemma}
%     \label{lemma: dynamics of conditional covariance}
%     For all $t \in [0,1)$, we have $\frac{d}{dt} \bbE (\bSigma_t) = - \frac{2t}{(1-t)^3} \bbE(\bSigma_t^2)$.
% \end{lemma}

We recall the definitions of $(\bA_s)_{s\ge 0}$ and $(\tilde \bSigma_{t(s)})_{s\ge 0}$ in Section \ref{sec: SL} and Section \ref{sec: RF covariance}. Note that $(\bA_s)_{s\ge 0}$ and $(\tilde \bSigma_{t(s)})_{s\ge 0}$ (see Section \ref{sec: SL} and \ref{sec: RF covariance}) are also equivalent in law. This is simply due to the fact that $\tilde \bSigma_{t(s)} = \Cov_{1 \mid t(s)}(\tilde X_{t(s)}) = \Cov_{1 \mid t(s)}(s \tilde X_{t(s)}/t(s))$, and the distributional equivalence of the two processes $(U_s)_{s\ge 0}$ and $(s \tilde X_{t(s)}/t(s))_{s\ge0}$. It is known that the random matrix valued process $(\bA_s)_{s\ge 0}$ obeys the following SDE:
\begin{equation}
\label{eq: SDE A_s}
d  \bA_s = -  \bA_s^2 ds + \cM_s^{(3)} dB_s,
\end{equation}
where $\cM_s^{(\ell)} = \bbE[(X_1 - \bbE[X_1 \mid U_s])^{\otimes \ell} \mid U_s]$. More details on this can be found in \citep[Section 4.2.1]{eldan2020taming} or \citep[Section 3]{alberts2025joint}. Since the two processes $(\bA_s)_{s\ge 0}$ and $(\tilde \bSigma_{t(s)})_{s\ge 0}$ are equivalent in law, SDE \eqref{eq: SDE A_s} immediately yields 
%\red{(SR: SDE (20) is not used anywhere)}
\begin{equation}
    \label{eq: SDE Sigma_t}
    d \tilde \bSigma_t = - \frac{2t}{(1-t)^3} \tilde \bSigma_t^2 dt + \tilde \cM_t^{(3)} dB_{t^2/(1-t)^2},
\end{equation}
where $\tilde \cM_t^{(\ell)} = \bbE[(X_1 - \bbE[X_1 \mid \tilde X_t])^{\otimes \ell} \mid \tilde X_t]$.
Now, an application of Ito's formula on \eqref{eq: SDE A_s} yields 
\[
d(\tr( \bA_s^2)) = -  2\innerprod{ \bA_s,  \bA_s^2}ds  + \innerprod{ \cM_s^{(3)}, \cM_s^{(3)}} ds + 2 \innerprod{\bA_s, \cM_s^{(3)}} dB_s
\]
Taking expectation on both side we arrive at the following inequality:
\[
d(\tr[\bbE(\bA_s^2)]) = - 2 \bbE\left[\innerprod{\bA_s, \bA_s^2}\right]ds + \bbE\left[\innerprod{ \cM_s^{(3)}, \cM_s^{(3)}}\right]ds.
\]
Finally, using the fact that $\bA_s$ and $\bSigma_{t(s)}$ have the same law, chain rule yields the following:
\begin{equation}
    \label{eq: Moment Sigma_t ODE}
    d(\tr[\bbE(\bSigma_{t}^2)]) = - \frac{4t}{(1-t)^3} \bbE\left[\innerprod{\bSigma_t, \bSigma_t^2}\right]dt + \frac{2t}{(1-t)^3} \bbE\left[\innerprod{ \tilde \cM_t^{(3)}, \tilde\cM_t^{(3)}}\right]dt.
\end{equation}
We first focus on the first term on the right-hand side of \eqref{eq: Moment Sigma_t ODE}. Using the symmetry of $\bSigma_t$, we observe that

\begin{equation}
    \label{eq: bound on the leading term of sde}
    \begin{aligned}
    &\bbE \left[\innerprod{\bSigma_t, \bSigma_t^2}\right] = \bbE \left[\tr(\bSigma_t^3)\right]\\
    & \le \bbE \left[\Norm{\bSigma_t}_{\op} \Norm{\bSigma_t}^2_F\right] \\
    & \le \bbE \left[\tr\left(\bSigma_t\right) \Norm{\bSigma_t}^2_F\right]\\
    & =  \bbE \left[\tr\left(\bSigma_t\right) \ind\{X_t \in \cT_t\} \Norm{\bSigma_t}^2_F\right] + \bbE \left[\tr\left(\bSigma_t\right) \ind\{X_t \notin \cT_t\}\Norm{\bSigma_t}^2_F\right]\\
    & \le C_3 \frac{(1 - t)^2}{t^2}(k \log N) \bbE \left[\Norm{\bSigma_t}_F^2\right] + \bbE \left[\tr\left(\bSigma_t\right) \ind\{X_t \notin \cT_t\}\Norm{\bSigma_t}^2_F\right],
    \end{aligned}
\end{equation}
where the last inequality follows due to \eqref{eq: trace_Cov bound}. Regarding the last term in above display, Assumption \ref{assumption: bounded support} and \eqref{eq: prob of exclusion} yields
\[
\bbE \left[\tr\left(\bSigma_t\right) \ind\{X_t \notin \cT_t\}\Norm{\bSigma_t}^2_F\right] \le 8 N^{3 c_R} \pr(X_t \notin \cT_t) \le \frac{1}{N^{10}}.
\]
Therefore, we arrive at 
\begin{equation}
    \label{eq: bound on Sigma_t^3}
    \bbE \left[\innerprod{\bSigma_t, \bSigma_t^2}\right] \le C_3 \frac{(1 - t)^2}{t^2}(k \log N) \bbE \left[\Norm{\bSigma_t}_F^2\right]  + \frac{1}{N^{10}}.
\end{equation}
Substituting this in \eqref{eq: bound on the leading term of sde} yields,
\begin{align*}
    d\bbE(\Norm{\bSigma_t}_F^2) & = -\frac{4t}{(1-t)^3} \bbE\left[\innerprod{\bSigma_t, \bSigma_t^2}\right]dt + \frac{2t}{(1-t)^3} \bbE\left[\innerprod{ \tilde \cM_t^{(3)}, \tilde \cM_t^{(3)}}\right]dt\\
    & \ge - \frac{4t}{(1-t)^3} \bbE\left[\innerprod{\bSigma_t, \bSigma_t^2}\right]dt\\
    &\ge - \frac{4 C_3 k \log N}{(1-t)t}\bbE \left[\Norm{\bSigma_t}_F^2\right] dt - \frac{4t}{(1-t)^3}. \frac{1}{N^{10}} dt\\
    & \ge - \frac{4 C_3 k \log N}{(1-t_i)t_{i-1}}\bbE \left[\Norm{\bSigma_t}_F^2\right] dt - \frac{4t_i}{(1-t_{i})^3}. \frac{1}{N^{10}} dt,
\end{align*}
where the last line follows holds provided $t \in [t_{i-1}, t_i]$ with $i> 1$. Therefore, we can conclude that
\begin{align*}
    & \exp\left\{\frac{4 C_3 k t\log N}{(1-t_i)t_{i-1}}\right\} \bbE (\Norm{\bSigma_{t}}_F^2) -  \exp\left\{\frac{4 C_3 k t_{i-1}\log N}{(1-t_i)t_{i-1}}\right\} \bbE (\Norm{\bSigma_{t_{i-1}}}_F^2)\\
    & \ge - \frac{4t_i}{(1-t_{i})^3}. \frac{1}{N^{10}} \int_{t_{i-1}}^t \exp\left\{\frac{4 C_3 k t\log N}{(1-t_i)t_{i-1}}\right\}dt\\
    & \ge - \frac{t_i t_{i-1}}{C_3 k (1-t_{i})^2 \log N}. \frac{1}{N^{10}} \left(\exp\left\{\frac{4 C_3 k t\log N}{(1-t_i)t_{i-1}}\right\} - \exp\left\{\frac{4 C_3 k t_{i-1}\log N}{(1-t_i)t_{i-1}}\right\}\right)
\end{align*}
Dividing both sides by $\exp\left\{\frac{4 C_3 k t_{i-1}\log N}{(1-t_i)t_{i-1}}\right\}$, we obtain
\begin{align*}
    & \exp\left\{\frac{4 C_3 k (t-t_{i-1})\log N}{(1-t_i)t_{i-1}}\right\} \bbE (\Norm{\bSigma_{t}}_F^2) -  \bbE (\Norm{\bSigma_{t_{i-1}}}_F^2)\\
    & \ge - \frac{t_i t_{i-1}}{C_3 k (1-t_{i})^2 \log N}. \frac{1}{N^{10}} \left(\exp\left\{\frac{4 C_3 k (t-t_{i-1})\log N}{(1-t_i)t_{i-1}}\right\} -1\right)\\
    & \ge -\frac{1}{C_3 k (1-t_i)^2 \log N}. \frac{1}{N^{10}} \left(\exp\left\{\frac{4 C_3 k (t-t_{i-1})\log N}{(1-t_i)t_{i-1}}\right\} -1\right).
\end{align*}
Now, recall that $\frac{\eta_{t_i}}{(1-t_i)t_i} \le 2 h$. Using this inequality coupled with \eqref{eq: value of h} gives us
\[
\frac{4 C_3 k (t-t_{i-1})\log N}{(1-t_i)t_{i-1}} \le \frac{4 C_3 k \eta_{i-1}\log N}{(1-t_i)t_{i-1}} \le 8 C_3 k h \log N< 1,
\]
provided $N$ is sufficiently large. Therefore, we have 
\begin{align*}
     \bbE (\Norm{\bSigma_{t_{i-1}}}_F^2) & \le \exp\left\{\frac{4 C_3 k (t-t_{i-1})\log N}{(1-t_i)t_{i-1}}\right\} \bbE (\Norm{\bSigma_{t}}_F^2)\\
     & \quad + \frac{1}{C_3 k (1-t_i)^2 \log N}. \frac{1}{N^{10}} \left(\exp\left\{\frac{4 C_3 k (t-t_{i-1})\log N}{(1-t_i)t_{i-1}}\right\} -1\right)\\
    & \le 3 \bbE (\Norm{\bSigma_{t}}_F^2) + \frac{24h }{(1-t_i)^2 N^{10}}.
\end{align*}
The inequality follows from $e^x -1 \le x$ for $ x\in (0,1)$, and $\frac{\eta_{i}}{(1-t_i)t_i} \le 2 h$ (\Cref{lemma: properties of time-steps}).
Substituting $t_{i-1}$ by $t_{i}$ in the above inequality, we get the following 
\[
\bbE (\Norm{\bSigma_{t_{i}}}_F^2) \le 3 \bbE (\Norm{\bSigma_{t}}_F^2) + \frac{24h }{(1-t_{i+1})^2 N^{10}}, \quad t \in [t_i , t_{i+1}], i \ge 1.
\]

Now we use Lemma \ref{lemma: RF ddt cov lemma} to get 
\begin{equation*}
%\label{eq: Main cov bound via trace}
    \begin{aligned}
        \bbE [\tr(\bSigma_{t_{i+1}})] - \bbE[\tr(\bSigma_{t_{i}})] & =  -\int_{t_{i}}^{t_{i+1}} \frac{2t}{(1-t)^3} \bbE(\Norm{\bSigma_t}_F^2)\\
        & \le - \int_{t_i}^{t_{i+1}} \frac{2t }{3(1-t)^3 }\bbE(\Norm{\bSigma_{t_i}}_F^2)dt + \frac{8 h \eta_{i}}{(1-t_{i+1})^2 N^{10}}\\
        & \le - \frac{2 \eta_{t_i}(t_{i+1}+ t_i)}{3(1-t_{i+1})^2 (1-t_i)^2}  \bbE(\Norm{\bSigma_{t_i}}_F^2) + \frac{8 h \eta_{i} }{(1-t_{i+1})^2 N^{10}}
    \end{aligned}
\end{equation*}

Rearranging the terms in the above inequality and using $t_{i+1} + t_i> 2t_i$, we finally get 
\begin{equation}
    \label{eq: Main cov bound via trace}
    \frac{ \eta_{t_i}t_i}{(1-t_{i+1})^2 (1-t_i)^2}  \bbE(\Norm{\bSigma_{t_i}}_F^2)  \le \frac{3}{4} \left(\bbE [\tr(\bSigma_{t_{i}})] - \bbE[\tr(\bSigma_{t_{i+1}})]\right) + \frac{6 h \eta_{i} }{(1-t_{i+1})^2 N^{10}}.
\end{equation}

\subsection{Properties of  time scheduling \eqref{eq: gemoetric time discretization}}


\begin{lemma}
    \label{lemma: properties of time-steps}
    The time-steps \eqref{eq: gemoetric time discretization} satisfies the following for $i \in \{1, \ldots, N-2\}$:
    \begin{enumerate}[label = (\alph*)]
    \item\label{item: part a} We have $\eta_i = h t_i \ind\bc{t_{i} \le 1/2} + h (1-t_{i+1}) \ind\bc{t_i > 1/2}$.
    
    \item\label{item: part b} We have $\frac{\eta_i}{1-t_i} \le h$ and $\frac{\eta_i}{1-t_{i+1}} \le h$.
    
     \item\label{item: part c} We have $\frac{ \eta_i^2 (1-t_i)^2 }{(1-t_{i+1})^2} \le h^2 $, as $\frac{\eta_i}{(1-t_{i+1})} \le h$ and $(1-t_i) \le 1$.
     
     \item\label{item: part d} We have $\sum_{i : t_i >0}\frac{ \eta_i^2 }{(1-t_{i+1})^2 t_i^2} = \sum_{i: 0 < t_{i} < 1/2 }\frac{ \eta_i^2 }{(1-t_{i+1})^2 t_i^2} + \sum_{i: t_{i}\ge 1/2 }\frac{ \eta_i^2 }{(1-t_{i+1})^2 t_i^2} \le  4 h^2 N$.
\end{enumerate}
\end{lemma}

\begin{proof}
{\color{white} Ghost line}
\begin{itemize}
    \item \textbf{Part \ref{item: part a}:} If $t_i < 1/2$, we have $\eta_i = t_{i+1} - t_i = (1+h)t_i - t_i = h t_i$. For the case $t_{i}\ge 1/2$, we have 
    \begin{align*}
        \eta_i &= t_{i+1} - t_i\\
        & = (1- t_i) - (1 - t_{i+1})\\
        & = h (1- t_{i+1}).
    \end{align*}

    \item \textbf{Part \ref{item: part b}:} For $t_i < 1/2$, we have $\frac{\eta_i}{1 - t_i} = ht_i/(1-t_i) \le h$. We use the fact that $\frac{t}{1-t}<1$ for $t \le 1/2$. Similarly, $\frac{\eta_i}{1 - t_{i+1}} = \frac{h t_i}{1 - t_{i+1}}$. Now, note that $t_i <1/2$ implies that $t_{i+1} \le 1/2$. Therefore, $\frac{t_i}{1 - t_{i+1}}\le 1$, and it yields $\frac{\eta_i}{1 - t_{i+1}} = \frac{h t_i}{1 - t_{i+1}} < h$.
    
    For $t_i\ge 1/2$, we have $\frac{\eta_i}{1-t_i} = \frac{h(1- t_{i+1})}{1 - t_i} < h$. Similarly, $\frac{\eta_i}{1-t_{i+1}} = \frac{h(1- t_{i+1})}{1 - t_{i+1}} = h$. This yields the coveted bound.

    \item \textbf{Part \ref{item: part c}:} The proof immediately follows from Part \ref{item: part b}.


    \item \textbf{Part \ref{item: part d}:} Fist, from Part \ref{item: part a}, we can note that $\frac{\eta_{i}}{t_i} \le h$. Also, Part \ref{item: part b} yields that $\frac{\eta_i}{1 - t_{i+1}} \le h$. Therefore, 
    \begin{align*}
        & \sum_{i : t_i >0}\frac{ \eta_i^2 }{(1-t_{i+1})^2 t_i^2}\\
        & = \sum_{i: 0 < t_{i} < 1/2 } \blue{\frac{ \eta_i^2 }{ t_i^2}}. \frac{1}{(1-t_{i+1})^2} + \sum_{i: t_{i}\ge 1/2 } \blue{\frac{ \eta_i^2 }{(1-t_{i+1})^2 }}. \frac{1}{t_i^2} \\
        & \le \sum_{i: 0 < t_{i} < 1/2 } 4 h^2 + \sum_{i: t_{i}\ge 1/2 } 4 h^2\\
        & \le 4h^2 N.
    \end{align*}
\end{itemize}
    

    


    


    % \paragraph{Part 3:}sds
\end{proof}

\begin{lemma}
\label{lemma: simple identity 1}
Let $t_{i+1} = t_i + \eta_i$ with $0 < t_i \le t_{i+1} \le 1$. Then
\[
\frac{(1-t_i)^2}{t_i^2}
-
\frac{(1-t_{i+1})^2}{t_{i+1}^2}
\le
\frac{2\eta_i}{t_i^3}.
\]
\end{lemma}

\begin{proof}
Define $f(t) := \frac{(1-t)^2}{t^2}$ for $t > 0$. A direct calculation gives
\[
f'(t) = \frac{2(t-1)}{t^3}.
\]
By the mean value theorem, there exists $\xi \in (t_i, t_{i+1})$ such that
\[
f(t_i) - f(t_{i+1})
= -f'(\xi)\,(t_{i+1}-t_i)
= \frac{2(1-\xi)}{\xi^3}\,\eta_i.
\]
Since $\xi \ge t_i$ and $\xi \le 1$, we have
\(
(1-\xi)/\xi^3 \le 1/t_i^3
\),
which implies the desired bound.
\end{proof}


%%%%%%%%%%%%%%%%%%%%%%%%%%%%%%%%%%%%%%%%%%%%%%%%%%%%%%%%%%%%%%%%%%%%%%%%%%%%%%%
%%%%%%%%%%%%%%%%%%%%%%%%%%%%%%%%%%%%%%%%%%%%%%%%%%%%%%%%%%%%%%%%%%%%%%%%%%%%%%%

\end{document}
